\documentclass[11pt]{article}

\usepackage[a4paper,margin=1in]{geometry}
\usepackage{amsmath,amssymb,amsthm,mathtools}
\usepackage{booktabs,array,longtable}
\usepackage{enumitem}
\usepackage{microtype}
\usepackage{xcolor}
\usepackage{hyperref}
\usepackage[nameinlink,capitalise,noabbrev]{cleveref}

\hypersetup{
  colorlinks=true,
  linkcolor=blue!55!black,
  citecolor=blue!55!black,
  urlcolor=blue!55!black
}

\newtheorem{theorem}{Theorem}[section]
\newtheorem{proposition}[theorem]{Proposition}
\newtheorem{lemma}[theorem]{Lemma}
\newtheorem{corollary}[theorem]{Corollary}
\newtheorem{definition}[theorem]{Definition}
\newtheorem{remark}[theorem]{Remark}

% The environments above share the theorem counter, so cleveref would name
% them all "Theorem"; these aliases restore the correct names.
\usepackage{etoolbox}
\AtBeginEnvironment{proposition}{\crefalias{theorem}{proposition}}
\AtBeginEnvironment{lemma}{\crefalias{theorem}{lemma}}
\AtBeginEnvironment{corollary}{\crefalias{theorem}{corollary}}
\AtBeginEnvironment{definition}{\crefalias{theorem}{definition}}
\AtBeginEnvironment{remark}{\crefalias{theorem}{remark}}

\numberwithin{equation}{section}

\newcommand{\one}{\mathbf 1}
\newcommand{\R}{\mathbb R}
\newcommand{\ip}[2]{\left\langle #1,#2\right\rangle}
\newcommand{\norm}[1]{\left\lVert #1\right\rVert}
\newcommand{\Tr}{\operatorname{Tr}}
\newcommand{\ord}{\operatorname{ord}}
\newcommand{\pos}[1]{\left(#1\right)_{+}}
\newcommand{\dd}{\,\mathrm d}
\newcommand{\E}{\mathbb E}
\newcommand{\GammaLaw}{\operatorname{Gamma}}
\newcommand{\BetaLaw}{\operatorname{Beta}}
\newcommand{\supp}{\operatorname{supp}}

\newcommand*{\tk}[1]{\textcolor{blue}{\textbf{[TK: #1]}}}

% Passages written or reordered in the 2026-08-07 pass, awaiting the author's
% reading.  Delete the environment (or make it a no-op) once they are cleared.
\definecolor{revised}{rgb}{0.56,0.30,0.30}
\newenvironment{rev}{\color{revised}}{}
\newcommand*{\rv}[1]{{\color{revised}#1}}
% Text that was only moved, not rewritten: delimited rather than coloured.
\newcommand*{\revbegin}{\textcolor{revised}{\textbf{[moved, text unchanged $\downarrow$]}}}
\newcommand*{\revend}{\textcolor{revised}{\textbf{[$\uparrow$ end moved]}}}

\title{A Complete Proof of the Goodman--Style Lower Bound\\
for Odd-Cycle Densities in Graphons}
\author{Consolidated working note for H.~Yang and collaborators}
\date{July 2026}


\begin{document}
\maketitle

\begin{abstract}
  Let $W:[0,1]^2\to[0,1]$ be a graphon of edge density $p$, and let $m\ge3$ be odd.
  We prove the conjectured Goodman--style inequality
  \[
    t(C_m,W)\ge p^m-p(1-p)^{m-1}
  \]
  for every $p\in[0,1]$.
  The proof is divided at $p=2/3$, or equivalently at the complement density $q=1/3$.

  For $p\ge2/3$, a two-sided Schur-complement identity cancels the unknown odd moments of the compressed complement operator.
  Expansion in complete homogeneous symmetric polynomials is then diagonalized by a Dirichlet average, converted to a beta integral, and smoothed by a Laplace transform into expectations of a shifted gamma random variable.
  The remaining positivity follows from a uniform odd-to-even gamma moment inequality.

  For $1/2<p<2/3$, the Hilbert--Schmidt bound allows at most one compression eigenvalue above $q=1-p$.
  Isolating this frontier eigenvalue and retaining both its direct coupling and the coupling forced into the safe spectral subspace leads to a one-dimensional Huber-type envelope.
  Its dual form is bounded below by a single function of one dimensionless ratio, and the resulting scalar inequality is settled by a division of the parameter region into two halves.
  The cycles $C_3,C_5,C_7$ are also given separate elementary proofs.

  A final section extracts the reusable mechanisms: one-step Jacobi coefficient stripping, first-return and cyclic excursion expansions, Dirichlet probabilistic polarization, beta--gamma/Laplace positivity transfer, the Laguerre differential equation hidden in the gamma moment lemma, and the one-dangerous-mode/Huber principle.
  It also explains why the same threshold $q=1/3$ emerges independently from the two sides.
\end{abstract}

\tableofcontents

\section{The theorem, notation, and proof map}\label{sec:intro}

For a finite graph $F$ and a graphon $W$, write
\[
  t(F,W)=\int_{[0,1]^{|V(F)|}}
  \prod_{ij\in E(F)}W(x_i,x_j)
  \prod_{i\in V(F)}\dd x_i.
\]
In particular,
\[
  p=t(K_2,W)=\int_{[0,1]^2}W(x,y)\dd x\dd y
\]
is the edge density.
We prove the following theorem.

\begin{theorem}[Odd-cycle Goodman bound]\label{thm:main}
  Let $W$ be a graphon of edge density $p\in[0,1]$.
  For every odd integer $m\ge3$,
  \begin{equation}\label{eq:main-goal}
    t(C_m,W)\ge p^m-p(1-p)^{m-1}.
  \end{equation}
\end{theorem}

Put
\[
  U=1-W,\qquad q=t(K_2,U)=1-p.
\]
If $p\le1/2$, then
\[
  p^m-pq^{m-1}=p\bigl(p^{m-1}-q^{m-1}\bigr)\le0,
\]
whereas $t(C_m,W)\ge0$.
Thus, only $p>1/2$, or $q<1/2$, needs proof.
The nontrivial range splits into
\[
  \underbrace{0\le q\le\frac{1}{3}}_{1 \ge p\ge2/3}
  \qquad\text{and}\qquad
  \underbrace{\frac{1}{3}<q<\frac{1}{2}}_{1/2<p<2/3}.
\]

\begin{corollary} \label{cor:main}
  For every graphon $W$ of edge density $p = 1 - 1 / k$ for some integer $k \ge 2$, and for every odd integer $m \ge 3$, the balanced complete $k$-partite graphon
  \[
    W_k(x, y) = \begin{cases}
      1 & \text{if } \lfloor kx \rfloor \neq \lfloor ky \rfloor,
      \\
      0 & \text{otherwise}.
    \end{cases}
  \]
  is the minimiser of $t(C_m, W)$ among all graphons of edge density $p$.
\end{corollary}

\subsection{Proof map}

Every part of the proof works with the same few quantities, so we introduce them first.
Let $T_U$ be the integral operator with kernel $U=1-W$, and split it against the constant function $\one$.
The constant direction contributes the scalar $q$, and the remainder is the mean-zero function
\[
  g=T_U\one-q\one
  =p-d_W,\qquad d_W(x)=\int_0^1W(x,y)\dd y,
\]
that is, the mean $p$ minus the degree function of $W$.
We also restrict $T_U$ to the mean-zero subspace $\one^\perp$, and write $A$ for the resulting self-adjoint operator, so that $\ip f{Ah}=\ip f{T_Uh}$ for all $f,h\perp\one$.

Throughout the paper, $\lambda_1,\lambda_2,\ldots$ denote the eigenvalues of $A$, and $P_i$ the orthogonal projection onto the eigenspace of $\lambda_i$.
We write
\[
  \mu=\sum_i\norm{P_ig}^2\delta_{\lambda_i}
\]
for the associated measure on the spectrum of $A$, so that $\int\lambda^j\dd\mu(\lambda)=\ip g{A^jg}$ for every $j\ge0$ and $\mu(\R)=\norm g^2$.

\begin{enumerate}[label=\textbf{Section \arabic*.},leftmargin=4.6em,start=2]
  \item \emph{The cases $m=3,5,7$.}
  We prove these cases separately so that when we consider the general odd $m$ in the later sections, we can assume $m\ge9$.
  The case $m=3$ is Goodman's original bound, in graphon form.
  For $m=5$ and $m=7$, the inclusion--exclusion principle applied to $\prod_{i=1}^m(1-U_i)$ expands $t(C_m,W)$ over the subsets of the $m$ edges.
  A proper subset selects a disjoint union of paths, and therefore contributes a product of the path densities
  \[
    x_j=t(P_j,U)=\ip\one{T_U^j\one},
  \]
  while the full subset contributes $t(C_m,U)$, for which $t(C_m,U)\le x_{m-1}$ because $0\le U\le1$.
  Each $x_j$ can in turn be written as a polynomial in $q$ and the moments
  \[
    s_j=\ip g{A^jg}=\int\lambda^j\dd\mu(\lambda).
  \]
  For $m=5$, the quantity to be shown nonnegative is a quadratic form in $g$ and $T_Ug$, which \Cref{subsec:C5} writes as a sum of two squares.
  For $m=7$, it becomes
  \[
    6s_1^2+\int R_{q,r}(\lambda)\dd\mu(\lambda),
    \qquad r=\norm g^2,
  \]
  where $R_{q,r}$ is a quartic in $\lambda$; \Cref{subsec:C7} writes it as an sum of squares, nonnegative for every real $\lambda$.

  \item In this section, we prove lemmas that are used throughout the remaining part of this paper.
  We first prove that, in proving \eqref{eq:main-goal}, it is enough to consider step graphons, so that every computation below is finite dimensional.
  We then show that the logarithm of a determinant generates the cycle densities: for any kernel $K$,
  \[
    -\log\det(I-zT_K)=\sum_{j\ge1}\frac{\Tr(T_K^j)}{j}z^j,
    \qquad
    \Tr(T_K^j)=t(C_j,K)\quad(j\ge2).
  \]
  We apply this to $U$ and to $W$, and factorize the two determinants by Schur complementation against the constant direction.
  This produces two further generating functions $L_-(z)$ and $L_+(z)$, which satisfy the following two equalities (\Cref{lem:cycle-identities}):
  \begin{align*}
    t(C_m,U)&=\phantom{-}\Tr(A^m)+q^m+m[z^m]L_+(z),\\
    t(C_m,W)&=-\Tr(A^m)+p^m+m[z^m]L_-(z).
  \end{align*}
  \item In this section, we prove \Cref{thm:main} for the case $p \ge 2/3$.
  We first eliminate the $\Tr(A^m)$ by addint the two identities:
  \[
    t(C_m,W)+t(C_m,U)=p^m+q^m+S_m,
    \qquad
    S_m=m[z^m]\bigl(L_-(z)+L_+(z)\bigr).
  \]
  Bounding $t(C_m,U)$ by the path density $x_{m-1}$, everything reduces to the nonnegativity of
  \[
    \Phi_m=q^{m-1}+S_m-x_{m-1}.
  \]

  Expanding $L_\pm$ and $x_{m-1}$ gives $\Phi_m$ as an integral against $\dd\mu^{\otimes r}$ of a symmetric polynomial (\Cref{prop:pge23-expansion}): with $n=m-2r$,
  \[
    \Phi_m=\sum_{r=1}^{(m-1)/2}\int\!\cdots\!\int
    P_{m,r}(q;\lambda_1,\ldots,\lambda_r)
    \prod_{i=1}^r\dd\mu(\lambda_i),
  \]
  \[
    P_{m,r}=\frac {m}{r}\Bigl[
    h_n(p^{\times r},-\lambda_1,\ldots,-\lambda_r)
    +h_n(q^{\times r},\lambda_1,\ldots,\lambda_r)\Bigr]
    -h_{n-1}(q^{\times(r+1)},\lambda_1,\ldots,\lambda_r),
  \]
  where $h_d$ is the complete homogeneous symmetric polynomial of degree $d$ and $a^{\times k}$ means $a$ repeated $k$ times.
  It therefore suffices to prove $P_{m,r}\ge0$ on $[-1/2,1/2]^r$.

  That polynomial is then simplified three times until it can be shown nonnegative with elementary methods.
  First, using the fact that a complete homogeneous polynomial is an expectation of a power of a random convex combination, so with $\Theta$ uniform on the simplex $\Delta^r$ and $\widetilde P_{m,r}(q,\ell)=P_{m,r}(q;\ell,\ldots,\ell)$ (\Cref{lem:pge23-dirichlet}),
  \[
    P_{m,r}(q;\lambda_1,\ldots,\lambda_r)
    =\E\,\widetilde P_{m,r}
    \Bigl(q,\sum_{i=1}^r\Theta_i\lambda_i\Bigr).
  \]
  Every such convex combination again lies in $[-1/2,1/2]$, so the $r$ variables collapse to one and it suffices to prove $\widetilde P_{m,r}(q,\ell)\ge0$ for $\ell\in[-1/2,1/2]$.
  Second, once we obtain this diagonal polynomial, we can simplify it further to an integral against a beta distribution, which for $\ell>0$ becomes (\Cref{lem:pge23-beta})
  \begin{align*}
    \widetilde P_{m,r}(q,\ell) &=\frac{\Gamma(m)}{r\Gamma(n)\Gamma(r)^2}\,\ell^{n+r}
    \int_0^\infty\frac{s^{r-1}}{(\ell+s)^m}
    \rho_{n,m}(q+s)\dd s,
    \\
    \rho_{n,m}(u)&=\frac {m}{n}\bigl(u^n+(1-u)^n\bigr)-u^{n-1}.
  \end{align*}
  The case $\ell\le0$ is proved separately in \Cref{lem:pge23-ell-negative}, where the argument of $\rho_{n,m}$ stays at most $q\le1/2$ and $\rho_{n,m}$ is nonnegative there.
  As $\rho_{n,m}$ is not nonnegative at every real number, the integrand cannot be handled pointwise, and we estimate the average instead.
  Third, the Laplace representation of $(\ell+s)^{-m}$ replaces that average by an expectation over a gamma variable $Y\sim\GammaLaw(r,1)$ (\Cref{lem:pge23-laplace-gamma}):
  \[
    \widetilde P_{m,r}(q,\ell)
    =\frac{\ell^{n+r}}{r\Gamma(n)\Gamma(r)}
    \int_0^\infty t^{m-r-1}e^{-\ell t}\,
    \E\,\rho_{n,m}\Bigl(q+\frac {Y}{t}\Bigr)\dd t.
  \]
  As all the remaining terms are nonnegative, it is enough to prove that $\E\rho_{n,m}(q+xY)\ge0$ for every $x\ge0$.
  Centering at $u=1/2$ and using that $n$ is odd, this reduces to a single inequality comparing consecutive shifted gamma moments,
  \[
    3j\,\E(zY-1)^{2j-1}\le(r+j)\,\E(zY-1)^{2j}
    \qquad(z\ge0),
  \]
  proved in \Cref{lem:pge23-gamma-moment} from a Stein identity and the differential equation it induces.
  The constant $3$ is where the assumption $p\ge2/3$ is used.

  \item In this section, we prove \Cref{thm:main} for the case $1/2<p<2/3$ and odd $m\ge9$.
  The main tool in this section is mostly spectral analysis of $A$.
  We start from the identity for $W$:
  \[
    t(C_m,W)=p^m-\Tr(A^m)+m[z^m]L_-(z),
  \]
  together with the nonnegativity of every coefficient of $L_-$ (\Cref{lem:plt23-one-sided-identity}).

  We first consider the eigenvalues of $A$ larger than $q$.
  If there were two of them, then $\sum_i\lambda_i^2>2q^2$, whereas $\sum_i\lambda_i^2\le pq$ by \Cref{lem:plt23-HS-bound}, a contradiction because $2q^2 > pq$ when $q>1/3$.
  Hence at most one eigenvalue exceeds $q$.
  If none does the theorem is immediate (\Cref{lem:plt23-eigenvalues-above-q}), so we may assume that exactly one does and write $\Lambda$ for it.
  The same bound gives $|\lambda_i|\le M:=\sqrt{pq-\Lambda^2}$ for every other eigenvalue.
  One can also derive an upper bound on $\Lambda$ by using its eigenfunction $\phi$: since $U\ge0$, replacing $\phi$ by $|\phi|$ can only increase $\ip\phi{T_U\phi}$, and expanding $|\phi|$ along $\one$ and $\phi$ gives
  \[
    \norm g^2\ge\frac{\Lambda(\Lambda-q)^2}{2(1-2\Lambda)},
  \]
  which with $2\norm g^2\le pq-\Lambda^2$, again from \Cref{lem:plt23-HS-bound}, gives $\Lambda^2+q\Lambda-q\le0$ (\Cref{lem:plt23-forced-variance}).

  Decomposing $g=\gamma\phi+g_s$ into its component along $\phi$ and its component in the remaining spectral subspace, the two nonconstant terms of the identity can be bounded.
  For $\Tr(A^m)$, the eigenvalues other than $\Lambda$ satisfy $\sum\lambda_i^m\le M^{m-2}\sum\lambda_i^2$, so \Cref{lem:plt23-HS-bound} gives
  \[
    \Tr(A^m)\le\Lambda^m+M^m-2M^{m-2}\bigl(\gamma^2+\norm{g_s}^2\bigr).
  \]
  For the generating function, keeping only the first term of $-\log(1-X)=\sum_{k\ge1}X^k/k$ discards nonnegative coefficients, and leaves
  \[
    m[z^m]L_-(z)\ge m\int k_m\dd\mu
    \ge mk_m(\Lambda)\gamma^2+mk_m(M)\norm{g_s}^2,
    \qquad
    k_m(\lambda)=\frac{p^{m-1}-\lambda^{m-1}}{p+\lambda}.
  \]
  Together they give (\Cref{prop:plt23-splitting-lower-bound})
  \[
    t(C_m,W)-\bigl(p^m-pq^{m-1}\bigr)
    \ge-R_m+\Lambda_m\gamma^2+M_m\norm{g_s}^2,
  \]
  where $R_m=\Lambda^m+M^m-pq^{m-1}$, $\Lambda_m=2M^{m-2}+mk_m(\Lambda)$ and $M_m=2M^{m-2}+mk_m(M)$, with $0<\Lambda_m<M_m$ because $k_m$ decreases and $M<\Lambda$.
  It therefore suffices to prove $\Lambda_m\gamma^2+M_m\norm{g_s}^2\ge R_m$.

  To prove this lower bound, we first derive bounds on $\gamma$ and $\norm{g_s}$ from the eigenfunction $\phi$.
  Decompose
  \[
    |\phi|=a_\phi\one+b\phi+k,
    \qquad a_\phi=\int_0^1|\phi|,
    \qquad k\perp\one,\phi,
  \]
  and set $K=\norm k^2$, so that $a_\phi^2+b^2+K=1$.
  One can then show (\Cref{lem:plt23-gamma-upper-bound,lem:plt23-gs-lower-bound}) that
  \[
    \gamma\le\frac{a_\phi^2-2\Lambda(1+b)}{2a_\phi},
    \qquad
    b\gamma+\ip{g_s}k\ge
    \frac{(\Lambda-q)a_\phi^2+(\Lambda-M)K}{2a_\phi},
  \]
  and the Cauchy--Schwarz inequality $\ip{g_s}k\le\norm{g_s}\sqrt K$ turns the second inequality into a lower bound on $\norm{g_s}^2$.
  We then solve the constrained minimization of $\Lambda_m\gamma^2+M_m\norm{g_s}^2$ under these two bounds.
  Minimizing first in $\gamma$, and then over the three numbers $a_\phi,b,K$ that describe $\phi$, gives (\Cref{thm:plt23-huber-elim})
  \[
    \Lambda_m\gamma^2+M_m\norm{g_s}^2\ge \Omega_m\psi(\xi,\rho),
    \qquad
    \psi(\xi,\rho)=\min_{0\le v\le1}
    \bigl\{\rho v^2+\pos{\xi-v+v^2}\bigr\},
  \]
  where $\Omega_m$, $\xi$, and $\rho$ depend on $q$, $\Lambda$, and $m$.
  The theorem in this region is thereby reduced to the scalar inequality $R_m\le \Omega_m\psi(\xi,\rho)$.
  Since $\pos x=\max_{0\le\lambda\le1}\lambda x$, the definition of $\psi$ is a nested optimization,
  \[
    \psi(\xi,\rho)=\min_{0\le v\le1}\ \max_{0\le\lambda\le1}
    \bigl\{\rho v^2+\lambda(\xi-v+v^2)\bigr\}.
  \]
  We interchange the minimum and the maximum by the strong duality of this problem and carry out the inner minimization in $v$, which gives (\Cref{prop:plt23-huber-dual})
  \[
    \psi(\xi,\rho)=\max_{0\le\lambda\le1}
    \left\{\lambda\xi-\frac{\lambda^2}{4(\rho+\lambda)}\right\}.
  \]
  Rather than choosing one value of $\lambda$, we keep both scales that occur in this maximum.
  Write
  \[
    \mathcal A=\frac{\Omega_m\rho\xi^2}{\Lambda^m},
    \qquad
    \mathcal B=\frac{\Omega_m\xi}{\Lambda^m},
    \qquad
    \theta=\rho\xi=\frac{\mathcal A}{\mathcal B}.
  \]
  Dividing the objective by $\xi$ and discarding $\lambda\xi$ from the denominator leaves a maximum in which only $\theta$ occurs,
  \[
    \frac{\Omega_m\psi(\xi,\rho)}{\Lambda^m}\ge\mathcal B\,d(\theta),
    \qquad
    d(\theta)=\max_{0\le\lambda\le1}\left\{\lambda-\frac{\lambda^2}{4\theta}\right\},
  \]
  and $d(\theta)$ equals $\theta$ for $\theta\le1/2$ and $1-1/(4\theta)$ for $\theta\ge1/2$.

  The goal $R_m\le \Omega_m\psi(\xi,\rho)$ involves only $q$, $\Lambda$, and $m$.
  From $p=1-q$ and $M^2=pq-\Lambda^2$, the two ratios
  \[
    v=\frac {p}{\Lambda}-1,
    \qquad
    \zeta=\frac{M^2}{\Lambda(\Lambda-q)}
  \]
  determine $q$, $p$, $\Lambda$, and $M$, so we may replace $q$ and $\Lambda$ by $v$ and $\zeta$.
  For every $m$, the pairs $(v,\zeta)$ to be considered are those satisfying
  \[
    0<v<1,
    \qquad
    \zeta(\zeta+v)\ge1
  \]
  (\Cref{subsec:plt23-variables}).
  Writing $F_N=R_m/\Lambda^m$, what has to be proved is that $F_N\le\mathcal B\,d(\theta)$ holds in that region for every odd $m\ge9$.

  The region is divided only at $v=1/2$, and on each half two bounds are proved at once.
  \begin{itemize}[leftmargin=1.6em,topsep=2pt,itemsep=1pt]
    \item For $v\le1/2$: $2F_N\le\mathcal A$ and $F_N/\mathcal A+2F_N/\mathcal B\le2$ (\Cref{prop:plt23-low-v}).
    \item For $v\ge1/2$: $F_N\le\mathcal A$ and $2F_N\le\mathcal B$ (\Cref{prop:plt23-high-v-A,prop:plt23-high-v-B}).
  \end{itemize}
  In both halves the first bound settles $\theta\le1/2$ and the second settles $\theta\ge1/2$, which gives $F_N\le\mathcal B\,d(\theta)$ (\Cref{prop:plt23-scalar}).
  All four bounds are proved by elementary analysis; the only inputs are two upper bounds for $F_N$, one bound for a product of elementary factors, and one comparison between a geometric sum and a power.
\end{enumerate}

We prove \Cref{thm:main} in \Cref{sec:assembly}, combining the parts above.
Separately from the proof, \Cref{sec:transferable} discusses the possible applicability of the methods used here to other problems in extremal combinatorics.

\section{The elementary cases \texorpdfstring{$m=3,5,7$}{m=3,5,7}}\label{sec:short-cycles}

Throughout this section, $W$ is any graphon, $p=t(K_2,W)$,
\[
  U=1-W,\qquad q=t(K_2,U)=1-p,
\]
and $T=T_U$ is the integral operator of $U$.
When $q>1/2$, the right side of \eqref{eq:main-goal} is nonpositive, so the result is trivial.
We may therefore assume
\begin{equation}\label{eq:q-short}
  0\le q\le\frac{1}{2}.
\end{equation}
The arguments below actually prove the three cases at every edge density.

\subsection{The triangle: Goodman's one-line argument}\label{subsec:C3}

The following is the graphon form of Goodman's classical counting argument~\cite{Goodman1959}.

\begin{proposition}\label{prop:C3}
  For every graphon $W$ with edge density $p=t(K_2,W)$,
  \[
    t(C_3,W)\ge p^3-p(1-p)^2.
  \]
\end{proposition}

\begin{proof}
  The target polynomial is
  \[
    p^3-p(1-p)^2=2p^2-p.
  \]
  Let $d_W(x)=\int_0^1W(x,y)\dd y$.
  Since $ab\ge a+b-1$ for $a,b\in[0,1]$,
  \[
    \int_0^1W(x,z)W(y,z)\dd z\ge d_W(x)+d_W(y)-1.
  \]
  Multiplying by $W(x,y)\ge0$, integrating in $(x,y)$, and applying Jensen's inequality gives
  \begin{align*}
    t(C_3,W)
    &\ge \int W(x,y)\bigl(d_W(x)+d_W(y)-1\bigr)\dd x\dd y\\
    &=2\int_0^1d_W(x)^2\dd x-p
    \ge2p^2-p,
  \end{align*}
  as desired.
\end{proof}

\subsection{Path moments in the complement}\label{subsec:path-moments}

Write
\begin{equation}\label{eq:T1-decomp}
  T\one=q\one+g,
  \qquad g\perp\one.
\end{equation}
This decomposition has the explicit pointwise form
\[
  g(x)=\int_0^1U(x,y)\dd y-q
  =\bigl(1-d_W(x)\bigr)-(1-p)
  =p-d_W(x).
\]
Thus $g$ is the negative of the deviation of the $W$-degree from its mean $p$.
For $j\ge1$, let $P_j$ denote the path with $j$ edges and put
\[
  x_j=t(P_j,U)=\ip{\one}{T^j\one}.
\]
Thus $x_1=q$.
The following first moments will be used for $C_5$.
Note that $T$ is a symmetric operator, so $\ip{f}{Tg}=\ip{Tf}{g}$ for all $f,g$.

\begin{lemma}\label{lem:path-to-4}
  With \eqref{eq:T1-decomp},
  \begin{align*}
    x_2&=q^2+\norm g^2,\\
    x_3&=q^3+2q\norm g^2+\ip g{Tg},\\
    x_4&=q^4+3q^2\norm g^2+2q\ip g{Tg}+\norm{Tg}^2.
  \end{align*}
\end{lemma}

\begin{proof}
  The first identity is
  \[
    x_2 = \langle T\one, T\one\rangle = \langle q\one + g, q\one + g\rangle = q^2 \norm{\one}^2 + \norm{g}^2 = q^2 + \norm{g}^2.
  \]
  Next, $T^2\one=T(q \one + g) = q^2\one + qg + Tg$, while
  \[
    \ip{\one}{Tg}=\ip{T\one}{g}= \ip{q\one + g}{g}=\norm g^2.
  \]
  Expanding $x_3$ and $x_4$ similarly gives the other two identities:
  \begin{align*}
    x_3 &= \ip{T^2\one}{T\one} = \ip{q^2\one + qg + Tg}{q\one + g} \\
    {} &= \ip{q^2\one}{q\one} + \ip{q^2\one}{g} + \ip{qg}{q\one} + \ip{qg}{g} + \ip{Tg}{q\one} + \ip{Tg}{g} \\
    {} &= q^3 + q \norm{g}^2 + q \ip{g}{T \one} + \ip{Tg}{g} \\
    {} &= q^3 + q \norm{g}^2 + q \ip{g}{q \one + g} + \ip{Tg}{g} \\
    {} &= q^3 + 2q \norm{g}^2 + \ip{g}{Tg},
    \\
    x_4 &= \ip{T^2 \one}{T^2 \one} = \ip{q^2\one + qg + Tg}{q^2\one + qg + Tg} \\
    {} &= \ip{q^2\one}{q^2\one} + 2 \ip{q^2\one}{qg} + 2 \ip{q^2\one}{Tg} + \ip{qg}{qg} + 2 \ip{qg}{Tg} + \ip{Tg}{Tg} \\
    {} &= q^4 + 2q^2 \ip{T \one}{g} + q^2 \norm{g}^2 + 2q \ip{g}{Tg} + \norm{Tg}^2 \\
    {} &= q^4 + 2q^2 \ip{q \one + g}{g} + q^2 \norm{g}^2 + 2q \ip{g}{Tg} + \norm{Tg}^2 \\
    {} &= q^4 + 3q^2 \norm{g}^2 + 2q \ip{g}{Tg} + \norm{Tg}^2.
  \end{align*}
\end{proof}

For $C_7$, let $H_0=\one^\perp$ and
\[
  A=P_{H_0}TP_{H_0}.
\]
Concretely, $H_0$ is the space of mean-zero functions, and its orthogonal projection is the centering operation
\[
  (P_{H_0}f)(x)=f(x)-\int_0^1f(y)\dd y.
\]
Thus $P_{H_0}$ removes the constant part of $f$, and $A$ is the restriction of $T$ to $H_0$, in the sense that $\ip f{Ah}=\ip f{Th}$ for all $f,h\in H_0$.
Define
\begin{equation}\label{eq:sj-def}
  s_j=\ip g{A^jg}\qquad(j\ge0).
\end{equation}
The kernel of $T$ is square-integrable, so $A$ is a compact self-adjoint operator and $H_0$ has an orthonormal basis of its eigenfunctions.
Write $\lambda_i$ for the distinct eigenvalues of $A$ and $P_i$ for the orthogonal projection onto the eigenspace of $\lambda_i$, so that $\sum_iP_i$ is the identity on $H_0$ and $A^j=\sum_i\lambda_i^jP_i$.
Then
\[
  s_j=\sum_i\lambda_i^j\norm{P_ig}^2,
\]
which is the $j$-th moment of the finite positive measure
\begin{equation}\label{eq:mu-short}
  \mu=\sum_i\norm{P_ig}^2\,\delta_{\lambda_i}.
\end{equation}
Rewriting $s_j$ in terms of this measure,
\begin{equation}\label{eq:sj-moment}
  s_j=\int\lambda^j\dd\mu(\lambda),
  \qquad s_0=\norm g^2=\mu(\R).
\end{equation}

We now express the path densities $x_j$ through $q$ and the moments $s_j$.
Split $T^n\one$ into its constant and mean-zero parts.
The constant part is $\ip\one{T^n\one}\one=x_n\one$, so
\[
  T^n\one=x_n\one+h_n,\qquad h_n\perp\one,
\]
where $x_0=1$ and $h_0=0$.
Now \eqref{eq:T1-decomp} and the definition of $A$ give the recursive formulas
\begin{equation}\label{eq:path-recursion}
  x_{n+1}=qx_n+\ip g{h_n},
  \qquad h_{n+1}=x_ng+Ah_n.
\end{equation}
For clarity, iterating the second recursion first gives
\begin{align*}
  h_1={}&g,\\
  h_2={}&qg+Ag,\\
  h_3={}&x_2g+qAg+A^2g,\\
  h_4={}&x_3g+x_2Ag+qA^2g+A^3g,\\
  h_5={}&x_4g+x_3Ag+x_2A^2g+qA^3g+A^4g,\\
  h_6={}&x_5g+x_4Ag+x_3A^2g+x_2A^3g
  +qA^4g+A^5g.
\end{align*}
More generally,
\[
  h_n=\sum_{j=0}^{n-1}x_{n-1-j}A^jg.
\]
If we take the inner product with $g$ and use the definition \eqref{eq:sj-def} of $s_j$,
\[
  \ip g{h_n}
  =\sum_{j=0}^{n-1}x_{n-1-j}\ip g{A^jg}
  =\sum_{j=0}^{n-1}x_{n-1-j}s_j.
\]
Substituting this into the first recursion of \eqref{eq:path-recursion} eliminates $h_n$, and leaves a recursion for the path densities alone,
\begin{equation}\label{eq:path-scalar-recursion}
  x_{n+1}=qx_n+\sum_{j=0}^{n-1}x_{n-1-j}s_j
  \qquad(n\ge1).
\end{equation}
We expand the next few terms that are used in the $C_5$ and $C_7$ calculations:
\begin{align}
  x_2={}& q^2+s_0, \notag\\
  x_3={}& q^3+2qs_0+s_1, \notag\\
  x_4={}& q^4+3q^2s_0+2qs_1+s_0^2+s_2, \notag\\
  x_5={}& q^5+4q^3s_0+3q^2s_1+3qs_0^2+2qs_2+2s_0s_1+s_3, \notag\\
  x_6={}& q^6+5q^4s_0+4q^3s_1+6q^2s_0^2+3q^2s_2
  +6qs_0s_1+2qs_3 \notag\\
  &\qquad{}+s_0^3+2s_0s_2+s_1^2+s_4. \label{eq:path-2-6}
\end{align}

\subsection{The pentagon: one completion of a square}\label{subsec:C5}

\begin{proposition}\label{prop:C5}
  For every graphon $W$ with edge density $p=t(K_2,W)$,
  \[
    t(C_5,W)\ge p^5-p(1-p)^4.
  \]
\end{proposition}

\begin{proof}
  Assume \eqref{eq:q-short}.
  Label the integration variables cyclically by $z_1,\ldots,z_5$, put $z_6=z_1$, and abbreviate
  \[
    U_i=U(z_i,z_{i+1})\qquad(1\le i\le5).
  \]
  Since $W=1-U$, the integrand defining $t(C_5,W)$ is $\prod_{i=1}^5(1-U_i)$.
  Ordinary inclusion--exclusion gives the exact pointwise expansion
  \begin{align*}
    \prod_{i=1}^5(1-U_i)
    ={}&1-\sum_{i=1}^5U_i
    +\sum_{1\le i<j\le5}U_iU_j
    -\sum_{1\le i<j<k\le5}U_iU_jU_k\\
    &+\sum_{i=1}^5\prod_{j\ne i}U_j
    -\prod_{i=1}^5U_i.
  \end{align*}
  We now integrate this identity over $(z_1,\ldots,z_5)$.
  Each single-edge term has integral $q=x_1$, so the first sum contributes $5q$.
  Among the $\binom52=10$ pairs of edges, five pairs are adjacent around the cycle; each forms a two-edge path and hence has integral $x_2$.
  The other five pairs are disjoint edges.
  Their variables separate, so each has integral $x_1^2=q^2$.
  Therefore
  \[
    \int\sum_{i<j}U_iU_j=5x_2+5q^2.
  \]

  Similarly, among the $\binom53=10$ triples, five consist of three consecutive edges and form a path with three edges, contributing $x_3$ each.
  Each of the other five consists of a two-edge path together with one disjoint edge.
  Homomorphism density is multiplicative over disjoint unions, so each of these contributes $x_2x_1=qx_2$.
  Thus
  \[
    \int\sum_{i<j<k}U_iU_jU_k=5x_3+5qx_2.
  \]
  There are five ways to select four edges; deleting any one edge from the cycle leaves a four-edge path, so every corresponding integral is $x_4$.
  Finally, selecting all five edges gives $c_5:=t(C_5,U)$.
  Combining these contributions with their alternating signs yields
  \begin{align*}
    t(C_5,1-U)
    ={}&1-5q+5x_2+5q^2-5x_3-5qx_2+5x_4-c_5.
  \end{align*}
  Since $0\le U\le1$, deleting one edge from the cycle integrand gives $c_5\le x_4$.
  Therefore
  \begin{equation}\label{eq:C5-IE}
    t(C_5,W)\ge1-5q+5q^2+5(1-q)x_2-5x_3+4x_4.
  \end{equation}
  Subtracting the lower bound
  \[
    (1-q)^5-(1-q)q^4=1-5q+10q^2-10q^3+4q^4
  \]
  from \eqref{eq:C5-IE} and using \Cref{lem:path-to-4} gives
  \begin{align}
    \Phi_5={}&4\norm{Tg}^2+(8q-5)\ip g{Tg}
    +(12q^2-15q+5)\norm g^2 \notag\\
    ={}&4\left\|Tg+\frac{8q-5}{8}g\right\|^2
    +\left(8q^2-10q+\frac{55}{16}\right)\norm g^2. \label{eq:Phi5}
  \end{align}
  The quadratic coefficient $8q^2-10q+\frac{55}{16}$ has discriminant
  \[
    100-4\cdot8\cdot\frac{55}{16}=-10<0
  \]
  and positive leading coefficient, so it is positive for every $q$.
  Thus $\Phi_5\ge0$, proving the claim.
\end{proof}

\subsection{The heptagon: a univariate sum of squares}\label{subsec:C7}

\begin{proposition}\label{prop:C7}
  For every graphon $W$ with edge density $p=t(K_2,W)$,
  \[
    t(C_7,W)\ge p^7-p(1-p)^6.
  \]
\end{proposition}

\begin{proof}
  Write $U=1-W$ and $q=1-p=t(K_2,U)$, and let $x_j=t(P_j,U)$ be the path densities of \eqref{eq:sj-def} above, so that $x_1=q$.
  If $p<1/2$, the right side of the claim is nonpositive and there is nothing to prove, so assume $p\ge1/2$, that is $q\le1/2$.
  As in the proof for $C_5$, label the integration variables cyclically by $z_1,\ldots,z_7$, put $z_8=z_1$, and write
  \[
    U_i=U(z_i,z_{i+1})\qquad(1\le i\le7).
  \]
  Expanding the seven factors $1-U_i$ gives the exact inclusion--exclusion identity
  \begin{equation}\label{eq:C7-subset-expansion}
    t(C_7,W)
    =\sum_{S\subseteq\{1,\ldots,7\}}(-1)^{|S|}
    \int_{[0,1]^7}\prod_{i\in S}U_i\prod_{j=1}^7\dd z_j.
  \end{equation}
  A proper subset $S$ selects a disjoint union of paths around the cycle, and its variables separate between the components, so its term is the corresponding product of path densities; the full subset gives $c_7=t(C_7,U)$.
  Collecting the subsets of each size, \eqref{eq:C7-subset-expansion} expands to
  \begin{align*}
    t(C_7,W)={}&1\\
    &-7x_1\\
    &+7x_2+14x_1^2\\
    &-7x_3-21x_1x_2-7x_1^3\\
    &+7x_4+14x_1x_3+7x_2^2+7x_1^2x_2\\
    &-7x_5-7x_1x_4-7x_2x_3\\
    &+7x_6\\
    &-c_7.
  \end{align*}
  The seven subsets of size one each contribute $x_1$; of the $21$ subsets of size two, seven consist of two adjacent edges and contribute $x_2$ while the other fourteen contribute $x_1^2$; and so on, the coefficients in each line summing to $\binom7k$.
  Because $0\le U_i\le1$, deleting one factor from the integrand of $c_7$ can only increase it; the remaining six factors form a six-edge path.
  Hence
  \[
    c_7=\int\prod_{i=1}^7U_i
    \le\int\prod_{i=1}^6U_i=x_6.
  \]
  Replacing $-c_7$ by $-x_6$ in the preceding identity and collecting the coefficients of the same path densities gives
  \begin{align}
    t(C_7,W)\ge{}&6x_6-7x_5+7(1-q)x_4+7(2q-1)x_3 \notag\\
    &+7(q^2-3q+1)x_2+7x_2^2-7x_2x_3 \notag\\
    &+1-7q+14q^2-7q^3. \label{eq:C7-IE}
  \end{align}
  Since
  \[
    (1-q)^7-(1-q)q^6
    =1-7q+21q^2-35q^3+35q^4-21q^5+6q^6,
  \]
  substitution of \eqref{eq:path-2-6} reduces the desired inequality to the nonnegativity of the polynomial
  \begin{align}
    \Phi_7={}&6s_4+(12q-7)s_3+(18q^2-21q+7)s_2 \notag\\
    &+(24q^3-42q^2+28q-7)s_1 \notag\\
    &+(30q^4-70q^3+70q^2-35q+7)s_0 \notag\\
    &+12s_0s_2+(36q-21)s_0s_1
    +(36q^2-42q+14)s_0^2+6s_0^3+6s_1^2. \label{eq:Phi7-expanded}
  \end{align}
  Write $r$ for the total mass $s_0=\mu(\R)$.
  Keeping one factor $r$ in each product and expressing every remaining $s_j$ by \eqref{eq:sj-moment} as an integral against $\mu$, the polynomial \eqref{eq:Phi7-expanded} becomes
  \begin{equation}\label{eq:Phi7-integral}
    \Phi_7=6s_1^2+\int R_{q,r}(\lambda)\dd\mu(\lambda),
  \end{equation}
  where
  \[
    R_{q,r}(\lambda)=P_q(\lambda)+rB_q(\lambda)+6r^2,
  \]
  \begin{align*}
    P_q(\lambda)={}&6\lambda^4+(12q-7)\lambda^3
    +(18q^2-21q+7)\lambda^2\\
    &+(24q^3-42q^2+28q-7)\lambda
    +30q^4-70q^3+70q^2-35q+7,\\
    B_q(\lambda)={}&12\lambda^2+(36q-21)\lambda+36q^2-42q+14.
  \end{align*}
  We show that both $P_q$ and $B_q$ are nonnegative.

  The minimum of $B_q$ as a quadratic in $\lambda$ occurs at $\lambda_*=(7-12q)/8$, and
  \[
    B_q(\lambda_*)=\frac{144q^2-168q+77}{16}\ge\frac{29}{16},
  \]
  because
  \[
    144q^2-168q+48=144\left(q-\frac{1}{2}\right)
    \left(q-\frac{2}{3}\right)\ge0
  \]
  on $[0,1/2]$.

  For $P_q$, set
  \begin{align*}
    D(q)&=288q^2-336q+119,\\
    C(q)&=24q^3-42q^2+28q-7,\\
    N(q)&=5184q^6-18144q^5+28602q^4-25802q^3\\
    &\qquad+13874q^2-4165q+539.
  \end{align*}
  Direct expansion gives the exact identity
  \begin{align}
    P_q(\lambda)
    ={}&6\left(\lambda^2+\left(q-\frac{7}{12}\right)\lambda\right)^2 \notag\\
    &+\frac{D(q)}{24}
    \left(\lambda+\frac{12C(q)}{D(q)}\right)^2
    +\frac{N(q)}{D(q)}. \label{eq:Pq-SOS}
  \end{align}
  Writing $y=1/2-q\ge0$,
  \[
    D(q)=288y^2+48y+23>0
  \]
  and
  \begin{align*}
    8N(q)={}&41472y^6+20736y^5+21456y^4+7984y^3\\
    &+2032y^2+316y+11>0.
  \end{align*}
  Thus $P_q(\lambda)\ge0$ for every real $\lambda$.
  Since $r\ge0$, $R_{q,r}(\lambda)\ge0$, and \eqref{eq:Phi7-integral} yields $\Phi_7\ge0$.
\end{proof}

The identity \eqref{eq:Pq-SOS} is an explicit sum-of-squares certificate for $P_q$.
For background on positivity of one-variable polynomials and on such certificates, see Powers--Reznick~\cite{PowersReznick2000}.

\iffalse
\begin{remark}[Pointwise nonnegativity versus operator positivity]
  The $C_5$ and $C_7$ arguments do not require the kernel operator $T_U$ to be positive semidefinite.
  Such positivity would mean that
  \[
    \ip f{T_Uf}
    =\int_{[0,1]^2}U(x,y)f(x)f(y)\dd x\dd y\ge0
    \qquad\text{for every }f\in L^2([0,1]).
  \]
  Although $U(x,y)\ge0$ pointwise, this condition need not hold, because $f(x)f(y)$ may be negative.
  For example, let $U$ be the balanced complete bipartite graphon, equal to $1$ between the two halves of $[0,1]$ and to $0$ within each half.
  If $f$ equals $1$ on the first half and $-1$ on the second, then $T_Uf=-\tfrac{1}{2}f$, and hence $\ip f{T_Uf}=-\tfrac{1}{2}\norm f^2<0$.

  What the proofs use instead is self-adjointness, which follows from the symmetry of $U$.
  The spectral theorem then associates with $g$ a nonnegative measure $\mu$ satisfying $s_j=\int\lambda^j\dd\mu(\lambda)$.
  Here nonnegative refers to the weights of the measure; its support may still contain negative eigenvalues $\lambda$.
  This is why the proof above shows that the relevant polynomials are nonnegative for every real $\lambda$, rather than only for $\lambda\ge0$.
  No positive-semidefiniteness assumption on $T_U$ is needed.
\end{remark}

\fi

\section{Reductions and notation used in both regions}\label{sec:common-operator}

\subsection{Step-graphon reduction}\label{subsec:step-reduction}

As described in the proof map above, the cycle densities appear as the coefficients of $-\log\det(I-zT_W)$.
To avoid taking determinants of operators on an infinite-dimensional space, we reduce the theorem to step graphons, whose operators have finite rank.
The approximation is the dyadic conditional expectation of \cite[Appendix~A.3.6]{Lovasz2012}, which has the property we need here: the approximating graphons have edge density exactly $p$.

\begin{lemma}[Finite-rank reduction]\label{lem:step-reduction}
  For each fixed odd $m\ge3$, it is enough to prove \eqref{eq:main-goal} for symmetric step graphons.
\end{lemma}

\begin{proof}
  For $n\ge1$, divide $[0,1]$ into the $2^n$ dyadic intervals
  \[
    I_{n,a}=\left[\frac{a}{2^n},\frac{a+1}{2^n}\right),
    \qquad 0\le a<2^n,
  \]
  with the convention that the last interval contains the endpoint $1$.
  Let $\mathcal A_n$ be the finite sigma-algebra consisting of all unions of these intervals.
  Thus the atoms of the product sigma-algebra $\mathcal A_n\otimes\mathcal A_n$ are precisely the dyadic rectangles $I_{n,a}\times I_{n,b}$.
  The partitions refine as $n$ increases, and dyadic intervals generate the Borel sigma-algebra.

  On each such rectangle, replace $W$ by its average.
  More explicitly, set
  \begin{equation}\label{eq:dyadic-step-approximation}
    w^{(n)}_{ab}
    =2^{2n}\int_{I_{n,a}\times I_{n,b}}W(x,y)\dd x\dd y,
    \qquad
    W_n(x,y)=w^{(n)}_{ab}
    \quad\text{if }(x,y)\in I_{n,a}\times I_{n,b}.
  \end{equation}
  This cell-averaging operation is the standard conditional expectation operation with respect to the sigma algebra $\mathcal A_n\otimes\mathcal A_n$, and is denoted by the notation $\E[\,-\, \mid\mathcal A_n\otimes\mathcal A_n]$.
  So, $W_n = \E[W\mid\mathcal A_n\otimes\mathcal A_n] $.
  Since $0\le W\le1$, every cell average lies in $[0,1]$.
  Moreover, symmetry of $W$ gives $w^{(n)}_{ab}=w^{(n)}_{ba}$.
  Hence $W_n$ is a symmetric step graphon.

  There is also a completely finite interpretation.
  The matrix $(w^{(n)}_{ab})_{0\le a,b<2^n}$ is the weighted adjacency matrix of a graph on $2^n$ equally weighted vertices (loops are allowed), and
  \[
    t(C_m,W_n)
    =\frac{1}{2^{nm}}
    \sum_{a_1,\ldots,a_m=0}^{2^n-1}
    w^{(n)}_{a_1a_2}w^{(n)}_{a_2a_3}\cdots
    w^{(n)}_{a_ma_1}.
  \]
  The same finite structure can be seen directly at the operator level.
  For any function $f$, if $x\in I_{n,a}$, then
  \[
    (T_{W_n}f)(x)
    =\int_0^1W_n(x,y)f(y)\dd y
    =\sum_{b=0}^{2^n-1}w^{(n)}_{ab}
    \int_{I_{n,b}}f(y)\dd y.
  \]
  Thus the output $T_{W_n}f$ is constant on each interval $I_{n,a}$, no matter how complicated the input $f$ is.
  Every such output is therefore a linear combination of the $2^n$ indicator functions $\one_{I_{n,0}},\ldots,\one_{I_{n,2^n-1}}$.
  The set of all possible outputs is called the \emph{range} of the operator, and its dimension is the \emph{rank}.
  Since the range is contained in a space of dimension $2^n$, the rank is at most $2^n$.
  This explains the term ``finite-rank reduction.''

  Two facts remain to be checked: the edge density is unchanged, and the cycle density converges.
  For the first, summing the averages in \eqref{eq:dyadic-step-approximation} over all cells gives
  \[
    \int_{[0,1]^2}W_n(x,y)\dd x\dd y
    =\sum_{a,b}|I_{n,a}|\,|I_{n,b}|\,w^{(n)}_{ab}
    =\int_{[0,1]^2}W(x,y)\dd x\dd y=p.
  \]
  Thus $t(K_2,W_n)=p$ for every $n$, not merely in the limit.

  Next we verify that $W_n\to W$ in $L^2([0,1]^2)$.
  Dyadic rectangle step functions are dense in this space.
  Therefore, given $\varepsilon>0$, we may choose an integer $N$ and a function $S$, constant on every rectangle $I_{N,a}\times I_{N,b}$, such that $\norm{W-S}_2<\varepsilon$.
  For $n\ge N$, the function $S$ is also constant on every level-$n$ rectangle.
  Averaging is an $L^2$-contraction (this is Jensen's inequality on each rectangle), so
  \[
    \norm{W_n-S}_2
    =\norm{\E[W-S\mid\mathcal A_n\otimes\mathcal A_n]}_2
    \le \norm{W-S}_2.
  \]
  Consequently
  \[
    \norm{W_n-W}_2
    \le \norm{W_n-S}_2+\norm{S-W}_2
    <2\varepsilon,
  \]
  which proves the asserted convergence.

  Finally, cycle densities are continuous under this convergence.
  Put $x_{m+1}=x_1$, and for a fixed $(x_1,\ldots,x_m)$, write
  \[
    a_i=W_n(x_i,x_{i+1}),\qquad b_i=W(x_i,x_{i+1}).
  \]
  The telescoping identity
  \[
    \prod_{i=1}^m a_i-\prod_{i=1}^m b_i
    =\sum_{j=1}^m
    \left(\prod_{i<j}a_i\right)(a_j-b_j)
    \left(\prod_{i>j}b_i\right)
  \]
  and the bounds $0\le a_i,b_i\le1$ show, after integration, that
  \begin{align*}
    \bigl|t(C_m,W_n)-t(C_m,W)\bigr|
    &\le \sum_{j=1}^m
    \int_{[0,1]^m}
    |W_n(x_j,x_{j+1})-W(x_j,x_{j+1})|
    \prod_{i=1}^m\dd x_i \\
    &=m\norm{W_n-W}_1
    \le m\norm{W_n-W}_2
    \longrightarrow0.
  \end{align*}
  Here the equality in the second line follows because all variables other than $x_j,x_{j+1}$ integrate out to $1$.

  Now suppose \eqref{eq:main-goal} has been proved for every symmetric step graphon.
  Applying it to $W_n$, whose edge density is exactly $p$, gives
  \[
    t(C_m,W_n)\ge p^m-p(1-p)^{m-1}.
  \]
  Letting $n\to\infty$ proves the same inequality for $W$, as required.
\end{proof}

\subsection{The block decomposition and the support of the spectral measure}
\label{subsec:block-decomposition}

From here to the end of \Cref{sec:plt23}, we assume $W$ to be a step graphon, which \Cref{lem:step-reduction} permits, and all determinants and traces are taken in the resulting finite-dimensional setting.
Fix a finite measurable partition
\[
  [0,1]=B_1\sqcup\cdots\sqcup B_N
\]
into sets of positive measure such that $W$, and hence $U=1-W$, is constant on each rectangle $B_a\times B_b$.

Let
\[
  \mathcal E=\operatorname{span}
  \{\one_{B_1},\ldots,\one_{B_N}\}\subseteq L^2([0,1]).
\]
This is the $N$-dimensional space of functions that are constant on every part $B_a$.
If a step kernel $K$ has value $K_{ab}$ on $B_a\times B_b$, then, for $x\in B_a$,
\begin{equation}\label{eq:step-operator-action}
  (T_Kf)(x)=\sum_{b=1}^N K_{ab}\int_{B_b}f(y)\dd y.
\end{equation}
Consequently, $T_Kf\in\mathcal E$ for every $f\in L^2$.
Moreover, if $f\in\mathcal E^\perp$, then every integral in \eqref{eq:step-operator-action} is zero, so $T_Kf=0$.
Hence $T_K$ is determined by its restriction to the finite-dimensional space $\mathcal E$.

The cycle densities of $K$ are computed from this finite matrix as follows.
In the orthonormal basis $e_a=\one_{B_a}/\sqrt{|B_a|}$, the matrix of $T_K|_{\mathcal E}$ has entries
\[
  \ip{e_a}{T_Ke_b}=K_{ab}\sqrt{|B_a||B_b|}.
\]
Here $|B_a|$ denotes the measure of $B_a$.
Multiplying this matrix around a cycle shows, for every $r\ge2$, that
\begin{align}
  \Tr(T_K^r)
  &=\sum_{a_1,\ldots,a_r=1}^N
  |B_{a_1}|\cdots|B_{a_r}|
  K_{a_1a_2}\cdots K_{a_ra_1} \notag\\
  &=t(C_r,K). \label{eq:finite-trace-cycle}
\end{align}

We now split $\mathcal E$ into the constant direction and its orthogonal complement, and decompose $T_U$ accordingly.
The constant function $\one$ lies in $\mathcal E$ and has norm one, because the underlying probability space has total mass one, so setting
\[
  \mathcal E_0=\mathcal E\cap\one^\perp
\]
gives the orthogonal decomposition $\mathcal E=\R\one\oplus\mathcal E_0$.
Put
\[
  g=T_U\one-q\one\in\mathcal E_0,
  \qquad
  A=P_{\mathcal E_0}T_UP_{\mathcal E_0},
\]
so that $A$ is a self-adjoint map on $\mathcal E_0$, equivalently a real symmetric matrix in an orthonormal basis of that space.
These two determine $T_U$: for $f=a\one+h$ with $h\in\mathcal E_0$,
\[
  T_Uf=\bigl(aq+\ip gh\bigr)\one+\bigl(ag+Ah\bigr),
\]
since $\ip\one{T_Uh}=\ip{T_U\one}h=\ip gh$ and $P_{\mathcal E_0}T_Uh=Ah$.
In block form,
\begin{equation}\label{eq:block-U}
  T_U=
  \begin{pmatrix}
    q&g^*\\
    g&A
  \end{pmatrix},
\end{equation}
where $g^*$ denotes the linear functional $h\mapsto\ip gh$.
This $A$ plays the role of the operator in \Cref{subsec:path-moments}: for a step graphon, the compression $P_{\one^\perp}T_UP_{\one^\perp}$ used there is zero on $\mathcal E^\perp$, and its restriction to $\mathcal E_0$ is $A$.


\begin{lemma}[Operator norm bound]\label{lem:half-norm}
  The spectrum of $A$ lies in $[-1/2,1/2]$; equivalently,
  \begin{equation}\label{eq:A-half}
    \norm A_{\mathrm{op}}\le\frac{1}{2}.
  \end{equation}
\end{lemma}

\begin{proof}
  Let $f\in\mathcal E_0$ with $\norm f=1$, and write $f=f_+-f_-$ for nonnegative functions $f_+$ and $f_-$ on $[0,1]$.
  Since $\int f=0$,
  \[
    \int f_+=\int f_-=:a,
    \qquad 2a=\lVert f\rVert_1\le\lVert f\rVert_2=1.
  \]
  We also have
  \begin{align*}
    \ip f{T_Uf}
    ={}&\ip {f_+}{T_Uf_+}+\ip {f_-}{T_Uf_-}
    -2\ip {f_+}{T_Uf_-}.
  \end{align*}
  Each of the three inner products on the right is nonnegative.
  Moreover,
  \begin{align*}
    \ip {f_+}{T_Uf_+} &= \int_{[0,1]^2}U(x,y)f_+(x)f_+(y)\dd x\dd y \le \int_{[0,1]^2}f_+(x)f_+(y)\dd x\dd y = \left(\int f_+\right)^2,
    \\
    \ip {f_-}{T_Uf_-} &= \int_{[0,1]^2}U(x,y)f_-(x)f_-(y)\dd x\dd y \le \int_{[0,1]^2}f_-(x)f_-(y)\dd x\dd y = \left(\int f_-\right)^2,
    \\
    \ip {f_+}{T_Uf_-} &= \int_{[0,1]^2}U(x,y)f_+(x)f_-(y)\dd x\dd y \le \int_{[0,1]^2}f_+(x)f_-(y)\dd x\dd y = \left(\int f_+\right)\left(\int f_-\right).
  \end{align*}
  because $f_+$ and $f_-$ are nonnegative and so replacing $U(x,y)$ by $1$ in each of the above three cases gives the corresponding product of integrals.
  For the upper bound, we discard the nonpositive cross term and bound the two diagonal terms; hence
  \[
    \ip f{T_Uf}
    \le\left(\int f_+\right)^2+\left(\int f_-\right)^2
    =2a^2\le\frac{1}{2}.
  \]
  For the lower bound, we instead discard the two nonnegative diagonal terms and bound the cross term, obtaining
  \[
    \ip f{T_Uf}
    \ge-2\left(\int f_+\right)\left(\int f_-\right)
    =-2a^2\ge-\frac{1}{2}.
  \]
  The upper and lower bounds just proved concern the quadratic form of $T_U$ on $\mathcal E_0$, and this is also the quadratic form of $A$: since $P_{\mathcal E_0}f=f$,
  \[
    \ip f{Af}=\ip f{P_{\mathcal E_0}T_UP_{\mathcal E_0}f}=\ip f{T_Uf},
  \]
  so $\bigl|\ip f{Af}\bigr|\le\tfrac{1}{2}$ for every unit $f\in\mathcal E_0$.
  As $A$ is self-adjoint, its operator norm is the supremum of $\bigl|\ip f{Af}\bigr|$ over such $f$, whence $\norm A_{\mathrm{op}}\le\tfrac{1}{2}$.
\end{proof}

In the remaining section, $A$ and $g$ will appear through the quantities of the form $\ip g{F(A)g}$, which can be expressed as an integral against a measure supported on the spectrum of $A$.
Specifically, since $A$ is a finite real symmetric matrix, it has an orthogonal eigenspace decomposition
\[
  A=\sum_i\lambda_iP_i,
\]
where the distinct $\lambda_i$ are its eigenvalues and $P_i$ is the orthogonal projection onto the corresponding eigenspace.
Define the finite positive measure
\begin{equation}\label{eq:mu-def}
  \mu=\sum_i\norm{P_i g}^2\,\delta_{\lambda_i},
\end{equation}
where $\delta_{\lambda_i}$ denotes a unit point mass at $\lambda_i$; thus $\mu$ assigns to each eigenvalue the weight measuring how much of $g$ lies in its eigenspace, and by \Cref{lem:half-norm} it is supported on $[-1/2,1/2]$.
For any function $F$ defined on the eigenvalues of $A$, orthogonality of the eigenspaces gives
\[
  \ip g{F(A)g}
  =\sum_i F(\lambda_i)\norm{P_i g}^2
  =\int F(\lambda)\dd\mu(\lambda).
\]

\subsection{Interpreting the power-series
calculations}\label{subsec:analytic-reading}

Several later arguments introduce an auxiliary variable $z$ and use the notation $[z^m]F(z)$; the exponent of $z$ is the length of a walk or cycle.
We do not assign a value to $z$; every $F$ below is a formal power series, that is, a sequence of real coefficients $(c_j)_{j\ge0}$ written as $F=\sum_{j\ge0}c_jz^j$, and $[z^m]F=c_m$.
We say that two such series are equal when all their coefficients are equal.

Four operations on these series are used below, and in each of them every coefficient of the result is computed from finitely many coefficients of the arguments.
Series are added coefficientwise.
They are multiplied by the Cauchy product,
\[
  [z^m](FG)=\sum_{i+j=m}\bigl([z^i]F\bigr)\bigl([z^j]G\bigr).
\]
A series with $c_0\ne0$ has an inverse, obtained by solving the equations for the coefficients of $FF^{-1}=1$ recursively, so we may divide by such a series.
Finally, if we write $\ord F$ for the order of $F$, that is the least $j$ with $c_j\ne0$, then the series
\[
  -\log(1-Y)=\sum_{k\ge1}\frac{Y^k}{k}
\]
is defined whenever $\ord Y\ge1$: in that case $\ord Y^k\ge k$, so only the terms with $k\le m$ contribute to $[z^m]$.

We build our series from the finite matrices of \Cref{subsec:block-decomposition}, using the two expansions in the next lemma.

\begin{lemma}[Series expansions used below]\label{lem:series-expansions}
  Let $X$ be a real symmetric $N\times N$ matrix.
  As formal power series in $z$:
  \begin{enumerate}[label=\textup{(\arabic*)},leftmargin=2.6em]
    \item \emph{(Resolvent.)}
    The matrix $I-zX$ is invertible over $\R[[z]]$ and
    \[
      (I-zX)^{-1}=\sum_{j\ge0}z^jX^j,
    \]
    so that $[z^j]\ip u{(I-zX)^{-1}v}=\ip u{X^jv}$ for fixed vectors $u,v$.
    \item \emph{(Log-determinant.)}
    The polynomial $\det(I-zX)$ has constant term $1$, and
    \[
      -\log\det(I-zX)=\sum_{j\ge1}\frac{\Tr(X^j)}{j}z^j,
    \]
    so that $m\,[z^m]\bigl(-\log\det(I-zX)\bigr)=\Tr(X^m)$.
  \end{enumerate}
\end{lemma}

\begin{proof}
  (1)  The constant term of $I-zX$ is $I$, so the matrix is invertible over $\R[[z]]$; multiplying $\sum_{j\ge0}z^jX^j$ by $I-zX$ cancels all positive powers of $z$ and leaves $I$.

  (2)  Both sides are unchanged under $X\mapsto Q^{\!\top}XQ$ with $Q$ orthogonal, so we may take $X$ diagonal with entries $\lambda_i$.
  Then $\det(I-zX)=\prod_i(1-z\lambda_i)$, and $-\log$ of a product is the sum of the $-\log$s, whence
  \[
    -\log\det(I-zX)
    =\sum_{i=1}^N\sum_{j\ge1}\frac{\lambda_i^j}{j}z^j
    =\sum_{j\ge1}\frac{\Tr(X^j)}{j}z^j,
  \]
  the interchange being a finite sum in each coefficient.
\end{proof}

For a first example of these expansions, take $X=T_W$ and $m=3$.
Item~(2) gives
\[
  -\log\det(I-zT_W)
  =\Tr(T_W)\,z+\frac{\Tr(T_W^2)}{2}z^2+\frac{\Tr(T_W^3)}{3}z^3+\cdots,
\]
whence
\[
  3\,[z^3]\bigl(-\log\det(I-zT_W)\bigr)=\Tr(T_W^3)=t(C_3,W)
\]
by \eqref{eq:finite-trace-cycle}.

For a second example, expanding $(I-zA)^{-1}=\sum_{j\ge0}z^jA^j$ by \Cref{lem:series-expansions}\,(1) and taking $F(\lambda)=\lambda^j$ in each coefficient,
\begin{equation}\label{eq:resolvents-common}
  \ip g{(I-zA)^{-1}g}
  =\sum_{j\ge0}z^j\int\lambda^j\dd\mu(\lambda)
  =\int_{-1/2}^{1/2}\frac{\dd\mu(\lambda)}{1-\lambda z}.
\end{equation}
Similarly, we have
\begin{equation}\label{eq:resolvent-plus}
  \ip g{(I+zA)^{-1}g}
  =\int_{-1/2}^{1/2}\frac{\dd\mu(\lambda)}{1+\lambda z}.
\end{equation}

By \Cref{lem:series-expansions}\,(2) and \eqref{eq:finite-trace-cycle}, every step kernel $K$ satisfies
\begin{equation}\label{eq:logdet-cycle}
  -\log\det(I-zT_K)=\sum_{j\ge1}\frac{\Tr(T_K^j)}{j}z^j,
  \qquad
  m\,[z^m]\bigl(-\log\det(I-zT_K)\bigr)=t(C_m,K)
\end{equation}
for every $m\ge2$.
One can now simplify this generating function by the block decomposition \eqref{eq:block-U}, as the following lemma states.

\begin{lemma}[Cycle densities of $U$ and $W$]\label{lem:cycle-identities}
  For every odd $m\ge3$, let
  \begin{align}
    L_-(z)&=-\log\left(1-\frac{z^2\ip g{(I+zA)^{-1}g}}{1-pz}\right),
    \label{eq:Lminus}\\
    L_+(z)&=-\log\left(1-\frac{z^2\ip g{(I-zA)^{-1}g}}{1-qz}\right).
    \label{eq:Lplus}
  \end{align}
  Then
  \begin{align}
    t(C_m,U)&=\phantom{-}\Tr(A^m)+q^m+m[z^m]L_+(z),\label{eq:cycle-identity-U}\\
    t(C_m,W)&=-\Tr(A^m)+p^m+m[z^m]L_-(z).\label{eq:cycle-identity-W}
  \end{align}
\end{lemma}

\begin{proof}
  We first put the two matrices $I-zT_U$ and $I-zT_W$ in block form.
  The operator $T_1$ of the constant kernel $1$ sends $a\one+h$, with $h\in\mathcal E_0$, to $a\one$, because $\int h=0$; thus $T_1$ has block matrix $\left(\begin{smallmatrix}1&0\\0&0\end{smallmatrix}\right)$ and $T_W=T_1-T_U$ has block matrix
  \[
    T_W=
    \begin{pmatrix}1&0\\0&0\end{pmatrix}
    -
    \begin{pmatrix}q&g^*\\g&A\end{pmatrix}
    =
    \begin{pmatrix}p&-g^*\\-g&-A\end{pmatrix}.
  \]
  Consequently
  \[
    I-zT_U=
    \begin{pmatrix}1-qz&-zg^*\\-zg&I-zA\end{pmatrix},
    \qquad
    I-zT_W=
    \begin{pmatrix}1-pz&zg^*\\zg&I+zA\end{pmatrix}.
  \]

  Both determinants decompose along this block structure: if $D$ is invertible, then
  \begin{equation}\label{eq:block-determinant-identity}
    \det\begin{pmatrix}\alpha&v^*\\v&D\end{pmatrix}
    =\det(D)\bigl(\alpha-v^*D^{-1}v\bigr),
  \end{equation}
  where the scalar $\alpha-v^*D^{-1}v$ is the Schur complement of $D$.
  To see this, multiply on the left by the block matrix $\left(\begin{smallmatrix}1&-v^*D^{-1}\\0&I\end{smallmatrix}\right)$, whose determinant is one:
  \begin{align*}
    \det\begin{pmatrix}\alpha&v^*\\v&D\end{pmatrix}
    &=\underbrace{\det\begin{pmatrix}1&-v^*D^{-1}\\0&I\end{pmatrix}}_{=1}
      \det\begin{pmatrix}\alpha&v^*\\v&D\end{pmatrix}\\
    &=\det\left[\begin{pmatrix}1&-v^*D^{-1}\\0&I\end{pmatrix}
                \begin{pmatrix}\alpha&v^*\\v&D\end{pmatrix}\right]\\
    &=\det\begin{pmatrix}\alpha-v^*D^{-1}v&v^*-v^*D^{-1}D\\v&D\end{pmatrix}\\
    &=\det\begin{pmatrix}\alpha-v^*D^{-1}v&0\\v&D\end{pmatrix}
     =\det(D)\bigl(\alpha-v^*D^{-1}v\bigr),
  \end{align*}
  where the last equality holds since the determinant of a block triangular matrix is the product of its diagonal-block determinants.

  To apply \eqref{eq:block-determinant-identity} to $I-zT_U$ and $I-zT_W$, one needs to check that the lower right block $I\mp zA$ is invertible, and this follows from \Cref{lem:series-expansions}\,(1) since both of these terms have $I$ as their constant term.
  Taking $D=I\mp zA$ and $v=\mp zg$ factorizes the two determinants into three parts,
  \begin{align*}
    \det(I-zT_U)&=\det(I-zA)\bigl(1-qz-z^2\ip g{(I-zA)^{-1}g}\bigr)\\
    &=\det(I-zA)\,(1-qz)
    \left(1-\frac{z^2\ip g{(I-zA)^{-1}g}}{1-qz}\right),\\
    \det(I-zT_W)&=\det(I+zA)\bigl(1-pz-z^2\ip g{(I+zA)^{-1}g}\bigr)\\
    &=\det(I+zA)\,(1-pz)
    \left(1-\frac{z^2\ip g{(I+zA)^{-1}g}}{1-pz}\right).
  \end{align*}
  Taking $-\log$ turns each product into a sum of three series,
  \begin{align*}
    -\log\det(I-zT_U)&=-\log\det(I-zA)-\log(1-qz)+L_+(z),\\
    -\log\det(I-zT_W)&=-\log\det(I+zA)-\log(1-pz)+L_-(z),
  \end{align*}
  and we extract the coefficient of $z^m$ from each summand.
  On the left, \eqref{eq:logdet-cycle} gives
  \[
    m\,[z^m]\bigl(-\log\det(I-zT_U)\bigr)=t(C_m,U),
    \qquad
    m\,[z^m]\bigl(-\log\det(I-zT_W)\bigr)=t(C_m,W).
  \]
  For the first summand on the right, \Cref{lem:series-expansions}\,(2) gives, for odd $m$,
  \[
    m\,[z^m]\bigl(-\log\det(I-zA)\bigr)=\Tr(A^m),
    \qquad
    m\,[z^m]\bigl(-\log\det(I+zA)\bigr)=\Tr\bigl((-A)^m\bigr)=-\Tr(A^m),
  \]
  and for the second summand,
  \[
    m\,[z^m]\bigl(-\log(1-qz)\bigr)=q^m,
    \qquad
    m\,[z^m]\bigl(-\log(1-pz)\bigr)=p^m.
  \]
  The last summands equal $m[z^m]L_\pm(z)$, which proves \eqref{eq:cycle-identity-U} and \eqref{eq:cycle-identity-W}.
\end{proof}

\begin{remark} \label{rem:analytic-reading}
  Alternatively, one may interpret these series as functions.
  If a series converges in a neighbourhood of the origin, its sum is a function there, and the coefficients are recovered from that function by differentiation,
  \begin{equation}\label{eq:coeff-def}
    [z^m]F=\frac{F^{(m)}(0)}{m!}.
  \end{equation}
  The arguments below are the same under either interpretations.
\end{remark}

\section{Above the threshold \texorpdfstring{$p\ge2/3$}{p>=2/3}: Dirichlet averages and gamma smoothing}
\label{sec:pge23}

Throughout this section,
\begin{equation}\label{eq:pge23-q}
  q = 1 - p, \qquad 0 \le q \le \frac{1}{3},
\end{equation}
and $m\ge3$ is odd.
We prove \eqref{eq:main-goal} for all such $m$.

\subsection{A two-sided Schur-complement identity and cancellation of the odd
trace}\label{subsec:pge23-two-sided-schur}

\begin{lemma}[Sum of the two cycle densities]\label{lem:pge23-two-sided-identity}
  For every odd $m\ge3$, let
  \begin{equation}\label{eq:pge23-Sm}
    S_m=m[z^m]\bigl(L_-(z)+L_+(z)\bigr).
  \end{equation}
  Then
  \begin{equation}\label{eq:pge23-two-sided-identity}
    t(C_m,W)+t(C_m,U)=p^m+q^m+S_m.
  \end{equation}
  Consequently,
  \begin{equation}\label{eq:pge23-one-sided-identity}
    t(C_m,W)-\bigl(p^m-pq^{m-1}\bigr)
    =q^{m-1}+S_m-t(C_m,U).
  \end{equation}
\end{lemma}

\begin{proof}
  Adding \eqref{eq:cycle-identity-U} and \eqref{eq:cycle-identity-W} cancels the two traces and leaves
  \[
    t(C_m,W)+t(C_m,U)=p^m+q^m+m[z^m]\bigl(L_-(z)+L_+(z)\bigr),
  \]
  which is \eqref{eq:pge23-two-sided-identity} by the definition \eqref{eq:pge23-Sm} of $S_m$.

  To obtain the second identity, rearrange the first:
  \begin{align*}
    t(C_m,W)-\bigl(p^m-pq^{m-1}\bigr)
    &=q^m+S_m-t(C_m,U)+pq^{m-1}\\
    &=(p+q)q^{m-1}+S_m-t(C_m,U)\\
    &=q^{m-1}+S_m-t(C_m,U),
  \end{align*}
  using $p+q=1$.
\end{proof}

\subsection{Replacing a complement cycle by a path}

Let $P_j$ denote the path with $j$ edges and set
\begin{equation}\label{eq:pge23-xj}
  x_j=t(P_j,U)=\ip\one{T_U^j\one}.
\end{equation}
As in \Cref{sec:short-cycles}, deleting one edge from the $m$-cycle leaves a path with $m-1$ edges, so $0\le U\le1$ gives
\begin{equation}\label{eq:pge23-cycle-path}
  t(C_m,U)\le x_{m-1}.
\end{equation}

By \eqref{eq:pge23-one-sided-identity} and \eqref{eq:pge23-cycle-path}, it is enough to prove
\begin{equation}\label{eq:pge23-Phi}
  \Phi_m:=q^{m-1}+S_m-x_{m-1}\ge0.
\end{equation}
Two of the three terms are already expressed through $q$ and $A$, so it remains to do the same for the path densities.
By \Cref{lem:series-expansions}\,(1),
\[
  \sum_{j\ge0}x_jz^j
  =\sum_{j\ge0}z^j\ip\one{T_U^j\one}
  =\ip\one{(I-zT_U)^{-1}\one}.
\]
Write $(I-zT_U)^{-1}\one=a\one+h$ with $h\in\mathcal E_0$, so that $\ip\one{(I-zT_U)^{-1}\one}=a$.
In the block form of $I-zT_U$ this is
\[
  \begin{pmatrix}1-qz&-zg^*\\-zg&I-zA\end{pmatrix}
  \begin{pmatrix}a\\h\end{pmatrix}
  =\begin{pmatrix}1\\0\end{pmatrix}.
\]
Multiply on the left by the triangular matrix of \eqref{eq:block-determinant-identity}, here $\left(\begin{smallmatrix}1&zg^*(I-zA)^{-1}\\0&I\end{smallmatrix}\right)$.
It fixes $\left(\begin{smallmatrix}1\\0\end{smallmatrix}\right)$ and produces the Schur complement in the upper left corner, so the system becomes
\[
  \begin{pmatrix}1-qz-z^2\ip g{(I-zA)^{-1}g}&0\\-zg&I-zA\end{pmatrix}
  \begin{pmatrix}a\\h\end{pmatrix}
  =\begin{pmatrix}1\\0\end{pmatrix}.
\]
Write
\begin{equation}\label{eq:Rplus}
  R_+(z):=\ip g{(I-zA)^{-1}g}=\sum_{j\ge0}z^j\ip g{A^jg} = \int \frac{\dd\mu(\lambda)}{1 - \lambda z},
\end{equation}
where the last equality follows from \eqref{eq:resolvents-common}.
The first row of the system shows
\begin{equation}\label{eq:pge23-path-gf}
  \sum_{j\ge0}x_jz^j=a=\frac{1}{1-qz-z^2R_+(z)}.
\end{equation}

\subsection{Expansion in complete homogeneous polynomials}

In this section, we will give the explicit form of $\Phi_m$ as an integral of a polynomial in $q$ and the eigenvalues of $A$.

For variables $y_1,\ldots,y_k$, the \emph{complete homogeneous symmetric polynomial} of degree $d$ is
\[
  h_d(y_1,\ldots,y_k)
  =\sum_{\substack{\alpha_1,\ldots,\alpha_k\ge0\\
  \alpha_1+\cdots+\alpha_k=d}}
  y_1^{\alpha_1}\cdots y_k^{\alpha_k},
  \qquad h_0=1,
\]
the sum of all monomials of degree $d$ in $y_1,\ldots,y_k$.
Equivalently,
\begin{equation}\label{eq:pge23-h-gf}
  \prod_{i=1}^k\frac{1}{1-y_i z}
  =\sum_{d\ge0}h_d(y_1,\ldots,y_k)z^d.
\end{equation}
The notation $a^{\times r}$ means that $a$ is repeated $r$ times.

\begin{proposition}[Exact coefficient expansion]\label{prop:pge23-expansion}
  For odd $m\ge3$,
  \begin{equation}\label{eq:pge23-expansion}
    \Phi_m=\sum_{r=1}^{(m-1)/2}
    \int\!\cdots\!\int
    P_{m,r}(q;\lambda_1,\ldots,\lambda_r)
    \prod_{i=1}^r\dd\mu(\lambda_i),
  \end{equation}
  where $n=m-2r$ and
  \begin{align}
    P_{m,r}(q;\lambda_1,\ldots,\lambda_r)
    ={}&\frac {m}{r}\Bigl[
    h_n(p^{\times r},-\lambda_1,\ldots,-\lambda_r)
    +h_n(q^{\times r},\lambda_1,\ldots,\lambda_r)
    \Bigr]\notag\\
    &-h_{n-1}(q^{\times(r+1)},\lambda_1,\ldots,\lambda_r).
    \label{eq:pge23-Pmr}
  \end{align}
\end{proposition}

\begin{proof}
  From \eqref{eq:Lplus},
  \begin{align*}
    m[z^m]L_+(z)
    &=m[z^m]\left(-\log\left(1-\frac{z^2R_+(z)}{1-qz}\right)\right)
    \\
    {}&=m[z^m]\sum_{r\ge1}\frac{1}{r}\left(\frac{z^2R_+(z)}{1-qz}\right)^r
    \\
    {}&=\sum_{r\ge1}\frac {m}{r} [z^{m-2r}]\frac{R_+(z)^r}{(1-qz)^r},
  \end{align*}
  where the second equality comes from $-\log(1-x)=\sum_{r\ge1}x^r/r$.
  The spectral representation $R_+(z) = \int \frac{\dd\mu(\lambda)}{1 - \lambda z}$ gives
  \begin{align*}
    R_+(z)^r &=\int\!\cdots\!\int
    \prod_{i=1}^r\frac{1}{1-\lambda_i z}
    \prod_{i=1}^r\dd\mu(\lambda_i),
    \\
    \frac{R_+(z)^r}{(1-qz)^r} &= \int\!\cdots\!\int \frac{1}{(1-qz)^r} \prod_{i=1}^r \frac{1}{1-\lambda_i z} \prod_{i=1}^r \dd\mu(\lambda_i).
  \end{align*}
  By \eqref{eq:pge23-h-gf}, the $r$-th coefficient is
  \[
    [z^{m-2r}] \frac{R_+(z)^r}{(1-qz)^r} = \int\!\cdots\!\int h_{m-2r}(q^{\times r}, \lambda_1, \ldots, \lambda_r) \prod_{i=1}^r \dd\mu(\lambda_i),
  \]
  which is the second term in brackets in \eqref{eq:pge23-Pmr}.
  The same argument for $L_-$ with $R_-(z) = \langle g, (I + zA)^{-1} g \rangle$ produces the first term.

  On the other hand, factoring $1-qz$ out of the denominator of \eqref{eq:pge23-path-gf} and expanding the geometric series gives
  \begin{align*}
    \sum_{j\ge0}x_jz^j
    &=\frac{1}{(1-qz)\left(1-\dfrac{z^2R_+(z)}{1-qz}\right)}
    \\
    {}&=\sum_{r\ge0}\frac{z^{2r}R_+(z)^r}{(1-qz)^{r+1}},
  \end{align*}
  so that
  \[
    x_{m-1}=\sum_{r\ge0}[z^{m-1-2r}]\frac{R_+(z)^r}{(1-qz)^{r+1}}.
  \]
  The term $r=0$ is $[z^{m-1}](1-qz)^{-1}=q^{m-1}$, which cancels the first term of \eqref{eq:pge23-Phi}.
  For $r\ge1$, the spectral representation of $R_+(z)=\int\frac{\dd\mu(\lambda)}{1-\lambda z}$ gives, as above,
  \[
    \frac{R_+(z)^r}{(1-qz)^{r+1}}
    =\int\!\cdots\!\int\frac{1}{(1-qz)^{r+1}}
    \prod_{i=1}^r\frac{1}{1-\lambda_iz}
    \prod_{i=1}^r\dd\mu(\lambda_i),
  \]
  and \eqref{eq:pge23-h-gf}, applied now to the $r+1$ copies of $q$ together with $\lambda_1,\ldots,\lambda_r$, yields
  \[
    [z^{m-1-2r}]\frac{R_+(z)^r}{(1-qz)^{r+1}}
    =\int\!\cdots\!\int
    h_{n-1}(q^{\times(r+1)},\lambda_1,\ldots,\lambda_r)
    \prod_{i=1}^r\dd\mu(\lambda_i),
  \]
  because $m-1-2r=n-1$.
  Subtracting these terms from the contributions of $L_-$ and $L_+$ gives \eqref{eq:pge23-expansion}.
  All three sums are finite and stop at $r=(m-1)/2$, since their coefficients are zero once $2r>m-1$.
\end{proof}

\subsection{Diagonalization by a Dirichlet average}

In this section, we will simplify the integrand $P_{m,r}$ of \eqref{eq:pge23-expansion} by replacing the eigenvalues $\lambda_1,\ldots,\lambda_r$ by a single variable $\ell$.
Define the diagonal polynomial
\begin{equation}\label{eq:pge23-Ptilde}
  \widetilde P_{m,r}(q,\ell)
  =P_{m,r}(q;\ell,\ldots,\ell).
\end{equation}

\begin{lemma}[Dirichlet diagonalization]\label{lem:pge23-dirichlet}
  Let $(\Theta_1,\ldots,\Theta_r)$ be uniformly distributed on the simplex
  \[
    \Theta_i\ge0,\qquad \sum_{i=1}^r\Theta_i=1.
  \]
  Then
  \begin{equation}\label{eq:pge23-dirichlet-identity}
    P_{m,r}(q;\lambda_1,\ldots,\lambda_r)
    =\E\,\widetilde P_{m,r}
    \left(q,\sum_{i=1}^r\Theta_i\lambda_i\right).
  \end{equation}
  Consequently,
  \begin{equation}\label{eq:pge23-diagonal-min}
    \min_{\lambda_i\in[-1/2,1/2]}
    P_{m,r}(q;\lambda_1,\ldots,\lambda_r)
    =\min_{\ell\in[-1/2,1/2]}\widetilde P_{m,r}(q,\ell).
  \end{equation}
\end{lemma}

\begin{proof}
  Write $S=\sum_{i=1}^r\Theta_i\lambda_i$.
  The proof rests on two properties of $h_j$.

  First, the moments of $S$ are the complete homogeneous polynomials, up to a constant:
  \begin{equation}\label{eq:pge23-dirichlet-moment}
    \E\,S^j
    =\frac{h_j(\lambda_1,\ldots,\lambda_r)}
    {\binom{j+r-1}{r-1}}.
  \end{equation}
  To see this, expand
  \[
    S^j=\Bigl(\sum_{i=1}^r\Theta_i\lambda_i\Bigr)^j
    =\sum_{\alpha_1+\cdots+\alpha_r=j}
    \frac{j!}{\alpha_1!\cdots\alpha_r!}
    \prod_{i=1}^r\Theta_i^{\alpha_i}\lambda_i^{\alpha_i},
  \]
  and take expectations.
  From
  \[
    \E\prod_{i=1}^r\Theta_i^{\alpha_i}
    =\frac{(r-1)!\,\alpha_1!\cdots\alpha_r!}{(j+r-1)!},
    \qquad \alpha_1+\cdots+\alpha_r=j,
  \]
  we obtain
  \begin{align*}
    \E\,S^j
    &=\sum_{\alpha_1+\cdots+\alpha_r=j}
    \frac{j!}{\alpha_1!\cdots\alpha_r!}
    \cdot\frac{(r-1)!\,\alpha_1!\cdots\alpha_r!}{(j+r-1)!}
    \prod_{i=1}^r\lambda_i^{\alpha_i}\\
    &=\frac{j!\,(r-1)!}{(j+r-1)!}
    \sum_{\alpha_1+\cdots+\alpha_r=j}\prod_{i=1}^r\lambda_i^{\alpha_i}
    =\frac{h_j(\lambda_1,\ldots,\lambda_r)}{\binom{j+r-1}{r-1}}.
  \end{align*}

  Second, setting all variables equal in the definition of $h_j$ gives the diagonal value
  \begin{equation}\label{eq:pge23-h-diagonal}
    h_j(\ell,\ldots,\ell)=\binom{j+r-1}{r-1}\ell^j,
  \end{equation}
  the binomial coefficient being the number of monomials of degree $j$ in $r$ variables.

  It remains to express the terms of $P_{m,r}$ through the $h_j(\lambda_1,\ldots,\lambda_r)$.
  In the generating function \eqref{eq:pge23-h-gf} of the $h_d$, the $k$ equal variables and the remaining ones give separate factors, so
  \[
    \sum_{d\ge0}h_d(a^{\times k},b_1,\ldots,b_r)z^d
    =\frac{1}{(1-az)^k}\prod_{i=1}^r\frac{1}{1-b_iz}
    =\Bigl(\sum_{s\ge0}\binom{s+k-1}{k-1}a^sz^s\Bigr)
    \Bigl(\sum_{j\ge0}h_j(b_1,\ldots,b_r)z^j\Bigr),
  \]
  and comparing the coefficients of $z^d$ gives the splitting identity
  \begin{equation}\label{eq:pge23-h-split}
    h_d(a^{\times k},b_1,\ldots,b_r)
    =\sum_{j=0}^d
    \binom{d-j+k-1}{k-1}a^{d-j}h_j(b_1,\ldots,b_r).
  \end{equation}
  Apply this to the three terms of $P_{m,r}$, the first two with $k=r$ and $a=p$ or $q$, the third with $k=r+1$ and $a=q$.
  Since $h_j$ is homogeneous of degree $j$, we have $h_j(-\lambda_1,\ldots,-\lambda_r)=(-1)^jh_j(\lambda_1,\ldots,\lambda_r)$, and therefore
  \begin{align}
    P_{m,r}(q;\lambda_1,\ldots,\lambda_r)
    ={}&\frac {m}{r}\sum_{j=0}^{n}\binom{n-j+r-1}{r-1}
    \bigl((-1)^jp^{n-j}+q^{n-j}\bigr)
    h_j(\lambda_1,\ldots,\lambda_r) \notag\\
    &-\sum_{j=0}^{n-1}\binom{n-1-j+r}{r}q^{n-1-j}
    h_j(\lambda_1,\ldots,\lambda_r). \label{eq:pge23-P-in-h}
  \end{align}

  Write the expansion \eqref{eq:pge23-P-in-h} as
  \[
    P_{m,r}(q;\lambda_1,\ldots,\lambda_r)
    =\sum_{j=0}^nc_j\,h_j(\lambda_1,\ldots,\lambda_r),
  \]
  where $c_j$ subsumes the two coefficients of $h_j$ above, which do not involve the $\lambda_i$.
  Setting $\lambda_1=\cdots=\lambda_r=\ell$ and using the diagonal value \eqref{eq:pge23-h-diagonal},
  \begin{equation}\label{eq:pge23-Ptilde-in-c}
    \widetilde P_{m,r}(q,\ell)
    =\sum_{j=0}^nc_j\binom{j+r-1}{r-1}\ell^j.
  \end{equation}
  On the other hand, the moment identity \eqref{eq:pge23-dirichlet-moment} gives
  \[
    P_{m,r}(q;\lambda_1,\ldots,\lambda_r)
    =\sum_{j=0}^nc_j\binom{j+r-1}{r-1}\E\,S^j
    =\E\sum_{j=0}^nc_j\binom{j+r-1}{r-1}S^j
    =\E\,\widetilde P_{m,r}(q,S),
  \]
  the last equality being $\widetilde P_{m,r}$ with $\ell$ replaced by $S$.

  For the second claim, if the $\lambda_i$ lie in $[-1/2,1/2]$ then so does every convex combination $S$, so the right side of the identity \eqref{eq:pge23-dirichlet-identity} is at least $\min_\ell\widetilde P_{m,r}(q,\ell)$; conversely $\lambda_1=\cdots=\lambda_r=\ell$ attains each diagonal value.
\end{proof}

Thus \eqref{eq:pge23-expansion} will be nonnegative once we prove
\begin{equation}\label{eq:pge23-diagonal-target}
  \widetilde P_{m,r}(q,\ell)\ge0
\end{equation}
for
\[
  0\le q\le\frac{1}{3},
  \quad 1\le r\le\frac{m-1}{2},
  \quad -\frac{1}{2}\le\ell\le\frac{1}{2}.
\]
This is the first major simplification: an $r$-variable symmetric polynomial is replaced by one scalar variable polynomial.

\subsection{A beta representation}

In this section, we will further simplify the diagonal polynomial $\widetilde P_{m,r}$ by representing it as an expectation over a beta distribution.

\begin{lemma}[Beta formula]\label{lem:pge23-beta}
  Let $1\le r\le\frac{m-1}{2}$ and put
  \begin{equation}\label{eq:pge23-n}
    n=m-2r,
  \end{equation}
  so that $n\ge1$ is odd.
  Let $X\sim\BetaLaw(r,r)$ and
  \begin{equation}\label{eq:pge23-rho}
    \rho_{n,m}(u)
    =\frac {m}{n}\bigl(u^n+(1-u)^n\bigr)-u^{n-1}.
  \end{equation}
  Then
  \begin{equation}\label{eq:pge23-beta-second}
    \widetilde P_{m,r}(q,\ell)
    =\binom{n+2r-1}{2r-1}\frac {n}{r}
    \E\left[
    X^n\rho_{n,m}\left(q+\ell\frac{1-X}{X}\right)
    \right].
  \end{equation}
  If $\ell>0$, then
  \begin{equation}\label{eq:pge23-beta-integral}
    \widetilde P_{m,r}(q,\ell)
    =\frac{\Gamma(m)}{r\Gamma(n)\Gamma(r)^2}\,
    \ell^{n+r}
    \int_0^\infty
    \frac{s^{r-1}}{(\ell+s)^m}\rho_{n,m}(q+s)\dd s.
  \end{equation}
\end{lemma}

\begin{proof}
  For arbitrary $a,b$ we first claim that
  \begin{equation}\label{eq:pge23-beta-h}
    h_n(a^{\times r},b^{\times r})
    =\binom{n+2r-1}{2r-1}
    \E\bigl(aX+b(1-X)\bigr)^n.
  \end{equation}
  Both sides are polynomials in $a$ and $b$, so it suffices to compare the coefficient of $a^kb^{n-k}$.
  On the left, the splitting identity \eqref{eq:pge23-h-split} and the diagonal value \eqref{eq:pge23-h-diagonal} give
  \[
    h_n(a^{\times r},b^{\times r})
    =\sum_{k=0}^n h_k(a^{\times r})\,h_{n-k}(b^{\times r})
    =\sum_{k=0}^n\binom{k+r-1}{r-1}\binom{n-k+r-1}{r-1}a^kb^{n-k}.
  \]
  On the right, the binomial theorem and the beta moments
  \[
    \E\bigl[X^k(1-X)^{n-k}\bigr]
    =\frac{\Gamma(2r)}{\Gamma(r)^2}
    \cdot\frac{\Gamma(r+k)\,\Gamma(r+n-k)}{\Gamma(2r+n)}
    =\frac{(2r-1)!}{((r-1)!)^2}
    \cdot\frac{(k+r-1)!\,(n-k+r-1)!}{(n+2r-1)!}
  \]
  give
  \[
    \binom{n+2r-1}{2r-1}\E\bigl(aX+b(1-X)\bigr)^n
    =\sum_{k=0}^n\binom{n+2r-1}{2r-1}\binom nk
    \E\bigl[X^k(1-X)^{n-k}\bigr]a^kb^{n-k}.
  \]
  The two coefficients agree, since writing every factor in factorials turns both into $(r+k-1)!\,(r+n-k-1)!/\bigl(k!\,(n-k)!\,((r-1)!)^2\bigr)$; on the right the factors $(2r-1)!$, $n!$, and $(2r+n-1)!$ cancel against those of $\binom{n+2r-1}{2r-1}\binom nk$.

  Next we differentiate both sides of \eqref{eq:pge23-beta-h} with respect to $a$.
  On the left, differentiating the generating function \eqref{eq:pge23-h-gf} in the form $\sum_{d\ge0}h_d(a^{\times r},b^{\times r})z^d=(1-az)^{-r}(1-bz)^{-r}$ gives
  \[
    \sum_{d\ge0}\frac{\partial h_d(a^{\times r},b^{\times r})}{\partial a}z^d
    =rz\,(1-az)^{-(r+1)}(1-bz)^{-r}
    =r\sum_{d\ge0}h_d(a^{\times(r+1)},b^{\times r})z^{d+1},
  \]
  so that $\partial_a h_n(a^{\times r},b^{\times r}) =r\,h_{n-1}(a^{\times(r+1)},b^{\times r})$.
  On the right, $\partial_a\bigl(aX+b(1-X)\bigr)^n =nX\bigl(aX+b(1-X)\bigr)^{n-1}$.
  Hence
  \begin{equation}\label{eq:pge23-beta-h-deriv}
    h_{n-1}(a^{\times(r+1)},b^{\times r})
    =\binom{n+2r-1}{2r-1}\frac {n}{r}
    \E\left[X\bigl(aX+b(1-X)\bigr)^{n-1}\right].
  \end{equation}

  By the definition \eqref{eq:pge23-Ptilde} of the diagonal polynomial,
  \[
    \widetilde P_{m,r}(q,\ell)
    =\frac {m}{r}\bigl[h_n(p^{\times r},(-\ell)^{\times r})
    +h_n(q^{\times r},\ell^{\times r})\bigr]
    -h_{n-1}(q^{\times(r+1)},\ell^{\times r}),
  \]
  and we apply \eqref{eq:pge23-beta-h} to the first two terms, with $(a,b)=(p,-\ell)$ and $(a,b)=(q,\ell)$, and \eqref{eq:pge23-beta-h-deriv} to the third, with $(a,b)=(q,\ell)$.
  Using $\frac {m}{r}=\frac {n}{r}\cdot\frac {m}{n}$, this gives
  \begin{align*}
    &\widetilde P_{m,r}(q,\ell)
    \\
    {} = &\binom{n+2r-1}{2r-1}\frac {n}{r}
    \left[
    \frac {m}{n}\E\Bigl(\bigl(pX-\ell(1-X)\bigr)^n
                       +\bigl(qX+\ell(1-X)\bigr)^n\Bigr)
    -\E\Bigl(X\bigl(qX+\ell(1-X)\bigr)^{n-1}\Bigr)
    \right].
  \end{align*}
  Abbreviating the two linear forms by
  \[
    V_X=qX+\ell(1-X),
    \qquad
    W_X=pX-\ell(1-X),
    \qquad\text{so that } V_X+W_X=X,
  \]
  this simplifies to
  \begin{equation}\label{eq:pge23-beta-first}
    \widetilde P_{m,r}(q,\ell)
    =\binom{n+2r-1}{2r-1}\frac {n}{r}
    \left[
    \frac {m}{n}\E\bigl(V_X^n+W_X^n\bigr)
    -\E\bigl(X V_X^{n-1}\bigr)
    \right].
  \end{equation}

  The three terms $V_X^n$, $W_X^n$, and $XV_X^{n-1}$ all have degree $n$ in $(V_X,W_X,X)$, and $V_X+W_X=X$.
  They therefore depend on $X$ only through the proportion of $X$ taken by $V_X$.
  Writing that proportion as
  \[
    u_X=\frac{V_X}{X}=q+\ell\,\frac{1-X}{X},
    \qquad
    V_X=Xu_X,
    \qquad
    W_X=X(1-u_X),
  \]
  the last equality again because $V_X+W_X=X$, we obtain
  \[
    V_X^n+W_X^n=X^n\bigl(u_X^n+(1-u_X)^n\bigr),
    \qquad
    XV_X^{n-1}=X\cdot X^{n-1}u_X^{n-1}=X^nu_X^{n-1}.
  \]
  By linearity, the two expectations in \eqref{eq:pge23-beta-first} combine into the expectation of
  \[
    \frac {m}{n}\bigl(V_X^n+W_X^n\bigr)-XV_X^{n-1}
    =X^n\left[\frac {m}{n}\bigl(u_X^n+(1-u_X)^n\bigr)-u_X^{n-1}\right]
    =X^n\rho_{n,m}(u_X),
  \]
  by the definition \eqref{eq:pge23-rho} of $\rho_{n,m}$, and this proves \eqref{eq:pge23-beta-second}.

  For $\ell>0$, expand the expectation in \eqref{eq:pge23-beta-second} as an integral against the beta density,
  \[
    \E\left[X^n\rho_{n,m}\left(q+\ell\frac{1-X}{X}\right)\right]
    =\frac{\Gamma(2r)}{\Gamma(r)^2}
    \int_0^1 x^{n+r-1}(1-x)^{r-1}
    \rho_{n,m}\left(q+\ell\frac{1-x}{x}\right)\dd x,
  \]
  and apply the change of variable
  \[
    s=\ell\frac{1-x}{x},
    \qquad x=\frac{\ell}{\ell+s},
    \qquad 1-x=\frac {s}{\ell+s},
    \qquad \dd x=-\frac{\ell}{(\ell+s)^2}\dd s,
  \]
  under which $x=1$ corresponds to $s=0$, and $x\to0^+$ to $s\to\infty$.
  Thus
  \begin{align*}
    &\int_0^1 x^{n+r-1}(1-x)^{r-1}
     \rho_{n,m}\left(q+\ell\frac{1-x}{x}\right)\dd x\\
    &\qquad=\int_\infty^0
     \left(\frac{\ell}{\ell+s}\right)^{n+r-1}
     \left(\frac{s}{\ell+s}\right)^{r-1}
     \rho_{n,m}(q+s)
     \left(-\frac{\ell}{(\ell+s)^2}\right)\dd s\\
    &\qquad=\int_0^\infty
     \frac{\ell^{n+r}s^{r-1}}{(\ell+s)^{n+2r}}\rho_{n,m}(q+s)\dd s
     =\ell^{n+r}\int_0^\infty
     \frac{s^{r-1}}{(\ell+s)^{m}}\rho_{n,m}(q+s)\dd s,
  \end{align*}
  where in the last line the minus sign has been used to interchange the endpoints of integration, and the exponent $n+2r$ has been replaced by $m$ according to \eqref{eq:pge23-n}.
  Altogether
  \[
    \widetilde P_{m,r}(q,\ell)
    =\binom{m-1}{2r-1}\frac {n}{r}\frac{\Gamma(2r)}{\Gamma(r)^2}\,
     \ell^{n+r}\int_0^\infty
     \frac{s^{r-1}}{(\ell+s)^{m}}\rho_{n,m}(q+s)\dd s,
  \]
  and since
  \[
    \binom{m-1}{2r-1}\frac {n}{r}
    \frac{\Gamma(2r)}{\Gamma(r)^2}
    =\frac{(m-1)!}{(2r-1)!\,n!}\cdot\frac {n}{r}
     \cdot\frac{(2r-1)!}{((r-1)!)^2}
    =\frac{\Gamma(m)}{r\Gamma(n)\Gamma(r)^2},
  \]
  this proves \eqref{eq:pge23-beta-integral}.
\end{proof}

\subsection{The case \texorpdfstring{$\ell\le0$}{l<=0}: pointwise nonnegativity}

\begin{lemma}\label{lem:pge23-ell-negative}
  If $0\le q\le1/2$ and $\ell\le0$, then $\widetilde P_{m,r}(q,\ell)\ge0$.
\end{lemma}

\begin{proof}
  From the definition \eqref{eq:pge23-rho} of $\rho_{n,m}$,
  \begin{align*}
    \rho_{n,m}(u)
    &=\frac {m}{n}\bigl(u^n+(1-u)^n\bigr)-u^{n-1} \nonumber \\
    &=\frac {m}{n}(1-u)\bigl((1-u)^{n-1}-u^{n-1}\bigr)
    +\left(\frac {m}{n}-1\right)u^{n-1},
  \end{align*}
  which can be shown by expanding $(1-u)\bigl((1-u)^{n-1}-u^{n-1}\bigr)=(1-u)^n-u^{n-1}+u^n$.
  Recall that $\frac{m}{n}-1=\frac{2r}{n}>0$ by \eqref{eq:pge23-n}, and $n-1$ is even.
  The second term is therefore always nonnegative, and the first is nonnegative whenever $1-u\ge0$ and $|1-u|\ge|u|$, which hold for $u \le 1/2$.

  Since $X\sim\BetaLaw(r,r)$ takes values in $(0,1)$, we have $(1-X)/X\ge0$, so $\ell\le0$ gives
  \[
    q+\ell\,\frac{1-X}{X}\le q\le\frac{1}{2}.
  \]
  Hence $X^n\rho_{n,m}\bigl(q+\ell(1-X)/X\bigr)\ge0$, and \eqref{eq:pge23-beta-second} gives $\widetilde P_{m,r}(q,\ell)\ge0$.
\end{proof}

So only the case $\ell>0$ remains.
There the argument $q+s$ of $\rho_{n,m}$ in \eqref{eq:pge23-beta-integral} ranges over $[q,\infty)$, and unlike the $\ell\le0$ case, $\rho_{n,m}$ can be negative on that range: for instance $\rho_{5,7}(0.7)<0$.
The integrand can therefore no longer be handled pointwise, and nonnegativity has to be derived from the integral as a whole.

\subsection{Laplace--gamma representation}

Since the constant and the factor $\ell^{n+r}$ in \eqref{eq:pge23-beta-integral} are positive, what has to be shown nonnegative for $\ell>0$ is
\[
  \int_0^\infty
  \frac{s^{r-1}}{(\ell+s)^m}\rho_{n,m}(q+s)\dd s.
\]
Here $\ell$ and $s$ are coupled through $(\ell+s)^{-m}$, which is what makes the integral difficult to handle.
They can be separated by an inverse Laplace transform: the inverse Laplace transform of $x\mapsto x^{-m}$ is $t\mapsto t^{m-1}/\Gamma(m)$, so that
\begin{equation}\label{eq:pge23-laplace-kernel}
  \frac{1}{(\ell+s)^m}
  =\frac{1}{\Gamma(m)}\int_0^\infty t^{m-1}e^{-t(\ell+s)}\dd t,
\end{equation}
which can be shown by evaluating the integral under the substitution $w=t(\ell+s)$.

\begin{lemma}[Laplace--gamma representation]\label{lem:pge23-laplace-gamma}
  Let $Y\sim\GammaLaw(r,1)$.
  If $\ell>0$, then
  \begin{equation}\label{eq:pge23-laplace-gamma}
    \widetilde P_{m,r}(q,\ell)
    =\frac{\ell^{n+r}}{r\Gamma(n)\Gamma(r)}
    \int_0^\infty
    t^{m-r-1}e^{-\ell t}
    \E\rho_{n,m}\left(q+\frac {Y}{t}\right)\dd t.
  \end{equation}
\end{lemma}

\begin{proof}
  Insert \eqref{eq:pge23-laplace-kernel} into \eqref{eq:pge23-beta-integral} and exchange the two integrals:
  \begin{align*}
    \widetilde P_{m,r}(q,\ell)
    &=\frac{\Gamma(m)}{r\Gamma(n)\Gamma(r)^2}\,\ell^{n+r}
      \int_0^\infty s^{r-1}\rho_{n,m}(q+s)
      \left(\frac{1}{\Gamma(m)}\int_0^\infty t^{m-1}e^{-t(\ell+s)}\dd t\right)
      \dd s\\
    &=\frac{\ell^{n+r}}{r\Gamma(n)\Gamma(r)^2}
      \int_0^\infty t^{m-1}e^{-\ell t}
      \left(\int_0^\infty s^{r-1}e^{-ts}\rho_{n,m}(q+s)\dd s\right)\dd t.
  \end{align*}
  This exchange is justified by Fubini's theorem, given the absolute integrability of the double integral.
  To verify absolute integrability, integrate the absolute value in $t$ first: by \eqref{eq:pge23-laplace-kernel},
  \[
    \int_0^\infty\!\!\int_0^\infty
    s^{r-1}t^{m-1}e^{-t(\ell+s)}\bigl|\rho_{n,m}(q+s)\bigr|\dd t\dd s
    =\Gamma(m)\int_0^\infty \frac{s^{r-1}}{(\ell+s)^m}
    \bigl|\rho_{n,m}(q+s)\bigr|\dd s.
  \]
  It remains to show that the integral on the right is finite.
  Recall that $\rho_{n,m}$ is a polynomial of degree $n-1$, since the terms $u^n$ and $-u^n$ of $u^n+(1-u)^n$ cancel for odd $n$.
  As $s\to\infty$, the integrand is therefore $O\bigl(s^{r+n-2-m}\bigr)=O\bigl(s^{-r-2}\bigr)$ by \eqref{eq:pge23-n}, which is integrable at infinity.
  As $s\to0$, the integrand is $O\bigl(s^{r-1}\bigr)$, which is integrable because $r\ge1$.

  It remains to compute the inner integral.
  Substituting $y=ts$, so that $s=y/t$ and $\dd s=\dd y/t$, gives
  \[
    \int_0^\infty s^{r-1}e^{-ts}\rho_{n,m}(q+s)\dd s
    =t^{-r}\int_0^\infty y^{r-1}e^{-y}
     \rho_{n,m}\left(q+\frac {y}{t}\right)\dd y
    =\Gamma(r)\,t^{-r}
     \E\rho_{n,m}\left(q+\frac {Y}{t}\right),
  \]
  where the last equality holds because $Y$ has density $y^{r-1}e^{-y}/\Gamma(r)$ on $(0,\infty)$.
  Substituting this for the inner integral obtained after the exchange gives \eqref{eq:pge23-laplace-gamma}.
\end{proof}

In \eqref{eq:pge23-laplace-gamma} the parameter $\ell$ occurs only in the factor $\ell^{n+r}$ and in the weight $t^{m-r-1}e^{-\ell t}$, both of which are positive; the expectation $\E\rho_{n,m}(q+Y/t)$ does not involve $\ell$ at all.
It is therefore enough to prove that this expectation is nonnegative for every $t>0$, and the conclusion then follows for all $\ell>0$ at once.
This is what the next two subsections do.

\subsection{The shifted-gamma moment inequality}

To prove that $\E\rho_{n,m}(q+xY)$ is nonnegative, we center $\rho_{n,m}$ at $u=1/2$ and expand.
This produces the terms with positive constant factors
\[
  2(r+j)v^{2j}-jv^{2j-1},
  \qquad j\ge1,
\]
evaluated at the random $v=q+xY-1/2$.
In each of them, the even power has nonnegative expectation, while the odd power has no definite sign; what is needed is therefore a bound of the odd moment by the even one.
The following lemma supplies such a bound, uniformly in the scaling of $Y$.

\begin{lemma}[Uniform odd-to-even gamma moment bound]\label{lem:pge23-gamma-moment}
  Let $r,j\ge1$ be integers and $Y\sim\GammaLaw(r,1)$.
  For every $z\ge0$,
  \begin{equation}\label{eq:pge23-gamma-moment}
    3j\,\E(zY-1)^{2j-1}
    \le(r+j)\,\E(zY-1)^{2j}.
  \end{equation}
\end{lemma}

We postpone its proof to the next subsection and first show exactly how it proves the nonnegativity of $\E\rho_{n,m}(q+xY)$.

\begin{proposition}[Shifted-gamma positivity]\label{prop:pge23-gamma-positive}
  Let $n\ge1$ be odd, $r\ge1$, and $m=n+2r$.
  If $0\le q\le1/3$, then for every $x\ge0$,
  \begin{equation}\label{eq:pge23-shifted-positive}
    \E\rho_{n,m}(q+xY)\ge0,
    \qquad Y\sim\GammaLaw(r,1).
  \end{equation}
\end{proposition}

\begin{proof}
  Set
  \[
    A_n(u)=u^n+(1-u)^n.
  \]
  We rewrite $\rho_{n,m}$ through $A_n$ and $A_n'$.
  From $A_n'(u)=n\bigl(u^{n-1}-(1-u)^{n-1}\bigr)$,
  \[
    (1-u)A_n'(u)
    =n(1-u)u^{n-1}-n(1-u)^n
    =nu^{n-1}-nu^n-n(1-u)^n
    =nu^{n-1}-nA_n(u),
  \]
  that is, $u^{n-1}=A_n(u)+\frac1n(1-u)A_n'(u)$.
  Substituting this into the definition \eqref{eq:pge23-rho} of $\rho_{n,m}$ and using $m-n=2r$,
  \begin{equation}\label{eq:pge23-rho-derivative}
    n\rho_{n,m}(u)=2rA_n(u)-(1-u)A_n'(u).
  \end{equation}

  Now expand around $u=1/2$.
  With $u=\frac12+v$, the binomial theorem gives
  \begin{equation}\label{eq:pge23-even-expansion}
    A_n\left(\frac{1}{2}+v\right)
    =\left(\frac12+v\right)^n+\left(\frac12-v\right)^n
    =2\sum_{2j\le n}\binom n{2j}\left(\frac12\right)^{n-2j}v^{2j}
    =\sum_{j=0}^{(n-1)/2}a_jv^{2j},
  \end{equation}
  where $a_j=2^{2j+1-n}\binom n{2j}>0$.
  Differentiating in $v$, we also have $A_n'\bigl(\frac12+v\bigr)=\sum_{j\ge1}2ja_jv^{2j-1}$.
  Substituting both into \eqref{eq:pge23-rho-derivative}, with $1-u=\frac12-v$, gives
  \begin{align}
    n\rho_{n,m}\left(\frac{1}{2}+v\right)
    &=2r\sum_{j\ge0}a_jv^{2j}
     -\left(\frac12-v\right)\sum_{j\ge1}2ja_jv^{2j-1} \nonumber \\
    &=2ra_0+
     \sum_{j=1}^{(n-1)/2}
     a_j\bigl(2rv^{2j}-jv^{2j-1}+2jv^{2j}\bigr) \nonumber \\
    &=2ra_0+
     \sum_{j=1}^{(n-1)/2}
     a_j\bigl(2(r+j)v^{2j}-jv^{2j-1}\bigr).\label{eq:pge23-rho-centered}
  \end{align}

  At $u=q+xY$, the $v$ in \eqref{eq:pge23-rho-centered} equals
  \[
    v=q+xY-\frac{1}{2}=xY-\delta=\delta(zY-1),
    \quad \text{where} \quad
    \delta=\frac{1}{2}-q\ge\frac{1}{2}-\frac{1}{3}=\frac{1}{6},
    \qquad
    z=\frac{x}{\delta},
  \]
  which is the form occurring in \Cref{lem:pge23-gamma-moment}.
  Taking expectations in \eqref{eq:pge23-rho-centered}, it suffices to bound each summand.
  For $j\ge1$,
  \begin{align*}
    &\E\left[
    2(r+j)\bigl(\delta(zY-1)\bigr)^{2j}
    -j\bigl(\delta(zY-1)\bigr)^{2j-1}
    \right]\\
    &\qquad=2(r+j)\delta^{2j}\E(zY-1)^{2j}
      -j\delta^{2j-1}\E(zY-1)^{2j-1}\\
    &\qquad\ge2(r+j)\delta^{2j}\E(zY-1)^{2j}
      -\frac{r+j}{3}\delta^{2j-1}\E(zY-1)^{2j}\\
    &\qquad=\delta^{2j}(r+j)
    \left(2-\frac{1}{3\delta}\right)
    \E(zY-1)^{2j}
    \ge0,
  \end{align*}
  where the inequality is \Cref{lem:pge23-gamma-moment} and the last step uses $2\delta\ge1/3$ together with $\E(zY-1)^{2j}\ge0$.
  Since every $a_j$ is positive, we conclude
  \[
    n\,\E\rho_{n,m}(q+xY)\ge2ra_0>0,
  \]
  which proves \eqref{eq:pge23-shifted-positive}.
\end{proof}

\begin{remark}[Tightness of the density requirement $p \ge 2/3$]\label{rem:pge23-threshold}
  This is the only place in the proof where the restriction $p\ge 2/3$ is used.
  As this inequality is tight, this proof only works for $p \ge 2/3$, and we require a different argument for $p<2/3$.
\end{remark}

\subsection{Proof of the gamma moment inequality}\label{subsec:pge23-gamma-moment-proof}

It remains to prove \Cref{lem:pge23-gamma-moment}: for integers $r,j\ge1$, for $Y\sim\GammaLaw(r,1)$, and for every $z\ge0$,
\begin{equation}
  3j\,\E(zY-1)^{2j-1}
  \le(r+j)\,\E(zY-1)^{2j}. \label{eq:pge23-gamma-moment2}
\end{equation}

\begin{proof}[Proof of \Cref{lem:pge23-gamma-moment}]
  The case $z=0$ is immediate, since the left hand side equals $-3j < 0$ and the right hand side equals $r+j > 0$.

  Now assume $z>0$, put $b=1/z$, and define
  \[
    F(b)=\E(Y-b)^{2j},
    \qquad\text{so that}\qquad
    F'(b)=-2j\,\E(Y-b)^{2j-1}.
  \]
  Since $zY-1=z(Y-b)$, dividing \eqref{eq:pge23-gamma-moment2} by $z^{2j}>0$ gives the equivalent inequality
  \begin{equation}\label{eq:pge23-gamma-H}
    H(b):=(r+j)F(b)+\frac{3b}{2}F'(b)\ge0.
  \end{equation}

  For $I_k(b)=\E(Y-b)^k$, apply the Stein identity of the gamma distribution $Y\sim\GammaLaw(r,1)$,
  \begin{equation}\label{eq:pge23-gamma-Stein}
    \E\bigl[Yf'(Y)\bigr]=\E\bigl[(Y-r)f(Y)\bigr],
  \end{equation}
  to $f(y)=(y-b)^k$.
  Its two sides are
  \begin{align*}
    \E[Y f'(Y)] &= k\E\bigl[(Y-b)(Y-b)^{k-1}+b(Y-b)^{k-1}\bigr] = kI_k(b)+kbI_{k-1}(b),\\
    \E[(Y-r)f(Y)] &= \E\bigl[(Y-b)(Y-b)^k + (b-r)(Y-b)^k\bigr] = I_{k+1}(b)-(r-b)I_k(b),
  \end{align*}
  so that
  \begin{equation}\label{eq:pge23-gamma-recurrence}
    I_{k+1}(b)=(r+k-b)I_k(b)+kbI_{k-1}(b).
  \end{equation}
  Taking $k=2j-1$, and using $I_{k+1} = F$, $I_{k}=-\frac{F'}{2j}$ and $I_{k-1}=\frac{F''}{2j(2j-1)}$, gives the differential equation
  \begin{equation}\label{eq:pge23-gamma-ODE}
    bF''(b)+(b-r-2j+1)F'(b)-2jF(b)=0.
  \end{equation}

  We next determine the sign of $H'$ at a zero of $H$.
  Differentiating \eqref{eq:pge23-gamma-H} and then eliminating $F''$ by \eqref{eq:pge23-gamma-ODE}, we obtain
  \[
    H'(b)=\left(r+j+\frac{3}{2}\right)F'(b)+\frac{3b}{2}F''(b)
    =3jF(b)+\left(\frac{5r}{2}+4j-\frac{3b}{2}\right)F'(b).
  \]
  At a zero of $H$ we have $F'(b)=-\frac{2(r+j)}{3b}F(b)$ by \eqref{eq:pge23-gamma-H}, and substituting this leaves
  \begin{equation}\label{eq:pge23-gamma-crossing}
    \frac{H'(b)}{F(b)}
    =\frac{3(r+4j)b-(r+j)(5r+8j)}{3b}.
  \end{equation}
  Set
  \begin{equation}\label{eq:pge23-gamma-bstar}
    b_*:=\frac{(r+j)(5r+8j)}{3(r+4j)}.
  \end{equation}
  Since $F(b)=\E(Y-b)^{2j}>0$, formula \eqref{eq:pge23-gamma-crossing} shows that at a zero $b$ of $H$ we have $H'(b)>0$ if $b>b_*$, and $H'(b)<0$ if $b<b_*$.
  Now
  \[
    H(0)=(r+j)\E Y^{2j}>0,
    \qquad
    H(b)=(r+4j)b^{2j}+O(b^{2j-1})
    \qquad(b\to\infty),
  \]
  so $H$ is positive at both ends of $[0,\infty)$.
  If $H$ were negative somewhere, it would have a first zero, where $H'\le0$, and a later zero, where $H'\ge0$; the first therefore lies below $b_*$ and the second above it, so the interval where $H$ is negative would contain $b_*$.
  It is therefore enough to prove
  \begin{equation}\label{eq:pge23-gamma-Hbstar}
    H(b_*)\ge0.
  \end{equation}

  We first write $H(b_*)$ as an integral.
  Since $F'(b)=-2j\E(Y-b)^{2j-1}$, \eqref{eq:pge23-gamma-H} gives
  \[
    H(b_*)=(r+j)\E(Y-b_*)^{2j}-3jb_*\E(Y-b_*)^{2j-1}
    =(r+j)\,\E\bigl[(Y-b_*)^{2j-1}(Y-b_*-c)\bigr],
  \]
  where $c=\frac{3jb_*}{r+j}$.
  Expanding the expectation as an integral and substituting $x=y-b_*$,
  \begin{equation}\label{eq:pge23-gamma-H-integral}
    H(b_*) =\frac{r+j}{\Gamma(r) e^{b_*}}\int_{-b_*}^{\infty}
    x^{2j-1}(x-c)(x+b_*)^{r-1}e^{-x}\dd x.
  \end{equation}

  Set
  \begin{equation}\label{eq:pge23-gamma-G}
    G(x)=x^{2j}(x+b_*)^re^{-x},
    \qquad -b_*\le x<\infty,
  \end{equation}
  so that
  \[
    G'(x)=x^{2j-1}(x+b_*)^{r-1}e^{-x}
    \bigl(-x^2+(r+2j-b_*)x+2jb_*\bigr).
  \]
  The quadratic factor here factors further as $(c-x)(x+d)$, where
  \[
    d=\frac{2jb_*}{c}=\frac{2(r+j)}{3}.
  \]
  Thus
  \begin{equation}\label{eq:pge23-gamma-Gprime}
    G'(x)=x^{2j-1}(x+b_*)^{r-1}e^{-x}(c-x)(x+d),
  \end{equation}
  and the integrand of \eqref{eq:pge23-gamma-H-integral} is $-G'(x)/(x+d)$, whence
  \[
    \frac{\Gamma(r)e^{b_*}}{r+j}H(b_*)
    =-\int_{-b_*}^{\infty}\frac{G'(x)}{x+d}\dd x.
  \]
  While $-d$ lies inside the range of integration, the quotient $-G'(x)/(x+d)$ is the integrand of \eqref{eq:pge23-gamma-H-integral} and is therefore harmless there.

  We integrate by parts on the two sides of $-d$.
  Away from $x=-d$,
  \[
    \frac{\dd}{\dd x}\left(-\frac{G(x)-G(-d)}{x+d}\right)
    =-\frac{G'(x)}{x+d}+\frac{G(x)-G(-d)}{(x+d)^2}
  \]
  holds, so for $0<\varepsilon<b_*-d$ we have
  \begin{align*}
    -\int_{-b_*}^{-d-\varepsilon}\frac{G'(x)}{x+d}\dd x
    &=\left[-\frac{G(x)-G(-d)}{x+d}\right]_{-b_*}^{-d-\varepsilon}
    -\int_{-b_*}^{-d-\varepsilon}\frac{G(x)-G(-d)}{(x+d)^2}\dd x,\\
    -\int_{-d+\varepsilon}^{\infty}\frac{G'(x)}{x+d}\dd x
    &=\left[-\frac{G(x)-G(-d)}{x+d}\right]_{-d+\varepsilon}^{\infty}
    -\int_{-d+\varepsilon}^{\infty}\frac{G(x)-G(-d)}{(x+d)^2}\dd x.
  \end{align*}
  Here $G(-b_*)=0$ and $G(x)$ converges to zero as $x\to\infty$, so the boundary terms at $-b_*$ and at $\infty$ contribute $G(-d)/(b_*-d)$; the boundary terms at $-d-\varepsilon$ and at $-d+\varepsilon$ vanish as $\varepsilon\downarrow0$, because $G(-d\pm\varepsilon)-G(-d) =O(\varepsilon^2)$ by $G'(-d)=0$; and $(G(x)-G(-d))/(x+d)^2$ extends continuously to $x=-d$ for the same reason, combining the two remaining integrals.
  Hence
  \begin{align}
    -\int_{-b_*}^{\infty}\frac{G'(x)}{x+d}\dd x
    &=\lim_{\varepsilon\downarrow0}
    \left(-\int_{-b_*}^{-d-\varepsilon}\frac{G'(x)}{x+d}\dd x
    -\int_{-d+\varepsilon}^{\infty}\frac{G'(x)}{x+d}\dd x\right)\notag\\
    &=\lim_{\varepsilon\downarrow0}
    \left(\left[-\frac{G(x)-G(-d)}{x+d}\right]_{-b_*}^{-d-\varepsilon}
    +\left[-\frac{G(x)-G(-d)}{x+d}\right]_{-d+\varepsilon}^{\infty}\right.\notag\\
    &\qquad\left.-\int_{-b_*}^{-d-\varepsilon}\frac{G(x)-G(-d)}{(x+d)^2}\dd x
    -\int_{-d+\varepsilon}^{\infty}\frac{G(x)-G(-d)}{(x+d)^2}\dd x\right)\notag\\
    &=\frac{G(-d)}{b_*-d}
    -\int_{-b_*}^{\infty}\frac{G(x)-G(-d)}{(x+d)^2}\dd x.
    \label{eq:pge23-gamma-ibp-max}
  \end{align}
  Here $b_*-d=r(r+j)/(r+4j)>0$ and $G(-d)=d^{2j}(b_*-d)^re^{d}>0$, so the first term is positive.
  For the second term, $(x+d)^{-2}$ is positive as well, so its sign is decided by that of $G(x)-G(-d)$.
  Consequently, if $G(x)-G(-d)\le0$ for every $x\ge-b_*$, that is, if $G(-d)$ is the global maximum of $G$, then the second term is nonnegative and \eqref{eq:pge23-gamma-Hbstar} follows.

  By \eqref{eq:pge23-gamma-Gprime}, the two nonzero local maxima occur at $-d$ and $c$.
  Their ratio is
  \begin{equation}\label{eq:pge23-gamma-max-ratio}
    \frac{G(c)}{G(-d)}
    =\left(\frac {c}{d}\right)^{2j}
    \left(\frac{b_*+c}{b_*-d}\right)^r e^{-(c+d)}.
  \end{equation}
  Let $t=j/r>0$.
  To see that $G$ attains its maximum at $-d$ we prove that the logarithm of \eqref{eq:pge23-gamma-max-ratio}, divided by $r$, namely
  \begin{align}
    L(t)={}&2t\log\left(
    \frac{3t(5+8t)}{2(1+t)(1+4t)}\right)
    +\log\left(
    \frac{(1+4t)(5+8t)}{3(1+t)}\right)\notag\\
    &-\frac{32t^2+25t+2}{3(1+4t)},
    \label{eq:pge23-gamma-L}
  \end{align}
  is negative for every $t>0$.
  Direct differentiation gives, with $A(t)=\frac{3t(5+8t)}{2(1+t)(1+4t)}$,
  \begin{align}
    L'(t)
    &=2\log A(t)
    -\frac{(32t^2+16t-7)(32t^2+25t+2)}
    {3(t+1)(4t+1)^2(8t+5)},\label{eq:pge23-gamma-Lprime}\\
    L''(t)
    &=-\frac{(32t^2+25t+2)
    (128t^4+224t^3+120t^2-28t-25)}
    {t(t+1)^2(4t+1)^3(8t+5)^2}.
    \label{eq:pge23-gamma-Lsecond}
  \end{align}
  Every factor of \eqref{eq:pge23-gamma-Lsecond} other than the quartic $P(t)=128t^4+224t^3+120t^2-28t-25$ is positive for $t>0$, so $L''$ and $P$ have opposite signs.
  The coefficients of $P$ change sign once, so by Descartes' rule of signs $P$ has at most one positive root.
  Since $P(0) = -25 < 0$ and $P(t)\to\infty$, it has exactly one positive root $t_0$, and $L'$ increases on $(0,t_0)$ and decreases on $(t_0,\infty)$.

  We evaluate $L'$ at three places to see its sign changes.
  As $t\downarrow0$ the argument $A(t)$ tends to $0$ while the rational term of \eqref{eq:pge23-gamma-Lprime} tends to $-14/15$, so $L'(t)\to-\infty$.
  At $t=1$ we have $A(1)=39/20$ and the rational term equals $2419/1950$, so the bound
  \[
    \log x\ge\frac{2(x-1)}{x+1}\qquad(x\ge1)
  \]
  gives $2\log A(1)\ge\frac{76}{59}$ and
  \[
    L'(1)\ge\frac{76}{59}-\frac{2419}{1950} \approxeq 0.047622... >0.
  \]
  At $t=3/2$ we have $A(3/2)=153/70$ and the rational term equals $19847/12495$; concavity of the logarithm at $x=2$ gives
  \[
    \log\frac{153}{70}\le\log2+\frac{1}{2}\left(\frac{153}{70}-2\right)
    =\log2+\frac{13}{140},
  \]
  so with $2 \log 2 < 7/5$,
  \[
    L'\left(\frac{3}{2}\right)
    <\frac{7}{5} + \frac{13}{70}-\frac{19847}{12495}\approxeq -0.002681... <0.
  \]

  Since $L'$ increases and then decreases, and is negative near $0$, positive at $t=1$ and negative at $t=3/2$, it has two roots: at some $t_1\in(0,1)$ with $L''(t_1) > 0$, and at some $t_2\in(1,3/2)$ with $L''(t_2) < 0$.
  Thus $L$ decreases on $(0,t_1)$, increases on $(t_1,t_2)$ and decreases on $(t_2,\infty)$, so its supremum over $(0,\infty)$ is the larger of $L(t_2)$ and the limit of $L$ at $0$.
  The latter is negative,
  \[
    \lim_{t\downarrow0}L(t)=\log\frac{5}{3}-\frac{2}{3}<0,
  \]
  and $t_2$ lies in $[1,3/2]$, so it remains to control $L$ on $[1,3/2]$.
  Let
  \[
    B(t)=\frac{(1+4t)(5+8t)}{3(1+t)}
  \]
  so that $L(t)=2t\log A(t)+\log B(t)-(32t^2+25t+2)/(3(4t+1))$.
  By concavity of $\log$,
  \[
    \log A\le\log2+\frac{A-2}{2},
    \qquad
    \log B\le\log13+\frac{B-13}{13}.
  \]
  Using $\log2<7/10$ and $\log13<13/5$, exact simplification gives
  \begin{equation}\label{eq:pge23-gamma-L-rational}
    L(t)<
    \frac{3(96t^3-191t^2-16t+46)}
    {130(t+1)(4t+1)}.
  \end{equation}
  On $[1,3/2]$ we have $96t-191<0$ and $t^2\ge1$, so the cubic numerator is at most $(96t-191)-16t+46=80t-145\le-25$, and $L(t)<0$ there as well.

  To summarize the proof: the signs of $L'$ and $L''$ in \eqref{eq:pge23-gamma-Lprime} and \eqref{eq:pge23-gamma-Lsecond} leave two candidates for the supremum of $L$, namely its limit at $0$, which is negative, and a single local maximum, which lies in $[1,3/2]$.
  Confined to that interval, $L$ is bounded by the rational function \eqref{eq:pge23-gamma-L-rational}, which is negative there; hence $L<0$ on all of $(0,\infty)$.
  The ratio \eqref{eq:pge23-gamma-max-ratio} is therefore smaller than $1$, so $G$ attains its maximum at $-d$ and the integrand $G(x)-G(-d)$ in \eqref{eq:pge23-gamma-ibp-max} is nonpositive.
  The right side of \eqref{eq:pge23-gamma-ibp-max} is then positive, giving $H(b_*)>0$, and the crossing argument yields $H(b)\ge0$ for every $b>0$, which is \eqref{eq:pge23-gamma-H} and hence implies \eqref{eq:pge23-gamma-moment2}.
\end{proof}

\subsection{Completion of the proof for \texorpdfstring{$p\ge2/3$}{p>=2/3}}

\begin{theorem}[The theorem for $p\ge2/3$]\label{thm:pge23-main}
  If $p\ge2/3$, then \eqref{eq:main-goal} holds for every odd $m\ge3$.
\end{theorem}

\begin{proof}
  Equivalently, $q\le1/3$.
  \Cref{lem:pge23-two-sided-identity}, the bound of $t(C_m,U)$ by the path density $x_{m-1}$, and the definition of $\Phi_m$ give
  \begin{align*}
    t(C_m,W)-\bigl(p^m-pq^{m-1}\bigr)
    &=q^{m-1}+S_m-t(C_m,U)\\
    &\ge q^{m-1}+S_m-x_{m-1}=\Phi_m,
  \end{align*}
  so the theorem follows from the nonnegativity of $\Phi_m$.
  By the expansion of \Cref{prop:pge23-expansion}, it is enough to show that $P_{m,r}\ge0$ on $[-1/2,1/2]^r$; and by \Cref{lem:pge23-dirichlet}, which presents $P_{m,r}$ as an expectation of the diagonal values $\widetilde P_{m,r}(q,\ell)$ at points $\ell$ of $[-1/2,1/2]$, it is enough to prove $\widetilde P_{m,r}(q,\ell)\ge0$ for $1\le r\le(m-1)/2$ and $-1/2\le\ell\le1/2$.

  When $\ell\le0$, \Cref{lem:pge23-ell-negative} proves $\widetilde P_{m,r}(q,\ell)\ge0$.
  When $\ell>0$, the beta formula \Cref{lem:pge23-beta} and the Laplace--gamma representation \Cref{lem:pge23-laplace-gamma} write $\widetilde P_{m,r}(q,\ell)$ as a positive weight times an average of the expectations $\E\rho_{n,m}(q+xY)$, each of which is nonnegative by \Cref{prop:pge23-gamma-positive}.
  Hence $\widetilde P_{m,r}(q,\ell)\ge0$ throughout that range, so $\Phi_m\ge0$.
\end{proof}

\section{Below the threshold \texorpdfstring{$1/2<p<2/3$}{1/2<p<2/3}: the single-outlier reduction for \texorpdfstring{$m\ge9$}{m>=9}}
\label{sec:plt23}

From now on, we assume
\begin{equation}\label{eq:plt23-range}
  \frac{1}{2}<p<\frac{2}{3}, \qquad \frac{1}{3}<q<\frac{1}{2},
\end{equation}
and $m\ge9$ is odd.
We continue with the block decomposition \eqref{eq:block-U} and the operator norm bound of \Cref{lem:half-norm}.

\subsection{Nonnegativity of the coefficients of \texorpdfstring{$L_-$}{L-}}\label{subsec:plt23-one-sided-schur}

Recall that
\begin{equation}\label{eq:plt23-one-sided-identity}
  t(C_m,W)=p^m-\Tr(A^m)+m[z^m]L_-(z),
\end{equation}
as proven in \Cref{lem:cycle-identities}.
Throughout \Cref{sec:plt23} we bound these terms directly.

\begin{lemma}[Nonnegative coefficients]\label{lem:plt23-one-sided-identity}
  Every coefficient of $L_-(z)$ is nonnegative.
  Moreover,
  \begin{equation}\label{eq:plt23-Lminus-linear-lower}
    [z^m]L_-(z)\ge\int\frac{p^{m-1}-\lambda^{m-1}}{p+\lambda}\dd\mu(\lambda).
  \end{equation}
\end{lemma}

\begin{proof}
  We first show that the coefficients of
  \[
    Y(z)=\frac{z^2\ip g{(I+zA)^{-1}g}}{1-pz}
  \]
  are nonnegative, which implies the same for $L_-=-\log(1-Y)=\sum_{k\ge1}Y^k/k$ by \eqref{eq:Lminus}.

  Expanding the two factors of $Y$,
  \[
    \ip g{(I+zA)^{-1}g}=\sum_{j\ge0}z^j\int(-\lambda)^j\dd\mu(\lambda),
    \qquad
    \frac{1}{1-pz}=\sum_{k\ge0}p^kz^k,
  \]
  where the first follows from \eqref{eq:resolvents-common} with $A$ replaced by $-A$ and $\mu=\sum_i\norm{P_ig}^2\delta_{\lambda_i}$ is the measure \eqref{eq:mu-def}, gives the coefficients of $Y$ as
  \[
    [z^r]Y(z)
    =\int\sum_{j+k=r-2}p^k(-\lambda)^j\dd\mu(\lambda)
    \qquad(r\ge2).
  \]
  Summing the geometric series inside the integral,
  \begin{equation}\label{eq:plt23-Y-coeff}
    [z^r]Y(z)
    =\int\frac{p^{r-1}-(-\lambda)^{r-1}}{p+\lambda}\dd\mu(\lambda).
  \end{equation}
  By \Cref{lem:half-norm}, the support of $\mu$ lies in $[-1/2,1/2]$, so $|\lambda|<p$ on the whole range of integration and both the numerator and the denominator on the right of \eqref{eq:plt23-Y-coeff} are positive.
  The coefficients of $Y$ are therefore nonnegative.

  For the second claim, the coefficients of $L_-(z)-Y(z)=\sum_{k\ge2}Y(z)^k/k$ are nonnegative as well, so
  \[
    [z^m]L_-(z)\ge[z^m]Y(z)
    =\int\frac{p^{m-1}-(-\lambda)^{m-1}}{p+\lambda}\dd\mu(\lambda)
    =\int\frac{p^{m-1}-\lambda^{m-1}}{p+\lambda}\dd\mu(\lambda),
  \]
  where the first equality follows from \eqref{eq:plt23-Y-coeff} and the second holds because $m-1$ is even.
\end{proof}


\subsection{At most one eigenvalue above \texorpdfstring{$q$}{q}}

We first prove the following bound on the sum of the squared eigenvalues of $A$, in which $\Tr(A^2)=\sum_i\lambda_i^2$.

\begin{lemma}[Hilbert--Schmidt bound]\label{lem:plt23-HS-bound}
  \begin{equation}\label{eq:plt23-HS-bound}
    \Tr(A^2)+2\norm g^2\le pq.
  \end{equation}
\end{lemma}

\begin{proof}
  Since $0\le U\le1$,
  \[
    \Tr(T_U^2)=\int_{[0,1]^2}U(x,y)^2\dd x\dd y
    \le\int_{[0,1]^2}U(x,y)\dd x\dd y=q.
  \]
  Expanding $T_U^2$ in the block decomposition \eqref{eq:block-U},
  \[
    \Tr(T_U^2)=\Tr\begin{pmatrix}q^2 + \norm g^2 & q g^* + g^* A \\ qg + Ag & g g^* + A^2\end{pmatrix}=q^2+2\norm g^2+\Tr(A^2).
  \]
  Subtracting $q^2$ from the both sides proves the claim.
\end{proof}

\begin{lemma}[Eigenvalues above $q$]\label{lem:plt23-eigenvalues-above-q}
  There is at most one eigenvalue of $A$ that is larger than $q$.
  If none of them is larger than $q$, then \eqref{eq:main-goal} holds.
\end{lemma}

\begin{proof}
  If two eigenvalues of $A$ were larger than $q$, then
  \[
    \Tr(A^2)=\sum_i\lambda_i^2>2q^2>pq,
  \]
  the last inequality because $2q>2/3>p$, which contradicts \Cref{lem:plt23-HS-bound}.
  This proves the first claim.

  Assume now that no eigenvalue of $A$ is larger than $q$.
  Then every eigenvalue $\lambda_i$ of $A$ satisfies
  \[
    \lambda_i^m\le q^{m-2}\lambda_i^2,
  \]
  which is immediate for $0\le\lambda_i\le q$, while for $\lambda_i<0$ the left side is negative and the right side is positive.
  Therefore,
  \[
    \Tr(A^m) = \sum_{i}\lambda_i^m\le q^{m-2} \sum_{i} \lambda_i^2 = q^{m-2} \Tr(A^2) \le q^{m-2} \bigl(\Tr(A^2) + 2\norm g^2\bigr) \le pq^{m-1},
  \]
  where we used \Cref{lem:plt23-HS-bound} in the last inequality.
  Since the coefficients of $L_-$ are nonnegative, \eqref{eq:plt23-one-sided-identity} gives
  \[
    t(C_m,W) = p^m - \Tr(A^m) + m[z^m]L_-(z) \ge p^m - pq^{m-1},
  \]
  proving \eqref{eq:main-goal} in this case.
\end{proof}

We henceforth assume that $A$ has exactly one eigenvalue
\begin{equation}\label{eq:plt23-Lambda-above-q}
  \Lambda>q.
\end{equation}
Define
\begin{equation}\label{eq:plt23-M-def}
  M=\sqrt{pq-\Lambda^2},
\end{equation}
which by \eqref{eq:plt23-HS-bound} controls the remaining eigenvalues:
\[
  \sum_{\lambda_i\ne\Lambda}\lambda_i^2\le pq-\Lambda^2=M^2,
\]
so in particular $|\lambda_i|\le M$ for every $\lambda_i\ne\Lambda$.

These quantities are ordered as
\begin{equation}\label{eq:plt23-M-Lambda-order}
  0\le M<q<\Lambda<\frac{1}{2}<p,
\end{equation}
where $M<q$ because $M^2<pq-q^2=q(1-2q)<q^2$ for $q>1/3$, and $\Lambda<1/2$ comes from \Cref{lem:half-norm}.

\subsection{A lower bound via splitting \texorpdfstring{$g$}{g} along the
eigenfunction of \texorpdfstring{$\Lambda$}{Lambda}}\label{subsec:plt23-splitting-lower-bound}

Let $\phi$ be a unit eigenfunction of $A$ for $\Lambda$, which is determined up to sign.
We will choose the sign of $\phi$ later.
Split $g$ into its component along $\phi$ and the rest,
\begin{equation}\label{eq:plt23-g-decomp}
  g=\gamma\phi+g_s,
  \qquad \gamma=\ip g\phi,
  \qquad g_s\perp\phi,
\end{equation}
so that $\norm g^2=\gamma^2+\norm{g_s}^2$.

\begin{proposition}[Lower bound via splitting $g$ along the eigenfunction of $\Lambda$]\label{prop:plt23-splitting-lower-bound}
  For odd $m$, let
  \begin{align}
    k_m(\lambda)&=\frac{p^{m-1}-\lambda^{m-1}}{p+\lambda},
    \label{eq:plt23-km-def}\\
    \Lambda_m&=2M^{m-2}+mk_m(\Lambda),\label{eq:plt23-Lambdam-def}\\
    M_m&=2M^{m-2}+mk_m(M),\label{eq:plt23-Mm-def}\\
    R_m&=\Lambda^m+M^m-pq^{m-1}.\label{eq:plt23-Rm-def}
  \end{align}
  Then
  \begin{align}
    t(C_m,W)-\bigl(p^m-pq^{m-1}\bigr)
    &= pq^{m-1} - \Tr(A^m)+m[z^m]L_-(z) \nonumber
    \\
    {} &\ge -R_m+\Lambda_m\gamma^2+M_m\norm{g_s}^2.\label{eq:plt23-splitting-lower-bound}
  \end{align}
  Moreover,
  \begin{equation}\label{eq:plt23-Lambdam-Mm-positive}
    \Lambda_m>0,
    \qquad
    M_m>0.
  \end{equation}
\end{proposition}

\begin{proof}
  We first bound $m[z^m]L_-(z)$ from below.
  Recall that
  \[
    \mu=\sum_i\norm{P_ig}^2\delta_{\lambda_i}
  \]
  by \eqref{eq:mu-def}, which decomposes $\norm g^2$ along the eigenfunctions of $A$.
  The eigenspace of $\Lambda$ is spanned by $\phi$, so it carries the weight $\ip g\phi^2=\gamma^2$, and every other eigenvalue lies in $[-M,M]$ and carries the weight of $g_s$.
  This allows us to write \eqref{eq:plt23-g-decomp} in terms of $\mu$ as
  \[
    \mu(\{\Lambda\})=\gamma^2,
    \qquad
    \int_{[-M,M]}\dd\mu(\lambda)=\norm{g_s}^2.
  \]
  Then, we have
  \[
    m[z^m]L_-(z)\ge m\int k_m(\lambda)\dd\mu(\lambda)
    =mk_m(\Lambda)\gamma^2
    +m\int_{[-M,M]}k_m(\lambda)\dd\mu(\lambda).
  \]
  where the first equality follows from \eqref{eq:plt23-Lminus-linear-lower} and the second by splitting the integral at $\Lambda$.
  The remaining integral can be bounded from below by $k_m(M)\norm{g_s}^2$, because $k_m$ attains its minimum over $[-M,M]$ at $\lambda=M$.
  To see this, for $0\le s<p$, we first have
  \[
    k_m(-s)=\frac{p^{m-1}-s^{m-1}}{p-s}
    >\frac{p^{m-1}-s^{m-1}}{p+s}=k_m(s),
  \]
  so it suffices to minimize over $0\le\lambda\le M$.
  In this range, $k_m$ is strictly decreasing: with $n=m-1\ge2$,
  \[
    k_m'(\lambda)
    =-\frac{p^n+np\lambda^{n-1}+(n-1)\lambda^n}{(p+\lambda)^2}<0,
  \]
  since every term of the numerator is nonnegative for $\lambda\ge0$ and $p^n>0$.
  Hence
  \begin{equation}\label{eq:plt23-km-min-at-M}
    k_m(\lambda)\ge k_m(M)
    \qquad\text{whenever }|\lambda|\le M,
  \end{equation}
  and altogether
  \begin{equation}\label{eq:plt23-Lminus-lower}
    m[z^m]L_-(z)
    \ge mk_m(\Lambda)\gamma^2+mk_m(M)\norm{g_s}^2.
  \end{equation}

  We next bound $\Tr(A^m)$.
  Splitting off the eigenvalue $\Lambda$,
  \[
    \Tr(A^m)=\Lambda^m+\sum_{\lambda_i\ne\Lambda}\lambda_i^m,
  \]
  where $|\lambda_i|\le M$ gives
  \[
    \sum_{\lambda_i\ne\Lambda}\lambda_i^m
    \le M^{m-2}\sum_{\lambda_i\ne\Lambda}\lambda_i^2.
  \]
  From \eqref{eq:plt23-HS-bound},
  \[
    \Lambda^2+\sum_{\lambda_i\ne\Lambda}\lambda_i^2+2\norm g^2
    =\sum_i\lambda_i^2+2\norm g^2\le pq,
  \]
  so $M^2=pq-\Lambda^2$ gives $\sum_{\lambda_i\ne\Lambda}\lambda_i^2\le M^2-2\norm g^2$.
  Summing up,
  \begin{align}
    -\Tr(A^m)
    &\ge-\Lambda^m-M^{m-2}\sum_{\lambda_i\ne\Lambda}\lambda_i^2\notag\\
    &\ge-\Lambda^m-M^{m-2}\bigl(M^2-2\norm g^2\bigr)\notag\\
    &=-\Lambda^m-M^m+2M^{m-2}\bigl(\gamma^2+\norm{g_s}^2\bigr),
    \label{eq:plt23-trace-upper}
  \end{align}
  where the last equality follows from \eqref{eq:plt23-g-decomp}.

  It remains to combine the two bounds.
  Subtracting $p^m-pq^{m-1}$ in \eqref{eq:plt23-one-sided-identity} gives the first line of \eqref{eq:plt23-splitting-lower-bound},
  \[
    t(C_m,W)-\bigl(p^m-pq^{m-1}\bigr)
    =pq^{m-1}-\Tr(A^m)+m[z^m]L_-(z),
  \]
  and substituting \eqref{eq:plt23-trace-upper} and \eqref{eq:plt23-Lminus-lower},
  \begin{align*}
    &pq^{m-1}-\Tr(A^m)+m[z^m]L_-(z)
    \\
    \ge{}&pq^{m-1}-\Lambda^m-M^m
    +2M^{m-2}\bigl(\gamma^2+\norm{g_s}^2\bigr) +mk_m(\Lambda)\gamma^2+mk_m(M)\norm{g_s}^2\\
    ={}&-\underbrace{(\Lambda^m+M^m-pq^{m-1})}_{\eqqcolon R_m}
    +\underbrace{\bigl(2M^{m-2}+mk_m(\Lambda)\bigr)}_{\eqqcolon \Lambda_m}\gamma^2
    +\underbrace{\bigl(2M^{m-2}+mk_m(M)\bigr)}_{\eqqcolon M_m}\norm{g_s}^2,
  \end{align*}
  which proves \eqref{eq:plt23-splitting-lower-bound}.

  Finally, since $k_m(p)=0$ and $k_m$ is strictly decreasing on $[0,p]$, $k_m$ is positive on $[0,p)$.
  From \eqref{eq:plt23-M-Lambda-order}, we have $M, \Lambda \in [0, p)$, so both $k_m(\Lambda)$ and $k_m(M)$ are positive.
  Since the other terms in the definitions of $\Lambda_m$ and $M_m$ are also positive, \eqref{eq:plt23-Lambdam-Mm-positive} holds.
\end{proof}

\subsection{Minimizing
\texorpdfstring{$\Lambda_m\gamma^2+M_m\norm{g_s}^2$}{Lambda\_m gamma\^{}2 + M\_m |g\_s|\^{}2}}
\label{subsec:plt23-minimization}

For fixed $q$ and $\Lambda$, let $S$ be the set of pairs $(\gamma,\norm{g_s})$ that are realizable by some graphon satisfying the assumptions we made so far.
By deriving inequalities on $\gamma$ and $\norm{g_s}$, we construct a larger set $\overline S\supseteq S$, over which the minimum can be computed, and
\begin{equation}\label{eq:plt23-overapprox}
  \Lambda_m\gamma^2+M_m\norm{g_s}^2
  \ge\min_{(\gamma,\norm{g_s})\in\overline S}
  \bigl(\Lambda_m\gamma^2+M_m\norm{g_s}^2\bigr).
\end{equation}
Specifically, \Cref{lem:plt23-gamma-upper-bound} bounds $\gamma$ from above and \Cref{lem:plt23-gs-lower-bound} bounds $\norm{g_s}$ from below.
Both are linear in $\gamma$ and $\norm{g_s}$, and $\Lambda_m\gamma^2+M_m\norm{g_s}^2$ is a convex quadratic in them by \eqref{eq:plt23-Lambdam-Mm-positive}, so the minimization in \eqref{eq:plt23-overapprox} is convex.

Both lemmas bound $\gamma$ and $\norm{g_s}$ through the decomposition of $|\phi|$ along $\one$ and $\phi$,
\begin{equation}\label{eq:plt23-absphi-decomp}
  |\phi|=a_\phi\one+b\phi+\phi_\perp,
  \qquad
  a_\phi=\int|\phi|,
  \qquad
  b=\ip{|\phi|}\phi,
\end{equation}
where $\phi_\perp$ is orthogonal to $\one$ and $\phi$.
We fix the sign of $\phi$, left open in \Cref{subsec:plt23-splitting-lower-bound}, so that $b\ge0$.
Note that
\begin{equation}\label{eq:plt23-absphi-normalization}
  a_\phi^2+b^2+\norm{\phi_\perp}^2=\norm{|\phi|}^2 = \norm{\phi}^2 = 1.
\end{equation}
Besides this identity, we will need a lower bound on $a_\phi$.
Let $\phi_+$ and $\phi_-$ be the positive and negative parts of $\phi$, so that $\phi=\phi_+-\phi_-$ and $|\phi|=\phi_++\phi_-$.
Since $\phi\perp\one$, we have $\int\phi=0$, and therefore
\begin{equation}\label{eq:plt23-phi-parts}
  \int\phi_+=\int\phi_-=\frac{a_\phi}{2}.
\end{equation}
Expanding $\ip\phi{T_U\phi}$ in these two parts and dropping the cross term, which is nonnegative because $U\ge0$,
\begin{align*}
  \Lambda=\ip\phi{T_U\phi}
  &=\iint U(x,y)\bigl(\phi_+(x)\phi_+(y)+\phi_-(x)\phi_-(y)\bigr)\dd x\dd y\\
  &\qquad-2\iint U(x,y)\phi_+(x)\phi_-(y)\dd x\dd y\\
  &\le\iint U(x,y)
  \bigl(\phi_+(x)\phi_+(y)+\phi_-(x)\phi_-(y)\bigr)\dd x\dd y\\
  &\le\left(\int\phi_+\right)^2+\left(\int\phi_-\right)^2
  =\frac{a_\phi^2}{2},
\end{align*}
the last inequality because $U\le1$ and the last equality by \eqref{eq:plt23-phi-parts}.
Thus
\begin{equation}\label{eq:plt23-aphi-lower}
  a_\phi^2\ge2\Lambda.
\end{equation}

\begin{lemma}[Upper bound for $\gamma$]\label{lem:plt23-gamma-upper-bound}
  \begin{equation}\label{eq:plt23-gamma-upper-bound}
    \gamma\le u_\gamma:=
    \frac{a_\phi^2-2\Lambda-2\Lambda b}{2a_\phi}.
  \end{equation}
\end{lemma}

\begin{proof}
  For every $x$, the following holds
  \[
    (T_U \phi)(x) = \int U(x, y) \phi(y) \dd y \le \int U(x, y) \phi_+(y) \dd y \le \int \phi_+(y) \dd y = \frac{a_\phi}{2},
  \]
  where we used $U \ge 0$ and $\phi \le \phi_+$ in the first inequality, and $U \le 1$ in the second.
  Taking the inner product of both sides with $\phi_+\ge0$ gives
  \begin{equation}\label{eq:plt23-phiplus-pairing}
    \ip{\phi_+}{T_U\phi}
    \le\frac{a_\phi}{2}\int\phi_+
    =\frac{a_\phi^2}{4}.
  \end{equation}
  To expand the left side, the block form \eqref{eq:block-U} gives
  \[
    T_U\phi=\begin{pmatrix}
      q & g^* \\
      g & A
    \end{pmatrix} \begin{pmatrix}
      0 \\ \phi
    \end{pmatrix}=\ip g\phi\,\one+A\phi=\gamma\one+\Lambda\phi,
  \]
  using $\gamma=\ip g\phi$ from \eqref{eq:plt23-g-decomp} and $A\phi=\Lambda\phi$.
  Hence, by \eqref{eq:plt23-phi-parts},
  \[
    \ip{\phi_+}{T_U\phi}
    =\gamma\ip{\phi_+}\one+\Lambda\ip{\phi_+}\phi
    =\frac{\gamma a_\phi}{2}+\Lambda\ip{\phi_+}\phi.
  \]
  Finally, from $\phi_+(x) \phi_-(x) = 0$, we have $\ip{\phi_+}\phi=\ip{\phi_+}{\phi_+} - \ip{\phi_+}{\phi_-} = \norm{\phi_+}^2$, and
  \[
    \norm{\phi_+}^2+\norm{\phi_-}^2=\norm\phi^2=1,
    \qquad
    \norm{\phi_+}^2-\norm{\phi_-}^2=\ip{|\phi|}\phi=b,
  \]
  so $\ip{\phi_+}\phi=(1+b)/2$.
  Substituting into \eqref{eq:plt23-phiplus-pairing} gives \eqref{eq:plt23-gamma-upper-bound}.
\end{proof}

\begin{lemma}[Lower bound for $\norm{g_s}$]\label{lem:plt23-gs-lower-bound}
  The following inequality holds:
  \begin{equation}\label{eq:plt23-gs-inner-lower}
    b\gamma+\ip{g_s}{\phi_\perp}
    \ge\frac{(\Lambda-q)a_\phi^2+(\Lambda-M)\norm{\phi_\perp}^2}{2a_\phi}
    \eqqcolon\kappa\ge0,
  \end{equation}
  where the last inequality holds because $q<\Lambda$ and $M<\Lambda$ by \eqref{eq:plt23-M-Lambda-order}, and $a_\phi=\ip\one{|\phi|}>0$.
  Consequently,
  \begin{equation}\label{eq:plt23-gs-norm-lower}
    \norm{g_s}\norm{\phi_\perp}\ge \ip{g_s}{\phi_\perp}\ge\kappa-b\gamma.
  \end{equation}
\end{lemma}

\begin{proof}
  From $|\phi|=\phi_++\phi_-$ and $\phi=\phi_+-\phi_-$, we have
  \begin{align*}
    \ip{|\phi|}{T_U|\phi|}
    &=\ip{\phi_+}{T_U\phi_+}+\ip{\phi_-}{T_U\phi_-}
    +\ip{\phi_+}{T_U\phi_-}+\ip{\phi_-}{T_U\phi_+},\\
    \ip\phi{T_U\phi}
    &=\ip{\phi_+}{T_U\phi_+}+\ip{\phi_-}{T_U\phi_-}
    -\ip{\phi_+}{T_U\phi_-}-\ip{\phi_-}{T_U\phi_+},
  \end{align*}
  so the two inner products differ only in the sign of their cross terms.
  Since $T_U$ is self-adjoint, $\ip{\phi_-}{T_U\phi_+}=\ip{\phi_+}{T_U\phi_-}$, and therefore
  \begin{align*}
    \ip{|\phi|}{T_U|\phi|}-\ip\phi{T_U\phi}
    &=2\ip{\phi_+}{T_U\phi_-}+2\ip{\phi_-}{T_U\phi_+}
    \\
    &=4\ip{\phi_+}{T_U\phi_-}
    \\
    &= 4 \iint \phi_+(x) U(x, y) \phi_-(y) \dd x \dd y
    \\
    &\ge0,
  \end{align*}
  where the inequality holds because $U\ge0$ and $\phi_\pm\ge0$.
  We also have
  \[
    \ip{\phi}{T_U \phi} = \ip{\phi}{\gamma \one + \Lambda \phi} = \Lambda
  \]
  since $T_U\phi=\gamma\one+\Lambda\phi$ and $\phi\perp\one$.
  Hence
  \begin{equation}\label{eq:plt23-absphi-lower}
    \ip{|\phi|}{T_U|\phi|}\ge\ip\phi{T_U\phi}=\Lambda.
  \end{equation}

  We expand the left side of \eqref{eq:plt23-absphi-lower}.
  In $|\phi|=a_\phi\one+b\phi+\phi_\perp$ the last two terms are orthogonal to $\one$, so the block form \eqref{eq:block-U} gives
  \[
    T_U|\phi|=\begin{pmatrix}
      q & g^* \\
      g & A
    \end{pmatrix} \begin{pmatrix}
      a_\phi \\ b\phi+\phi_\perp
    \end{pmatrix}
    =\bigl(a_\phi g+A(b\phi+\phi_\perp)\bigr)+\bigl(qa_\phi+\ip g{b\phi+\phi_\perp}\bigr)\one
    .
  \]
  Taking the inner product of this with $|\phi|$, and using $\ip{|\phi|}\one=a_\phi$, $\one\perp g,\phi,\phi_\perp$, and that $A$ maps into $\one^\perp$,
  \begin{align*}
    &\ip{|\phi|}{T_U|\phi|}
    \\
    {} = &\ip{a_\phi \one}{a_\phi g + A(b\phi + \phi_\perp)}+ \ip{b\phi+\phi_\perp}{a_\phi g+A(b\phi+\phi_\perp)}+\bigl(qa_\phi+\ip g{b\phi+\phi_\perp}\bigr) \ip{|\phi|}\one
    \\
    {} =&\ip{b\phi+\phi_\perp}{a_\phi g+A(b\phi+\phi_\perp)}+\bigl(qa_\phi+\ip g{b\phi+\phi_\perp}\bigr) \ip{|\phi|}\one
    \\
    {}=&\ip{b\phi+\phi_\perp}{a_\phi g+A(b\phi+\phi_\perp)}+a_\phi\bigl(qa_\phi+\ip g{b\phi+\phi_\perp}\bigr)
    \\
    {} =&qa_\phi^2+2a_\phi\ip g{b\phi+\phi_\perp}
    +\ip{b\phi+\phi_\perp}{A(b\phi+\phi_\perp)}.
  \end{align*}
  The two remaining terms expand by $g=\gamma\phi+g_s$, $A\phi=\Lambda\phi$ and $\phi_\perp\perp\phi$:
  \begin{align*}
    \ip g{b\phi+\phi_\perp}
    &=b\ip g\phi+\ip{\gamma\phi+g_s}{\phi_\perp}
    =b\gamma+\ip{g_s}{\phi_\perp},\\
    \ip{b\phi+\phi_\perp}{A(b\phi+\phi_\perp)}
    &=b^2\ip\phi{A\phi}+2b\Lambda\ip{\phi_\perp}\phi
    +\ip{\phi_\perp}{A\phi_\perp}
    =\Lambda b^2+\ip{\phi_\perp}{A\phi_\perp},
  \end{align*}
  which gives
  \[
    \ip{|\phi|}{T_U |\phi|} = q a_\phi^2 + 2 a_\phi\bigl(b \gamma + \ip{g_s}{\phi_\perp}\bigr) + \Lambda b^2 + \ip{\phi_\perp}{A \phi_\perp}.
  \]
  Every eigenvalue of $A$ other than $\Lambda$ is at most $M$ in absolute value, so $\ip{\phi_\perp}{A\phi_\perp}\le M\norm{\phi_\perp}^2$, and therefore
  \[
    \Lambda\le qa_\phi^2+2a_\phi\bigl(b\gamma+\ip{g_s}{\phi_\perp}\bigr)
    +\Lambda b^2+M\norm{\phi_\perp}^2.
  \]
  By \eqref{eq:plt23-absphi-normalization}, $\Lambda-\Lambda b^2=\Lambda\bigl(a_\phi^2+\norm{\phi_\perp}^2\bigr)$, so rearranging gives
  \[
    (\Lambda-q)a_\phi^2+(\Lambda-M)\norm{\phi_\perp}^2
    \le2a_\phi\bigl(b\gamma+\ip{g_s}{\phi_\perp}\bigr),
  \]
  which is \eqref{eq:plt23-gs-inner-lower}.
  The consequence follows from \eqref{eq:plt23-gs-inner-lower} and Cauchy--Schwarz, $\norm{g_s}\norm{\phi_\perp}\ge\ip{g_s}{\phi_\perp}\ge\kappa-b\gamma$.
\end{proof}

By \eqref{eq:plt23-gamma-upper-bound} and \eqref{eq:plt23-gs-norm-lower}, the pair $(\gamma,\norm{g_s})$ lies in
\begin{equation}\label{eq:plt23-Sbar}
  \overline S=\bigl\{(\overline\gamma,\overline{\norm{g_s}}):
  \overline\gamma\le u_\gamma,\quad
  b\overline\gamma+\overline{\norm{g_s}}\norm{\phi_\perp}\ge\kappa,\quad
  \overline{\norm{g_s}}\ge0\bigr\},
\end{equation}
where $u_\gamma, b, \norm{\phi_\perp}$, and $\kappa$ do not depend on $\gamma$ or $\norm{g_s}$.
Hence
\begin{equation}\label{eq:plt23-min-over-Sbar}
  \Lambda_m\gamma^2+M_m\norm{g_s}^2
  \ge\min_{(\overline\gamma,\overline{\norm{g_s}})\in\overline S}
  \bigl(\Lambda_m\overline\gamma^2+M_m\overline{\norm{g_s}}^2\bigr),
\end{equation}
and the right side is the minimum of a convex quadratic over a polyhedron.

Recall that by \eqref{eq:plt23-splitting-lower-bound} it suffices to prove $\Lambda_m\gamma^2+M_m\norm{g_s}^2\ge R_m$ for \eqref{eq:main-goal}, and that $R_m$ depends only on the eigenvalues $\Lambda$ and $M$, together with $q$ and $m$.
However, the right side of \eqref{eq:plt23-min-over-Sbar} depends on the eigenfunction $\phi$, through the polyhedron $\overline S$, that is through $a_\phi$, $b$, $\norm{\phi_\perp}$ and $u_\gamma$.
To obtain a lower bound comparable with $R_m$, we eliminate all of these eigenfunction quantities in the following theorem.

\begin{theorem}[Reduction to $\psi(\xi,\rho)$]\label{thm:plt23-huber-elim}
  Let
  \begin{align}
    \Omega_m&=\frac{M_m(\Lambda-M)\sqrt{2\Lambda}\,(1-2\Lambda)^2}{4\Lambda^2},
    \label{eq:plt23-C-xi-rho}\\
    \xi&=\frac{4\Lambda^2(\Lambda-q)}{(1-2\Lambda)^2},
    \label{eq:plt23-xi-def}\\
    \rho&=\frac{\Lambda_m}{M_m}\frac{\sqrt\Lambda}{2\sqrt2\,(\Lambda-M)},
    \label{eq:plt23-rho-def}\\
    \psi(\xi,\rho)&=
    \min_{0\le v\le1}
    \left\{\rho v^2+\pos{\xi-v+v^2}\right\}.
    \label{eq:plt23-psi-def}
  \end{align}
  where $\Omega_m$, $\xi$ and $\rho$ are positive by \eqref{eq:plt23-M-Lambda-order} and \eqref{eq:plt23-Lambdam-Mm-positive}.
  Then
  \begin{equation}\label{eq:plt23-huber-payment}
    \Lambda_m\gamma^2+M_m\norm{g_s}^2
    \ge \Omega_m\psi(\xi,\rho).
  \end{equation}
\end{theorem}

\begin{proof}
  We will first remove the dependence on $\overline{\norm{g_s}}$ by minimizing \eqref{eq:plt23-min-over-Sbar} over it.
  Since the objective increases in $\overline{\norm{g_s}}$, one should choose $\overline{\norm{g_s}}$ as small as $\overline S$ allows.
  When $\phi_\perp\ne0$, one should choose $\overline{\norm{g_s}}\norm{\phi_\perp}=\pos{\kappa-b\overline\gamma}$, so that $\overline{\norm{g_s}}^2=\pos{\kappa-b\overline\gamma}^2/\norm{\phi_\perp}^2$ and
  \begin{equation}\label{eq:plt23-one-d-t}
    \min_{(\overline\gamma,\overline{\norm{g_s}})\in\overline S}
    \bigl(\Lambda_m\overline\gamma^2+M_m\overline{\norm{g_s}}^2\bigr)
    =\min_{\overline\gamma\le u_\gamma}
    \left\{
    \Lambda_m\overline\gamma^2
    +\frac{M_m}{\norm{\phi_\perp}^2}\pos{\kappa-b\overline\gamma}^2
    \right\}.
  \end{equation}
  When $\phi_\perp=0$, the variable $\overline{\norm{g_s}}$ does not appear in the constraint, so it should be $0$, and \eqref{eq:plt23-Sbar} becomes a restriction on $\overline\gamma$ alone, $b\overline\gamma\ge\kappa$.
  Then
  \begin{equation}\label{eq:plt23-one-d-t-zero}
    \min_{(\overline\gamma,\overline{\norm{g_s}})\in\overline S}
    \bigl(\Lambda_m\overline\gamma^2+M_m\overline{\norm{g_s}}^2\bigr)
    =\min_{\overline\gamma\le u_\gamma,\ b\overline\gamma\ge\kappa}
    \Lambda_m\overline\gamma^2.
  \end{equation}

  Under the assumption that $\overline\gamma\ge0$, we now provide bounds on $\overline\gamma$ and $b$ that are independent of the eigenfunction.
  First, from \eqref{eq:plt23-gamma-upper-bound} we have
  \[
    2\overline\gamma a_\phi+2\Lambda b
    \le a_\phi^2-2\Lambda=1-2\Lambda-b^2-\norm{\phi_\perp}^2
    \le1-2\Lambda-b^2,
  \]
  and since $a_\phi\ge\sqrt{2\Lambda}$ by \eqref{eq:plt23-aphi-lower}, the assumed hypothesis $\overline\gamma\ge0$ gives $2\overline\gamma a_\phi\ge2\overline\gamma\sqrt{2\Lambda}$, so
  \begin{equation}\label{eq:plt23-b-gamma-budget}
    b^2+2\Lambda b+2\overline\gamma\sqrt{2\Lambda}\le1-2\Lambda.
  \end{equation}

  We can use this inequality to obtain two bounds on $\overline\gamma$ and $b$ separately.
  For $\overline\gamma$, since $\phi$ was chosen so that $b\ge0$, the first two terms on the left of \eqref{eq:plt23-b-gamma-budget} are nonnegative, so
  \begin{equation}\label{eq:plt23-gamma-range}
    0\le\overline\gamma\le\frac{1-2\Lambda}{2\sqrt{2\Lambda}}.
  \end{equation}
  For $b$, viewing \eqref{eq:plt23-b-gamma-budget} as a quadratic inequality in $b$ gives
  \[
    b^2+2\Lambda b-\bigl(1-2\Lambda-2\overline\gamma\sqrt{2\Lambda}\bigr)\le0,
  \]
  whose discriminant is nonnegative by \eqref{eq:plt23-gamma-range}.
  So the larger root of this quadratic, which is $\sqrt{(1-\Lambda)^2-2\overline\gamma\sqrt{2\Lambda}}-\Lambda$, upper bounds $b$:
  \begin{equation}\label{eq:plt23-beta-gamma}
    b\le\beta(\overline\gamma):=
    \sqrt{(1-\Lambda)^2-2\overline\gamma\sqrt{2\Lambda}}-\Lambda.
  \end{equation}
  Therefore, whenever $b$ occurs with a nonnegative coefficient, we can replace it by $\beta(\overline\gamma)$.
  To simplify the optimization further, we now change the optimization variable from $\overline\gamma$ to $\overline v$, which allows us to change the constraint into $\overline v\in[0,1]$: define
  \begin{equation}\label{eq:plt23-v-def}
    \overline v=\frac{2\overline\gamma\sqrt{2\Lambda}}{1-2\Lambda},
    \qquad\text{so that}\qquad
    \overline\gamma=\frac{\overline v(1-2\Lambda)}{2\sqrt{2\Lambda}}
    \quad\text{and}\quad
    \Lambda_m\overline\gamma^2=\Omega_m\rho\overline v^2.
  \end{equation}
  Here $\overline v \in [0,1]$ by \eqref{eq:plt23-gamma-range}, the interval over which the minimum in \eqref{eq:plt23-psi-def} is taken, and the last identity can be shown by expanding the definitions in \eqref{eq:plt23-C-xi-rho} and \eqref{eq:plt23-rho-def}.
  Using the inequality $b\le\beta(\overline\gamma)$ of \eqref{eq:plt23-beta-gamma}, we can also bound $b\overline\gamma$:
  \begin{align}
    b\overline\gamma
    &\le\overline\gamma\beta(\overline\gamma)
    =\frac{\overline v(1-2\Lambda)}{2\sqrt{2\Lambda}}
      \left(\sqrt{(1-\Lambda)^2-2\overline\gamma\sqrt{2\Lambda}}-\Lambda\right)
    \notag\\
    &=\frac{\overline v(1-2\Lambda)}{2\sqrt{2\Lambda}}
      \cdot\frac{1-2\Lambda-2\overline\gamma\sqrt{2\Lambda}}
      {\sqrt{(1-\Lambda)^2-2\overline\gamma\sqrt{2\Lambda}}+\Lambda}
    \notag\\
    &\le\frac{\overline v(1-2\Lambda)}{2\sqrt{2\Lambda}}
      \cdot\frac{1-2\Lambda-2\overline\gamma\sqrt{2\Lambda}}{2\Lambda}
    \notag\\
    &=\frac{(1-2\Lambda)^2\overline v(1-\overline v)}{4\Lambda\sqrt{2\Lambda}},
    \label{eq:plt23-gamma-beta-v}
  \end{align}
  where the first inequality uses $\overline\gamma\ge0$, the second equality follows from rationalizing the numerator, using $(1-\Lambda)^2-\Lambda^2=1-2\Lambda$, and the second inequality holds since the numerator is nonnegative and the denominator is at least $2\Lambda$.

  With these conditional bounds in hand, we now drop the assumption $\overline\gamma\ge0$ and prove \eqref{eq:plt23-huber-payment} by a case analysis on $\phi_\perp$ and $u_\gamma$.

  \textbf{Case 1: $\phi_\perp=0$.}
  The constraint $b\overline\gamma\ge\kappa$ of \eqref{eq:plt23-one-d-t-zero} simplifies to
  \[
    b\overline\gamma\ge\kappa
    =\frac{(\Lambda-q)a_\phi^2+(\Lambda-M)\norm{\phi_\perp}^2}{2a_\phi}
    =\frac{(\Lambda-q)a_\phi}{2}>0,
  \]
  where the last inequality holds since $\Lambda>q$ by \eqref{eq:plt23-M-Lambda-order} and $a_\phi=\ip\one{|\phi|}>0$.
  Since $\phi$ was chosen so that $b\ge0$, this forces $\overline\gamma>0$, so \eqref{eq:plt23-gamma-range}--\eqref{eq:plt23-gamma-beta-v} hold.
  With \eqref{eq:plt23-gamma-beta-v}, we have
  \[
    \frac{(\Lambda-q)\sqrt{2\Lambda}}{2}
    \le\frac{(\Lambda-q)a_\phi}{2}
    \le b\overline\gamma
    \le\frac{(1-2\Lambda)^2\overline v(1-\overline v)}{4\Lambda\sqrt{2\Lambda}},
  \]
  which simplifies to $\xi-\overline v+\overline v^2\le0$.
  Thus the second term in \eqref{eq:plt23-psi-def} is zero at $v=\overline v$, and every pair $(\overline\gamma,\overline{\norm{g_s}})$ in $\overline S$ satisfies
  \[
    \Lambda_m\overline\gamma^2
    =\Omega_m\rho\overline v^2
    =\Omega_m\bigl(\rho\overline v^2+\pos{\xi-\overline v+\overline v^2}\bigr)
    \ge \Omega_m\psi(\xi,\rho),
  \]
  where the first equality is \eqref{eq:plt23-v-def}.
  With \eqref{eq:plt23-one-d-t-zero} this proves \eqref{eq:plt23-huber-payment} in the case $\phi_\perp=0$.

  \textbf{Case 2: $\phi_\perp\ne0$ and $u_\gamma\le0$.}
  Every pair $(\overline\gamma,\overline{\norm{g_s}})$ in $\overline S$ satisfies $\overline\gamma\le u_\gamma\le0$, which implies $\kappa-b\overline\gamma\ge\kappa\ge0$ together with \eqref{eq:plt23-gs-inner-lower} and $b\ge0$, giving the lower bound on the right side of \eqref{eq:plt23-one-d-t}:
  \[
    \min_{\overline\gamma\le u_\gamma}
    \left\{
    \Lambda_m\overline\gamma^2
    +\frac{M_m}{\norm{\phi_\perp}^2}\pos{\kappa-b\overline\gamma}^2
    \right\}
    \ge\frac{M_m}{\norm{\phi_\perp}^2}\kappa^2.
  \]
  Expanding the definition of $\kappa$, we have
  \[
    \frac{\kappa^2}{\norm{\phi_\perp}^2}
    =\frac{\bigl((\Lambda-q)a_\phi^2+(\Lambda-M)\norm{\phi_\perp}^2\bigr)^2}
    {4a_\phi^2\norm{\phi_\perp}^2}\ge(\Lambda-q)(\Lambda-M),
  \]
  where the inequality is equivalent to $\bigl((\Lambda-q)a_\phi^2-(\Lambda-M)\norm{\phi_\perp}^2\bigr)^2\ge0$.
  On the other hand, taking $v=0$ in \eqref{eq:plt23-psi-def} gives
  \[
    \Omega_m\psi(\xi,\rho)
    \le \Omega_m\xi=M_m(\Lambda-M)(\Lambda-q)\sqrt{2\Lambda}
    \le M_m(\Lambda-M)(\Lambda-q),
  \]
  where the last inequality comes from $\Lambda\le1/2$ in \eqref{eq:plt23-M-Lambda-order}.
  Thus \eqref{eq:plt23-huber-payment} holds in this case.

  \textbf{Case 3: $\phi_\perp\ne0$ and $u_\gamma>0$.}
  In the minimization of \eqref{eq:plt23-one-d-t}, any $\overline\gamma<0$ has a larger objective value than at $\overline\gamma=0$, hence we may assume that $\overline\gamma\ge0$, and \eqref{eq:plt23-gamma-range}--\eqref{eq:plt23-gamma-beta-v} hold.
  We now bound the objective on the right of \eqref{eq:plt23-one-d-t} from below at each $\overline\gamma$ with $0\le\overline\gamma\le u_\gamma$.
  Its first term already has the form required in \eqref{eq:plt23-psi-def},
  \[
    \Lambda_m\overline\gamma^2=\Omega_m\rho\overline v^2,
  \]
  by \eqref{eq:plt23-v-def}, so it is enough to consider the second term, which involves $\phi$ through $\kappa$, $b$ and $\norm{\phi_\perp}$.

  Since $2\Lambda\le a_\phi^2\le1$ by \eqref{eq:plt23-aphi-lower} and \eqref{eq:plt23-absphi-normalization},
  \begin{align}
    \kappa
    &=\frac{(\Lambda-q)a_\phi^2+(\Lambda-M)\norm{\phi_\perp}^2}{2a_\phi}
    =\frac{(\Lambda-q)a_\phi}{2}+\frac{(\Lambda-M)\norm{\phi_\perp}^2}{2a_\phi}\notag\\
    &\ge\frac{(\Lambda-q)\sqrt{2\Lambda}}{2}
    +\frac{(\Lambda-M)\norm{\phi_\perp}^2}{2},\label{eq:plt23-H-lower}
  \end{align}
  and in this lower bound for $\kappa$, only the second term involves $\phi_\perp$.
  Define
  \[
    w=\frac{(\Lambda-q)\sqrt{2\Lambda}}{2}-b\overline\gamma, \qquad \text{ so that } \kappa-b\overline\gamma\ge w+ \frac{(\Lambda-M)\norm{\phi_\perp}^2}{2}.
  \]
  where the inequality follows from \eqref{eq:plt23-H-lower}.
  Since $x\mapsto\pos x^2$ is nondecreasing,
  \[
    \frac{M_m}{\norm{\phi_\perp}^2}\pos{\kappa-b\overline\gamma}^2
    \ge M_m\,\frac{\pos{w+(\Lambda-M)\norm{\phi_\perp}^2/2}^2}{\norm{\phi_\perp}^2},
  \]
  and the right side involves $\phi_\perp$ only through $\norm{\phi_\perp}^2$.
  Bounding it below by its infimum over all positive values of $\norm{\phi_\perp}$ therefore removes $\phi_\perp$, and for every $w$ that infimum is
  \begin{equation}\label{eq:plt23-K-AMGM}
    \inf_{\norm{\phi_\perp}>0}
    \frac{\pos{w+(\Lambda-M)\norm{\phi_\perp}^2/2}^2}{\norm{\phi_\perp}^2}
    =2(\Lambda-M)\pos w.
  \end{equation}
  To see this, consider the two cases $w>0$ and $w\le0$ separately.
  When $w>0$, multiplying by $\norm{\phi_\perp}^2$ turns the inequality $\ge$ in \eqref{eq:plt23-K-AMGM} into
  \[
    \pos{w+\frac{(\Lambda-M)\norm{\phi_\perp}^2}{2}}^2
    \ge2(\Lambda-M)\pos w\norm{\phi_\perp}^2,
  \]
  which is the arithmetic--geometric mean inequality, with equality at $\norm{\phi_\perp}^2=2w/(\Lambda-M)$.
  When $w\le0$, the right side of \eqref{eq:plt23-K-AMGM} is zero since $\pos w=0$, and the left side is also zero: since $w+(\Lambda-M)\norm{\phi_\perp}^2/2\le(\Lambda-M)\norm{\phi_\perp}^2/2$ and $x\mapsto\pos x^2$ is nondecreasing,
  \[
    0\le\frac{\pos{w+(\Lambda-M)\norm{\phi_\perp}^2/2}^2}{\norm{\phi_\perp}^2}
    \le\frac{(\Lambda-M)^2\norm{\phi_\perp}^2}{4},
  \]
  and the right side tends to $0$ as $\norm{\phi_\perp}\to0$.
  Hence the second term of the objective is at least $2M_m(\Lambda-M)\pos w$.

  Now the remaining dependence on the eigenfunction is $b$ in $w$, which can be removed with \eqref{eq:plt23-gamma-beta-v}:
  \begin{align}
    \pos w
    &\ge\pos{\frac{(\Lambda-q)\sqrt{2\Lambda}}{2}
    -\frac{(1-2\Lambda)^2\overline v(1-\overline v)}{4\Lambda\sqrt{2\Lambda}}}\notag\\
    &=\pos{\frac{\sqrt{2\Lambda}\,(1-2\Lambda)^2}{8\Lambda^2}
    \left(\frac{4\Lambda^2(\Lambda-q)}{(1-2\Lambda)^2}-\overline v(1-\overline v)\right)}\notag\\
    &=\pos{\frac{\sqrt{2\Lambda}\,(1-2\Lambda)^2}{8\Lambda^2}
    \bigl(\xi-\overline v+\overline v^2\bigr)}\notag\\
    &=\frac{\sqrt{2\Lambda}\,(1-2\Lambda)^2}{8\Lambda^2}
    \pos{\xi-\overline v+\overline v^2},\label{eq:plt23-one-d-gamma}
  \end{align}
  where the second equality is the definition \eqref{eq:plt23-xi-def} of $\xi$.
  Hence the second term of the objective is at least
  \[
    2M_m(\Lambda-M)\frac{\sqrt{2\Lambda}\,(1-2\Lambda)^2}{8\Lambda^2}
    \pos{\xi-\overline v+\overline v^2}
    =\Omega_m\pos{\xi-\overline v+\overline v^2}
  \]
  by the definition \eqref{eq:plt23-C-xi-rho} of $\Omega_m$.
  Both terms of the objective are now bounded from below, so
  \begin{align*}
    \Lambda_m\gamma^2+M_m\norm{g_s}^2
    &\ge\min_{0\le\overline\gamma\le u_\gamma}
    \left\{
    \Lambda_m\overline\gamma^2
    +\frac{M_m}{\norm{\phi_\perp}^2}\pos{\kappa-b\overline\gamma}^2
    \right\}\\
    &\ge\min_{0\le\overline\gamma\le u_\gamma}
    \Omega_m\bigl(\rho\overline v^2+\pos{\xi-\overline v+\overline v^2}\bigr),
  \end{align*}
  where the first inequality combines \eqref{eq:plt23-min-over-Sbar} and \eqref{eq:plt23-one-d-t} with the reduction to $\overline\gamma\ge0$.

  By \eqref{eq:plt23-gamma-range} the range of this minimum lies in $[0,(1-2\Lambda)/(2\sqrt{2\Lambda})]$, which \eqref{eq:plt23-v-def} carries onto $[0,1]$, so
  \begin{align*}
    \min_{0\le\overline\gamma\le u_\gamma}
    \Omega_m\bigl(\rho\overline v^2+\pos{\xi-\overline v+\overline v^2}\bigr)
    &\ge\min_{0\le\overline\gamma\le\frac{1-2\Lambda}{2\sqrt{2\Lambda}}}
    \Omega_m\bigl(\rho\overline v^2+\pos{\xi-\overline v+\overline v^2}\bigr)\\
    &=\Omega_m\min_{0\le v\le1}
    \bigl\{\rho v^2+\pos{\xi-v+v^2}\bigr\}
    =\Omega_m\psi(\xi,\rho)
  \end{align*}
  by \eqref{eq:plt23-psi-def}, which proves \eqref{eq:plt23-huber-payment} in this final case.
\end{proof}

By \Cref{thm:plt23-huber-elim}, it is now enough to show that $\Omega_m\psi(\xi,\rho)\ge R_m$ for every triple $(q,\Lambda,M)$ that a graphon realizes.
Since $M=\sqrt{pq-\Lambda^2}$ is determined by $q$ and $\Lambda$, we now restrict the pairs $(q,\Lambda)$ that can occur, by the following lower bound on $\norm g$.

\begin{lemma}[Variance lower bound]\label{lem:plt23-forced-variance}
  One has
  \begin{equation}\label{eq:plt23-forced-variance}
    \norm g^2\ge
    \frac{\Lambda(\Lambda-q)^2}{2(1-2\Lambda)}.
  \end{equation}
  Consequently,
  \begin{equation}\label{eq:plt23-alpha-ceiling-poly}
    \Lambda^2+q\Lambda-q\le0,
  \end{equation}
  and by solving this inequality, we have
  \begin{equation}\label{eq:plt23-alpha-ceiling}
    q<\Lambda\le r(q):=
    \frac{\sqrt{q^2+4q}-q}{2}.
  \end{equation}
  Moreover $M>0$.
\end{lemma}

\begin{proof}
  Recall from \eqref{eq:plt23-absphi-lower} that $\ip{|\phi|}{T_U|\phi|}\ge\Lambda$.
  In the decomposition $|\phi| = a_\phi \one + b \phi + \phi_\perp$ of \eqref{eq:plt23-absphi-decomp}, the term $b\phi+\phi_\perp$ is orthogonal to $\one$ and has squared norm $b^2+\norm{\phi_\perp}^2=1-a_\phi^2$ by \eqref{eq:plt23-absphi-normalization}.
  Therefore
  \begin{align*}
    \Lambda
    &\le\ip{|\phi|}{T_U|\phi|}\\
    &=qa_\phi^2+2a_\phi\ip g{b\phi+\phi_\perp}
    +\ip{b\phi+\phi_\perp}{A(b\phi+\phi_\perp)}\\
    &\le qa_\phi^2+2a_\phi\norm g\sqrt{1-a_\phi^2}
    +\Lambda(1-a_\phi^2),
  \end{align*}
  where the last inequality bounds $\ip g{b\phi+\phi_\perp}$ by Cauchy--Schwarz and $\ip{b\phi+\phi_\perp}{A(b\phi+\phi_\perp)}$ by the maximality of $\Lambda$.
  Subtracting $\Lambda(1-a_\phi^2)$ and dividing by $a_\phi>0$ gives
  \[
    (\Lambda-q)a_\phi
    \le2\norm g\sqrt{1-a_\phi^2}.
  \]
  Squaring both sides proves \eqref{eq:plt23-forced-variance}:
  \[
    \norm g^2\ge\frac{(\Lambda-q)^2}{4}\cdot\frac{a_\phi^2}{1-a_\phi^2} \ge \frac{\Lambda(\Lambda-q)^2}{2(1 - 2\Lambda)},
  \]
  where we used $a_\phi^2 \ge 2\Lambda$ by \eqref{eq:plt23-aphi-lower} and that $x \mapsto x/(1-x)$ is increasing on $(0,1)$.

  For the second claim, \eqref{eq:plt23-HS-bound} gives $2\norm g^2\le pq-\Lambda^2$, so \eqref{eq:plt23-forced-variance} and $1-2\Lambda>0$ give
  \[
    \Lambda(\Lambda-q)^2\le\bigl(pq-\Lambda^2\bigr)(1-2\Lambda).
  \]
  With $p=1-q$, this is equivalent to $(\Lambda+q-1)(\Lambda^2+q\Lambda-q)\ge0$, and since $\Lambda+q-1<0$ by \eqref{eq:plt23-M-Lambda-order}, this proves \eqref{eq:plt23-alpha-ceiling-poly}.

  Finally, because $\Lambda>q$, the right side of \eqref{eq:plt23-forced-variance} is strictly positive.
  Thus $g\ne0$, and since $\Lambda$ is an eigenvalue of $A$, the bound \eqref{eq:plt23-HS-bound} gives $\Lambda^2\le \Tr(A^2) < \Tr(A^2) + 2\norm{g}^2 \le pq$.
  Consequently $M>0$.
\end{proof}

Together, \Cref{thm:plt23-huber-elim} and \Cref{lem:plt23-forced-variance} show that \eqref{eq:main-goal} follows once
\begin{equation}\label{eq:plt23-scalar-target}
  R_m\le \Omega_m\psi(\xi,\rho)
\end{equation}
is proved for every pair $(q,\Lambda)$ satisfying \eqref{eq:plt23-range} and \eqref{eq:plt23-alpha-ceiling}.

What is needed is a lower bound on $\psi(\xi,\rho)$, but \eqref{eq:plt23-psi-def} defines $\psi$ as a minimum, and bounding a minimum from below requires controlling the point at which it is attained.
The next proposition writes $\psi$ instead as a maximum over a variable $\lambda\in[0,1]$.
Substituting any particular $\lambda$ then gives a lower bound on $\psi$, so the value of $\lambda$ may be chosen to suit each situation.

\begin{proposition}[Dual form of $\psi$]\label{prop:plt23-huber-dual}
  For $\xi\ge0$ and $\rho>0$,
  \begin{equation}\label{eq:plt23-psi-dual}
    \psi(\xi,\rho)=
    \max_{0\le\lambda\le1}
    \left\{
    \lambda\xi-\frac{\lambda^2}{4(\rho+\lambda)}
    \right\}.
  \end{equation}
\end{proposition}

\begin{proof}
  Using the following equivalent form of the positive part,
  \[
    \pos x=\max_{0\le\lambda\le1}\lambda x.
  \]
  the objective in \eqref{eq:plt23-psi-def} can be rewritten as
  \[
    (\rho+\lambda)v^2-\lambda v+\lambda\xi.
  \]
  Since the objective is convex in $v$ and linear in $\lambda$, and the two domains are compact convex intervals, one can apply the minmax theorem:
  \[
    \min_{0\le v\le1}\max_{0\le\lambda\le1} \left\{ (\rho+\lambda)v^2-\lambda v+\lambda\xi \right\} = \max_{0\le\lambda\le1}\min_{0\le v\le1} \left\{ (\rho+\lambda)v^2-\lambda v+\lambda\xi \right\}.
  \]
  The inner optimization over $v$ is attained at
  \[
    v^*=\frac{\lambda}{2(\rho+\lambda)}\in[0,1/2],
  \]
  and substitution gives \eqref{eq:plt23-psi-dual}.
\end{proof}


\subsection{Reduction to a normalized inequality}\label{subsec:plt23-variables}

We retain all notation from \Cref{thm:plt23-huber-elim} and \Cref{prop:plt23-huber-dual}.
This subsection reduces the scalar target \eqref{eq:plt23-scalar-target} to a single inequality between two explicit quantities, stated as \Cref{prop:plt23-reduction} at the end.
The point of the normalization is to keep the two natural scales of the dual expression simultaneously, rather than choosing one value of $\lambda$ and then splitting the parameter region according to which value was chosen.

Define
\begin{equation}\label{eq:plt23-AB-theta}
  \mathcal A:=\frac{\Omega_m\rho\xi^2}{\Lambda^m},
  \qquad
  \mathcal B:=\frac{\Omega_m\xi}{\Lambda^m},
  \qquad
  \theta:=\rho\xi=\frac{\mathcal A}{\mathcal B},
  \qquad
  F_N:=\frac{R_m}{\Lambda^m}.
\end{equation}
The first three are positive.
In these terms the target \eqref{eq:plt23-scalar-target} reads $F_N\le\Omega_m\psi(\xi,\rho)/\Lambda^m$.

We now express $F_N$, $\mathcal A$, and $\mathcal B$ in variables adapted to the problem.
Set
\[
  N:=m-2,
\]
so that $N\ge7$ is odd.
Throughout the rest of this section a tilde denotes division by $\Lambda$:
\[
  \tilde q:=\frac q\Lambda,
  \qquad
  \tilde M:=\frac M\Lambda,
  \qquad
  \tilde p:=\frac p\Lambda,
  \qquad
  u:=1-\tilde q=\frac{\Lambda-q}{\Lambda}.
\]
The relation $M^2=pq-\Lambda^2$ gives
\[
  \tilde p=\frac{1+\tilde M^2}{\tilde q}.
\]
As in the preceding part of the manuscript,
\[
  0<M<q<\Lambda<\frac12<p,
\]
so
\[
  0<\tilde M<\tilde q<1,
  \qquad
  1<\tilde p<2.
\]
We also need the consequence of the frontier ceiling $\Lambda^2+q\Lambda-q\le0$ from \Cref{lem:plt23-forced-variance}.
Writing
\[
  D:=1+\tilde q^2+\tilde M^2,
\]
one has $\Lambda=\tilde q/D$ and $q=\tilde q^2/D$.
Therefore the frontier ceiling is equivalent to
\begin{equation}\label{eq:plt23-frontier-chart}
  \tilde M^2\ge \tilde q(1-\tilde q)=\tilde qu.
\end{equation}
Moreover, $q>1/3$ gives $D<3\tilde q^2$.
Since $D\ge1+\tilde q$ by \eqref{eq:plt23-frontier-chart},
\[
  3\tilde q^2-\tilde q-1>0,
\]
and hence
\[
  \tilde q>\frac{1+\sqrt{13}}6>\frac34.
\]

We now use the two variables that will occur in the bounds below:
\[
  \zeta:=\frac{\tilde M^2}{u}=\frac{M^2}{\Lambda(\Lambda-q)},
  \qquad
  v:=\tilde p-1=\frac{p-\Lambda}{\Lambda}.
\]
Writing $H_\zeta:=\zeta+1+v$, the definitions solve explicitly as
\begin{equation}\label{eq:plt23-zeta-v-inverse}
  u=\frac v{H_\zeta},
  \qquad
  \tilde q=\frac{\zeta+1}{H_\zeta},
  \qquad
  \tilde M^2=\frac{\zeta v}{H_\zeta}.
\end{equation}
Indeed, $\tilde M^2=\zeta u$, while $\tilde p=(1+\tilde M^2)/\tilde q$ and $\tilde q=1-u$ give $1+v=(1+\zeta u)/(1-u)$, which is equivalent to the first identity in \eqref{eq:plt23-zeta-v-inverse}; the other two follow immediately.
The ceiling \eqref{eq:plt23-frontier-chart} becomes
\[
  \zeta(\zeta+v)\ge1.
\]
In particular,
\begin{equation}\label{eq:plt23-zeta-v-consequences}
  0<v<1,
  \qquad
  \zeta\ge \tilde q>\frac34,
  \qquad
  \tilde M^2<v.
\end{equation}
The last inequality follows directly from $\tilde M^2=\zeta v/H_\zeta$ and $H_\zeta>\zeta$.

It is also useful to express $\Lambda$ without $\tilde q$.
Put
\[
  \Delta:=\tilde p^2+1+\tilde M^2.
\]
Since $p+q=1$, $p=\tilde p\Lambda$, and $q=\tilde q\Lambda=(1+\tilde M^2)\Lambda/\tilde p$,
\begin{equation}\label{eq:plt23-Delta-identities}
  \Lambda=\frac{\tilde p}{\Delta},
  \qquad
  q=\frac{1+\tilde M^2}{\Delta},
  \qquad
  1-2\Lambda=\frac{v^2+\tilde M^2}{\Delta}.
\end{equation}
Two elementary consequences will be used repeatedly.
First,
\begin{equation}\label{eq:plt23-Delta-upper-sigma}
  \Delta<\tilde p(v+2),
\end{equation}
for this is equivalent to $1+\tilde M^2<\tilde p$, that is, to $\tilde M^2<v$.
Second, since $q>1/3$,
\begin{equation}\label{eq:plt23-Delta-upper-three}
  \Delta<3(1+\tilde M^2).
\end{equation}
Finally, $q<\Lambda$ and $p=\tilde p\Lambda$ imply
\begin{equation}\label{eq:plt23-Lambda-lower}
  1=p+q<(\tilde p+1)\Lambda=(v+2)\Lambda,
  \qquad\text{so}\qquad
  \Lambda>\frac1{v+2}.
\end{equation}

Since $m=N+2$, $M=\tilde M\Lambda$, $q=\tilde q\Lambda$, and $p=\tilde p\Lambda=(1+\tilde M^2)\Lambda/\tilde q$, the definition $R_m=\Lambda^m+M^m-pq^{m-1}$ gives
\begin{equation}\label{eq:plt23-F-formula}
  F_N
  =1-\tilde q^N-\tilde M^2(\tilde q^N-\tilde M^N).
\end{equation}
Put
\[
  T_N:=N-\zeta(\tilde q^N-\tilde M^N).
\]
Because $1-\tilde q^N\le N(1-\tilde q)=Nu$ and $\tilde M^2=\zeta u$,
\begin{equation}\label{eq:plt23-F-T}
  F_N\le uT_N.
\end{equation}
Thus $F_N\le0$ already proves \eqref{eq:plt23-scalar-target}, because $\psi(\xi,\rho)\ge0$.
Whenever convenient below we may therefore assume
\begin{equation}\label{eq:plt23-F-positive}
  F_N>0,
\end{equation}
which by \eqref{eq:plt23-F-T} implies $T_N>0$.

Next normalize the two coefficients $\Lambda_m$ and $M_m$ of \Cref{prop:plt23-splitting-lower-bound}:
\[
  \widehat{\Lambda}_N:=\frac{\Lambda_m}{\Lambda^N},
  \qquad
  \widehat M_N:=\frac{M_m}{\Lambda^N}.
\]
A direct substitution into the definitions of $k_m$, $\Lambda_m$, and $M_m$ gives
\begin{align}
  \widehat{\Lambda}_N
  &=2\tilde M^N+(N+2)K_\Lambda,
  &
  K_\Lambda
  &:=\frac{\tilde p^{N+1}-1}{\tilde p+1},
  \label{eq:plt23-KA-def}\\
  \widehat M_N
  &=2\tilde M^N+(N+2)K_M,
  &
  K_M
  &:=\frac{\tilde p^{N+1}-\tilde M^{N+1}}{\tilde p+\tilde M}.
  \label{eq:plt23-KL-def}
\end{align}
Since $N+1$ is even,
\begin{align}
  K_\Lambda
  &=v\sum_{j=0}^{(N-1)/2}\tilde p^{2j},
  \label{eq:plt23-KA-sum}\\
  K_M
  &=(\tilde p-\tilde M)
    \sum_{j=0}^{(N-1)/2}
      \tilde p^{N-1-2j}\tilde M^{2j}.
  \label{eq:plt23-KL-sum}
\end{align}
For brevity set
\[
  S_N:=\sum_{j=0}^{(N-1)/2}\tilde p^{2j}.
\]
There are $(N+1)/2$ summands and their geometric mean is $\tilde p^{(N-1)/2}$.
Hence the arithmetic--geometric mean inequality gives
\begin{equation}\label{eq:plt23-SN-AMGM}
  S_N\ge \frac{N+1}{2}\tilde p^{(N-1)/2}.
\end{equation}

The definitions \eqref{eq:plt23-C-xi-rho}--\eqref{eq:plt23-rho-def} of $\Omega_m$, $\xi$, and $\rho$ now give closed forms for $\mathcal A$ and $\mathcal B$.
Multiplying $\Omega_m$ by $\xi$ cancels the factor $(1-2\Lambda)^2/(4\Lambda^2)$, and multiplying $\Omega_m$ by $\rho$ cancels $M_m(\Lambda-M)$, so that
\[
  \Omega_m\xi=M_m(\Lambda-M)\sqrt{2\Lambda}\,(\Lambda-q),
  \qquad
  \Omega_m\rho=\frac{\Lambda_m(1-2\Lambda)^2}{8\Lambda},
\]
and hence, using $\xi^2=16\Lambda^4(\Lambda-q)^2/(1-2\Lambda)^4$,
\[
  \Omega_m\rho\xi^2=\frac{2\Lambda^3\Lambda_m(\Lambda-q)^2}{(1-2\Lambda)^2}.
\]
Substituting $\Lambda-q=\Lambda u$, $\Lambda-M=\Lambda(1-\tilde M)$, $\Lambda_m=\Lambda^N\widehat{\Lambda}_N$, $M_m=\Lambda^N\widehat M_N$, and $m=N+2$ turns these into
\begin{equation}\label{eq:plt23-AB-formulas}
  \mathcal A
  =\frac{2\Lambda^3u^2}{(1-2\Lambda)^2}\,\widehat{\Lambda}_N,
  \qquad
  \mathcal B
  =u(1-\tilde M)\sqrt{2\Lambda}\,\widehat M_N.
\end{equation}
Each of $F_N$, $\mathcal A$, and $\mathcal B$ carries a factor $u$, so we also set
\[
  G:=\frac{F_N}{u},
  \qquad
  X:=\frac{\mathcal A}{u},
  \qquad
  Y:=\frac{\mathcal B}{u},
\]
and note that
\begin{equation}\label{eq:plt23-GXY-ratios}
  \frac{F_N}{\mathcal A}=\frac GX,
  \qquad
  \frac{F_N}{\mathcal B}=\frac GY.
\end{equation}

The reduction can now be stated.

\begin{proposition}[Reduction to a normalized inequality]\label{prop:plt23-reduction}
  Let
  \begin{equation}\label{eq:plt23-dtheta}
    d(\theta):=
    \begin{cases}
      \theta, & 0<\theta\le\frac12,\\[1mm]
      1-\dfrac{1}{4\theta}, & \theta\ge\frac12.
    \end{cases}
  \end{equation}
  If $F_N\le\mathcal B\,d(\theta)$, then the scalar target \eqref{eq:plt23-scalar-target} holds.
\end{proposition}

\begin{proof}
  Dividing the objective in \eqref{eq:plt23-psi-dual} by $\xi$ gives
  \begin{align*}
    \frac{\Omega_m\psi(\xi,\rho)}{\Lambda^m}
    &=\mathcal B
      \max_{0\le\lambda\le 1}
        \left\{\lambda-\frac{\lambda^2}{4(\theta+\lambda\xi)}\right\}\\
    &\ge
      \mathcal B
      \max_{0\le\lambda\le 1}
        \left\{\lambda-\frac{\lambda^2}{4\theta}\right\},
  \end{align*}
  because $\theta+\lambda\xi\ge\theta$, so replacing the denominator by $\theta$ makes the subtracted term larger.
  In the last maximum the unconstrained maximizer is $\lambda=2\theta$, so the maximum is attained at $\lambda=2\theta$ when $2\theta\le1$ and at $\lambda=1$ when $2\theta\ge1$, and its value is $d(\theta)$ in both cases.
  Hence
  \[
    \frac{\Omega_m\psi(\xi,\rho)}{\Lambda^m}
    \ge\mathcal B\,d(\theta)
    \ge F_N,
  \]
  and multiplying by $\Lambda^m$ gives \eqref{eq:plt23-scalar-target}.
\end{proof}

\subsection{Two elementary bounds for the defect}

We only need two of the bounds used in the former proof.

\begin{lemma}[Defect bounds]\label{lem:plt23-T-bounds}
  For every admissible point and every odd $N\ge7$,
  \begin{align}
    T_N
    &\le N(1+v)-\zeta\bigl(1-v^{N/2}\bigr),
    \label{eq:plt23-T-bound-one}\\
    T_N
    &<N(1-\tilde M^{N+1}).
    \label{eq:plt23-T-bound-two}
  \end{align}
\end{lemma}

\begin{proof}
  For the first inequality, Bernoulli gives $\tilde q^N=(1-u)^N\ge1-Nu$.
  Also $\tilde M^2<v$ by \eqref{eq:plt23-zeta-v-consequences}, hence $\tilde M^N<v^{N/2}$ and $\zeta u=\tilde M^2<v$.
  Therefore
  \begin{align*}
    T_N
    &=N-\zeta(\tilde q^N-\tilde M^N)\\
    &\le N-\zeta+N\zeta u+\zeta\tilde M^N\\
    &\le N(1+v)-\zeta\bigl(1-v^{N/2}\bigr).
  \end{align*}

  For the second inequality, $D<3\tilde q^2$ gives
  \[
    \tilde M^2<2\tilde q^2-1<(2\tilde q-1)^2,
  \]
  so $\tilde M<2\tilde q-1$, equivalently
  \[
    \tilde q-\tilde M>1-\tilde q=u.
  \]
  Convexity of $x\mapsto x^N$ on $[0,\infty)$ gives
  \[
    \tilde q^N-\tilde M^N
    \ge N\tilde M^{N-1}(\tilde q-\tilde M)
    >Nu\tilde M^{N-1}.
  \]
  Multiplying by $\zeta=\tilde M^2/u$ yields
  \[
    \zeta(\tilde q^N-\tilde M^N)>N\tilde M^{N+1},
  \]
  which is exactly \eqref{eq:plt23-T-bound-two}.
\end{proof}

\subsection{An elementary product bound}

The next lemma replaces all of the polynomial positivity checks that appeared in the former small-$v$ analysis.

\begin{lemma}[Uniform product bound]\label{lem:plt23-product}
  Let $N\ge7$ be odd.
  For every $0\le y\le1/\sqrt2$,
  \begin{equation}\label{eq:plt23-product}
    (1-y)(1-y+y^2)(1+y^2)^{N-2}
    \ge
    \frac{N}{N+2}\sqrt{1+\frac{y^2}{2}}.
  \end{equation}
\end{lemma}

\begin{proof}
  We first prove the elementary logarithmic bound
  \begin{equation}\label{eq:plt23-log-product}
    \log\bigl((1-y)(1-y+y^2)\bigr)
    \ge -2y-\frac{y^2}{8}
    \qquad
    \left(0\le y\le\frac1{\sqrt2}\right).
  \end{equation}
  Set
  \[
    h(y)
    :=\log\bigl((1-y)(1-y+y^2)\bigr)+2y+\frac{y^2}{8}.
  \]
  A direct differentiation gives
  \[
    h'(y)
    =\frac{y\,\mathcal Q(y)}
           {4(y-1)(y^2-y+1)},
    \qquad
    \mathcal Q(y):=y^3+6y^2-2y-1.
  \]
  On $(0,1)$ the denominator is negative.
  The cubic $\mathcal Q$ has exactly one positive zero: its derivative $3y^2+12y-2$ has exactly one positive zero, so $\mathcal Q$ first decreases and then increases, while $\mathcal Q(0)=-1$ and $\mathcal Q(1/\sqrt2)>0$.
  Consequently $h$ first increases and then decreases on $[0,1/\sqrt2]$, and its minimum occurs at an endpoint.
  Clearly $h(0)=0$.
  At the other endpoint, write $s=1/\sqrt2$, so that
  \[
    (1-s)(1-s+s^2)=2-\frac{5\sqrt2}{4}.
  \]
  To verify $h(s)>0$, it is enough to show
  \[
    \left(2-\frac{5\sqrt2}{4}\right)
    e^{\sqrt2+1/16}>1.
  \]
  Since $e^x\ge\sum_{j=0}^5x^j/j!$ for $x\ge0$, direct expansion gives
  \begin{align*}
   &\left(2-\frac{5\sqrt2}{4}\right)
    \sum_{j=0}^5\frac{(\sqrt2+1/16)^j}{j!}-1\\
   &\hspace{2cm}
    =\frac{-9799832+11426027\sqrt2}{503316480}>0.
  \end{align*}
  Thus \eqref{eq:plt23-log-product} holds.

  Write $\mathcal H_N$ for the quotient of the two sides of \eqref{eq:plt23-product},
  \[
    \mathcal H_N(y)
    :=\frac{N+2}{N}
      \frac{(1-y)(1-y+y^2)(1+y^2)^{N-2}}
           {\sqrt{1+y^2/2}},
  \]
  so that \eqref{eq:plt23-product} is the statement $\log\mathcal H_N(y)\ge0$.
  Using $\log(1+t)\ge t-t^2/2$ and $\log(1+t)\le t$ for $t\ge0$, together with \eqref{eq:plt23-log-product},
  \begin{equation}\label{eq:plt23-Psi-lower}
    \log\mathcal H_N(y)
    \ge
    \Psi_N(y):=
    \frac{2(N-1)}{N^2}
    -2y
    +\left(N-\frac{19}{8}\right)y^2
    -\frac{N-2}{2}y^4.
  \end{equation}

  First suppose $0\le y\le1/4$.
  Then $y^4\le y^2/16$, so
  \[
    \Psi_N(y)
    \ge
    \frac{2(N-1)}{N^2}-2y+\alpha_Ny^2,
    \qquad
    \alpha_N:=\frac{31N-74}{32}.
  \]
  The quadratic on the right has global minimum $2(N-1)/N^2-1/\alpha_N$.
  This is nonnegative because
  \begin{align*}
    \frac{2(N-1)}{N^2}\ge\frac1{\alpha_N}
    &\iff
    2(N-1)(31N-74)\ge32N^2\\
    &\iff
    15N^2-105N+74\ge0,
  \end{align*}
  and the last expression is positive at $N=7$ and increasing for $N\ge7$.

  It remains to consider $1/4\le y\le1/\sqrt2$.
  For odd $N\ge7$,
  \[
    \frac{\mathcal H_{N+2}(y)}{\mathcal H_N(y)}
    =\frac{N(N+4)}{(N+2)^2}(1+y^2)^2.
  \]
  When $y\ge1/4$, the right side is at least
  \[
    \frac{N(N+4)}{(N+2)^2}\left(\frac{17}{16}\right)^2>1,
  \]
  because
  \[
    289N(N+4)-256(N+2)^2
    =33N^2+132N-1024>0
    \qquad (N\ge7).
  \]
  Thus it is enough to prove the claim for $N=7$, where \eqref{eq:plt23-Psi-lower} reads
  \[
    \Psi_7(y)=\frac{12}{49}-2y+\frac{37}{8}y^2-\frac52y^4.
  \]
  Now
  \[
    \Psi_7'(y)=-2+\frac{37}{4}y-10y^3,
    \qquad
    \Psi_7''(y)=\frac{37}{4}-30y^2.
  \]
  Hence $\Psi_7'$ first increases and then decreases on our interval, so its minimum there occurs at an endpoint.
  We have
  \[
    \Psi_7'\left(\frac14\right)=\frac5{32}>0,
    \qquad
    \Psi_7'\left(\frac1{\sqrt2}\right)
    =-2+\frac{17}{4\sqrt2}>0.
  \]
  Therefore $\Psi_7$ is increasing on $[1/4,1/\sqrt2]$, and
  \[
    \Psi_7(y)\ge\Psi_7\left(\frac14\right)=\frac{607}{25088}>0.
  \]
  This completes the proof.
\end{proof}

\subsection{The range \texorpdfstring{$0<v\le1/2$}{0<v<=1/2}}

We now prove two simultaneous bounds.
The first controls the quadratic scale $\mathcal A$, while the second keeps both $\mathcal A$ and $\mathcal B$ at once.

\begin{proposition}[The low-$v$ bounds]\label{prop:plt23-low-v}
  Suppose $0<v\le1/2$.
  Then
  \begin{align}
    2F_N&\le\mathcal A,
    \label{eq:plt23-low-v-A}\\
    \frac{F_N}{\mathcal A}+2\frac{F_N}{\mathcal B}&\le2.
    \label{eq:plt23-low-v-both}
  \end{align}
\end{proposition}

\begin{proof}
  If $F_N\le0$, both inequalities are immediate.
  Assume \eqref{eq:plt23-F-positive}.
  Put
  \[
    z:=\frac\zeta{\tilde p},
    \qquad
    c:=1-v^{N/2}.
  \]
  Since $\zeta>3/4$ and $\tilde p\le3/2$,
  \[
    z>\frac12.
  \]
  Because $\zeta=\tilde pz$ and $1+v=\tilde p$, \eqref{eq:plt23-zeta-v-inverse} becomes
  \begin{equation}\label{eq:plt23-low-v-u-Mtilde}
    u=\frac{v}{\tilde p(1+z)},
    \qquad
    \tilde M^2=\frac{vz}{1+z}.
  \end{equation}
  By \eqref{eq:plt23-F-T} and \eqref{eq:plt23-T-bound-one},
  \begin{equation}\label{eq:plt23-G-low-v}
    G\le T_N\le\tilde pP,
    \qquad
    P:=N-cz.
  \end{equation}
  Since $F_N>0$, the left side is positive, so $P>0$.

  We first bound $X$ from below.
  From \eqref{eq:plt23-AB-formulas}, \eqref{eq:plt23-Delta-identities}, and \eqref{eq:plt23-low-v-u-Mtilde},
  \begin{align*}
    X
    &=\frac{2\Lambda^3u}{(1-2\Lambda)^2}\,\widehat{\Lambda}_N\\
    &=\frac{2\tilde p^2(1+z)}
            {\Delta\,v\,(v+\tilde pz)^2}\,\widehat{\Lambda}_N.
  \end{align*}
  Indeed,
  \[
    v^2+\tilde M^2
    =v\left(v+\frac{z}{1+z}\right)
    =\frac{v(v+\tilde pz)}{1+z}.
  \]
  By \eqref{eq:plt23-KA-def} and \eqref{eq:plt23-KA-sum}, $\widehat{\Lambda}_N\ge(N+2)vS_N$.
  Using also \eqref{eq:plt23-Delta-upper-sigma},
  \begin{equation}\label{eq:plt23-X-low-v-1}
    X
    >\frac{2\tilde p(N+2)S_N(1+z)}
            {(v+2)(v+\tilde pz)^2}.
  \end{equation}
  We claim
  \[
    (v+2)(v+\tilde pz)^2
    \le2\tilde p^3z(1+z).
  \]
  The difference between the right and left sides is
  \[
    v\tilde p^2z^2+2\tilde pz-v^2(v+2),
  \]
  which is positive because $z\ge1/2$, hence $2\tilde pz\ge1$, while $v^2(v+2)\le5/8$.
  Thus \eqref{eq:plt23-X-low-v-1} implies
  \begin{equation}\label{eq:plt23-X-low-v-2}
    X\ge\frac{(N+2)S_N}{\tilde p^2z}.
  \end{equation}
  Combining \eqref{eq:plt23-G-low-v}, \eqref{eq:plt23-X-low-v-2}, and \eqref{eq:plt23-SN-AMGM}, with $k=(N-1)/2\ge3$, gives
  \begin{equation}\label{eq:plt23-G-over-X}
    \frac GX
    \le
    \frac{\tilde p^3zP}{(N+2)S_N}
    \le
    \frac{2\tilde p^{3-k}zP}{(N+1)(N+2)}
    \le
    \frac{2zP}{(N+1)(N+2)}.
  \end{equation}
  Since $zP=z(N-cz)\le N^2/(4c)$,
  \begin{equation}\label{eq:plt23-G-over-X-c}
    \frac GX
    \le\frac{N^2}{2c(N+1)(N+2)}.
  \end{equation}
  Now $v\le1/2$ gives $c\ge1-2^{-N/2}$.
  For $N\ge7$, $2^{N/2}\ge N$, and
  \[
    \frac1N\le\frac{3N+2}{(N+1)(N+2)}.
  \]
  Hence
  \[
    c\ge
    1-\frac{3N+2}{(N+1)(N+2)}
    =\frac{N^2}{(N+1)(N+2)}.
  \]
  Substitution into \eqref{eq:plt23-G-over-X-c} gives $G/X\le1/2$, which is \eqref{eq:plt23-low-v-A}.

  We next retain both scales.
  From \eqref{eq:plt23-G-over-X} and $P\le N$,
  \begin{equation}\label{eq:plt23-G-over-X-z}
    \frac GX\le\frac{2z}{N+2}.
  \end{equation}
  For $Y$, \eqref{eq:plt23-AB-formulas}, \eqref{eq:plt23-KL-def}, and the first term of the sum in \eqref{eq:plt23-KL-sum} give
  \begin{align*}
    Y
    &=(1-\tilde M)\sqrt{2\Lambda}\,\widehat M_N\\
    &\ge
    (N+2)(1-\tilde M)(\tilde p-\tilde M)\tilde p^{N-1}
    \sqrt{\frac2{v+2}},
  \end{align*}
  where we used \eqref{eq:plt23-Lambda-lower}.
  Consequently,
  \begin{equation}\label{eq:plt23-2G-over-Y}
    \frac{2G}{Y}
    \le
    \frac{2P}{N+2}
    \frac{\sqrt{(v+2)/2}}
         {(1-\tilde M)(\tilde p-\tilde M)\tilde p^{N-2}}.
  \end{equation}
  Put $y=\sqrt v$, so $0<y\le1/\sqrt2$.
  Since $\tilde M\le y$ and $x\mapsto(1-x)(\tilde p-x)$ is decreasing on $[0,y]$, \Cref{lem:plt23-product} gives
  \begin{align}
   (1-\tilde M)(\tilde p-\tilde M)\tilde p^{N-2}
   &\ge
   (1-y)(1-y+y^2)(1+y^2)^{N-2}\notag\\
   &\ge
   \frac{N}{N+2}\sqrt{1+\frac{y^2}{2}}
   =\frac{N}{N+2}\sqrt{\frac{v+2}{2}}.
   \label{eq:plt23-product-applied}
  \end{align}
  Also
  \[
    c=1-v^{N/2}\ge\frac{N}{N+2}.
  \]
  Indeed, $v^{N/2}\le2^{-N/2}\le2/(N+2)$ for $N\ge7$.
  Therefore
  \begin{equation}\label{eq:plt23-P-bound}
    \frac PN=1-\frac{cz}{N}
    \le1-\frac{z}{N+2}
    =\frac{N+2-z}{N+2}.
  \end{equation}
  Substitution of \eqref{eq:plt23-product-applied} and \eqref{eq:plt23-P-bound} into \eqref{eq:plt23-2G-over-Y} gives
  \[
    \frac{2G}{Y}
    \le\frac{2(N+2-z)}{N+2}.
  \]
  Adding this to \eqref{eq:plt23-G-over-X-z} proves
  \[
    \frac GX+2\frac GY\le2,
  \]
  which is \eqref{eq:plt23-low-v-both} by \eqref{eq:plt23-GXY-ratios}.
\end{proof}

\subsection{The range \texorpdfstring{$1/2\le v<1$}{1/2<=v<1}}

The second half is even more direct.
Here the admissibility condition $q>1/3$ forces $\tilde M$ to be sufficiently close to $\sqrt v$; this is exactly what compensates for the apparently weak factors in $\mathcal A$ and $\mathcal B$.

Put
\[
  \tau:=\frac{\tilde M^2}{v},
  \qquad 0<\tau<1.
\]
Since $\tilde M^2=v\tau$, the inequality \eqref{eq:plt23-Delta-upper-three} gives
\[
  (1+v)^2+1+v\tau<3(1+v\tau),
\]
so
\[
  \tau>\frac{v^2+2v-1}{2v}\ge v^2.
\]
The last inequality follows from
\[
  v^2+2v-1-2v^3
  =(1-v)(1+v)(2v-1)\ge0.
\]
Hence
\begin{equation}\label{eq:plt23-vtau-square}
  (v+\tau)^2\le\tau(1+v)^2=\tau\tilde p^2,
\end{equation}
for
\[
  \tau\tilde p^2-(v+\tau)^2
  =(1-\tau)(\tau-v^2)\ge0.
\]

\begin{proposition}[The high-$v$ quadratic bound]\label{prop:plt23-high-v-A}
  Suppose $1/2\le v<1$.
  Then
  \begin{equation}\label{eq:plt23-high-v-A}
    F_N\le\mathcal A.
  \end{equation}
\end{proposition}

\begin{proof}
  Again the claim is immediate if $F_N\le0$, so assume \eqref{eq:plt23-F-positive}.
  From \eqref{eq:plt23-AB-formulas}, \eqref{eq:plt23-Delta-identities}, $u=v(1-\tau)/\tilde p$, and \eqref{eq:plt23-KA-def}--\eqref{eq:plt23-KA-sum},
  \[
    X
    =\frac{2\tilde p^2(1-\tau)}{\Delta(v+\tau)^2}
     \left(
       \frac{2\tilde M^N}{v}+(N+2)S_N
     \right).
  \]
  Using \eqref{eq:plt23-Delta-upper-sigma} and \eqref{eq:plt23-vtau-square}, and dropping the positive first term in parentheses,
  \begin{equation}\label{eq:plt23-X-high-v-lower}
    X>
    \frac{2(N+2)S_N(1-\tau)}{\tau\tilde p(v+2)}.
  \end{equation}

  The positivity of $F_N$ supplies a useful lower bound on $1-\tau$.
  From \eqref{eq:plt23-F-formula},
  \[
    (1+\tilde M^2)\tilde q^N<1+\tilde M^{N+2},
  \]
  so
  \[
    \tilde q^N
    <1-\frac{\tilde M^2}{1+\tilde M^2}(1-\tilde M^N).
  \]
  On the other hand, Bernoulli gives $\tilde q^N=(1-u)^N\ge1-Nu$.
  Comparing the two inequalities, and substituting $u=v(1-\tau)/\tilde p$ and $\tilde M^2=v\tau$, yields
  \[
    N(1+\tilde M^2)(1-\tau)
    >\tau\tilde p(1-\tilde M^N).
  \]
  Therefore \eqref{eq:plt23-X-high-v-lower} gives
  \begin{equation}\label{eq:plt23-X-high-v-final}
    X>
    \frac{2(N+2)S_N}{N(1+\tilde M^2)(v+2)}
    (1-\tilde M^N).
  \end{equation}

  By \eqref{eq:plt23-F-T} and \eqref{eq:plt23-T-bound-two},
  \begin{equation}\label{eq:plt23-G-high-v}
    G<T_N<N(1-\tilde M^{N+1}).
  \end{equation}
  For $0\le x<1$,
  \begin{equation}\label{eq:plt23-ratio-powers}
    \frac{1-x^N}{1-x^{N+1}}\ge\frac{N}{N+1}.
  \end{equation}
  Indeed, after clearing denominators this is $1-(N+1)x^N+Nx^{N+1}\ge0$, whose derivative is $-N(N+1)x^{N-1}(1-x)\le0$ and whose value at $x=1$ is zero.
  Combining \eqref{eq:plt23-X-high-v-final}, \eqref{eq:plt23-SN-AMGM}, and \eqref{eq:plt23-ratio-powers}, and using $1+\tilde M^2\le\tilde p$, it is enough for $X\ge G$ that
  \begin{equation}\label{eq:plt23-high-v-condition}
    \tilde p^{(N-3)/2}
    \ge\frac{N(v+2)}{N+2}.
  \end{equation}
  For $N\ge9$, the left side is at least $(3/2)^3=27/8>3$, whereas the right side is less than $3$.
  For $N=7$, \eqref{eq:plt23-high-v-condition} is
  \[
    (1+v)^2\ge\frac{7(v+2)}9.
  \]
  It holds at $v=1/2$, and the difference has derivative $2(1+v)-7/9>0$ thereafter.
  Thus $G\le X$, which is exactly \eqref{eq:plt23-high-v-A}.
\end{proof}

We next prove the bound for $\mathcal B$.
The only auxiliary fact is a simple comparison between a geometric sum and a power.

\begin{lemma}[A geometric-sum bound]\label{lem:plt23-geometric-sum}
  Let $N\ge7$ be odd and $1/\sqrt2\le x<1$.
  Then
  \begin{equation}\label{eq:plt23-geometric-sum}
    \frac{N+2}{2N}
    \sqrt{\frac2{2+x^2}}
    (1-x+x^2)(1+x^2)^{N-1}
    \ge
    \sum_{j=0}^N x^j.
  \end{equation}
\end{lemma}

\begin{proof}
  Write $N=2k+1$, so $k\ge3$, and put
  \[
    y:=x^2,
    \qquad
    C:=(1+y)^k.
  \]
  Then
  \begin{equation}\label{eq:plt23-geometric-sum-bound}
    \sum_{j=0}^N x^j
    =(1+x)\sum_{j=0}^k y^j
    \le (1+x)\bigl(C-(k-1)y\bigr).
  \end{equation}
  Indeed, $(1+y)^k-(k-1)y=1+y+\sum_{j=2}^k\binom{k}{j}y^j$, and every remaining binomial coefficient is at least one.

  Suppose first $k\ge4$, equivalently $N\ge9$.
  Since $x<1$,
  \[
    \sqrt{\frac2{2+x^2}}>\frac45,
    \qquad
    1-x+x^2\ge\frac34,
    \qquad
    \frac{N+2}{2N}>\frac12.
  \]
  Hence the left side of \eqref{eq:plt23-geometric-sum} is greater than
  \[
    \frac3{10}C^2.
  \]
  On the other hand, since $1+x\le2$ and $y\ge1/2$,
  \[
    (1+x)\bigl(C-(k-1)y\bigr)
    \le2C-(k-1).
  \]
  It remains to show
  \begin{equation}\label{eq:plt23-C-quadratic}
    \frac3{10}C^2\ge2C-(k-1).
  \end{equation}
  Here $C\ge(3/2)^k$.
  The left side minus the right side has derivative $3C/5-2$ in $C$, which is positive throughout this range because $C\ge(3/2)^4>10/3$; it is therefore enough to check \eqref{eq:plt23-C-quadratic} at $C=(3/2)^k$.
  At $k=4$ and $C=(3/2)^4$, the difference is
  \[
    \frac{1443}{2560}>0.
  \]
  Moreover, if $C=(3/2)^k$, the change in this difference under $k\mapsto k+1$ is
  \[
    \frac38C^2-C+1>0.
  \]
  Thus \eqref{eq:plt23-C-quadratic} holds for all $k\ge4$.

  It remains to take $k=3$, that is $N=7$.
  Define
  \[
    f(x):=\sqrt{\frac2{2+x^2}}(1-x+x^2).
  \]
  A direct differentiation gives
  \[
    f'(x)
    =\frac{\sqrt2\,(x^3+3x-2)}{(x^2+2)^{3/2}}>0
    \qquad\left(x\ge\frac1{\sqrt2}\right).
  \]
  At the left endpoint,
  \[
    f\left(\frac1{\sqrt2}\right)^2-\frac{49}{100}
    =\frac{171}{100}-\frac{6\sqrt2}{5}>0,
  \]
  because $171^2>2\cdot120^2$.
  Hence $f(x)>7/10$ on the whole interval.
  Therefore the left side of \eqref{eq:plt23-geometric-sum} is at least
  \[
    \frac9{14}\cdot\frac7{10}\,C^2
    =\frac9{20}C^2.
  \]
  By \eqref{eq:plt23-geometric-sum-bound}, the right side is at most $2(C-1)$.
  The difference $\frac9{20}C^2-2(C-1)$ has derivative $9C/10-2$, which is positive since $C\ge(3/2)^3=27/8$, and at $C=27/8$ the difference is
  \[
    \frac{481}{1280}>0.
  \]
  Hence the claim holds also for $N=7$.
\end{proof}

\begin{proposition}[The high-$v$ linear bound]\label{prop:plt23-high-v-B}
  Suppose $1/2\le v<1$.
  Then
  \begin{equation}\label{eq:plt23-high-v-B}
    2F_N\le\mathcal B.
  \end{equation}
\end{proposition}

\begin{proof}
  The claim is trivial if $F_N\le0$, so assume $F_N>0$.
  From \eqref{eq:plt23-AB-formulas}, \eqref{eq:plt23-KL-def}, \eqref{eq:plt23-KL-sum}, and \eqref{eq:plt23-Lambda-lower},
  \[
    Y\ge
    (N+2)(1-\tilde M)(\tilde p-\tilde M)\tilde p^{N-1}
    \sqrt{\frac2{v+2}}.
  \]
  Together with \eqref{eq:plt23-G-high-v}, it is enough to prove
  \begin{equation}\label{eq:plt23-high-v-B-sufficient}
    \frac{N+2}{2N}
    \sqrt{\frac2{v+2}}
    (\tilde p-\tilde M)\tilde p^{N-1}
    \ge
    \sum_{j=0}^N\tilde M^j,
  \end{equation}
  where we cancelled the positive factor $1-\tilde M$ using $1-\tilde M^{N+1}=(1-\tilde M)\sum_{j=0}^N\tilde M^j$.

  For fixed $v$, the left side of \eqref{eq:plt23-high-v-B-sufficient} decreases with $\tilde M$, while the right side increases.
  Since $\tilde M^2<v$, it is enough to put
  \[
    x:=\sqrt v,
    \qquad
    \tilde M=x,
    \qquad
    \tilde p=1+x^2.
  \]
  Then \eqref{eq:plt23-high-v-B-sufficient} becomes exactly \eqref{eq:plt23-geometric-sum}.
  This proves \eqref{eq:plt23-high-v-B}.
\end{proof}

\subsection{Completion of the scalar inequality}

We can now finish without any further subdivision of the admissible region.

\begin{proposition}[The scalar target]\label{prop:plt23-scalar}
  For every admissible pair $(q,\Lambda)$ in the intermediate-density region $1/3<q<1/2$, and every odd $m\ge9$, the scalar target \eqref{eq:plt23-scalar-target} holds.
\end{proposition}

\begin{proof}
  If $F_N=R_m/\Lambda^m\le0$, the claim is immediate.
  Assume $F_N>0$, and recall $\theta=\mathcal A/\mathcal B$.
  By \Cref{prop:plt23-reduction} it is enough to prove $F_N\le\mathcal B\,d(\theta)$.

  Suppose first $0<v\le1/2$.
  \Cref{prop:plt23-low-v} gives
  \[
    \frac{F_N}{\mathcal B}
    \le
    \min\left\{
      \frac\theta2,
      \frac{2\theta}{1+2\theta}
    \right\}.
  \]
  If $\theta\le1/2$, the first term on the right is at most $\theta=d(\theta)$.
  If $\theta\ge1/2$, the second term is at most $d(\theta)$, because
  \[
    1-\frac1{4\theta}-\frac{2\theta}{1+2\theta}
    =\frac{2\theta-1}{4\theta(1+2\theta)}\ge0.
  \]
  Thus $F_N\le\mathcal B\,d(\theta)$ in this range.

  Now suppose $1/2\le v<1$.
  \Cref{prop:plt23-high-v-A,prop:plt23-high-v-B} give
  \[
    \frac{F_N}{\mathcal B}
    \le\min\{\theta,1/2\}.
  \]
  If $\theta\le1/2$, the right side is at most $\theta=d(\theta)$.
  If $\theta\ge1/2$, it is at most $1/2\le d(\theta)$.
  Hence $F_N\le\mathcal B\,d(\theta)$ again, and \Cref{prop:plt23-reduction} proves \eqref{eq:plt23-scalar-target}.
\end{proof}

\subsection{Completion of the intermediate-density region}

We state explicitly how the preceding scalar argument reconnects to the graphon proof.
Assume
\[
  \frac12<p<\frac23,
  \qquad
  \frac13<q<\frac12,
\]
and let $m\ge9$ be odd.
If the compression $A$ has no eigenvalue above $q$, \Cref{lem:plt23-eigenvalues-above-q} already proves the desired odd-cycle inequality.
Otherwise there is a unique frontier eigenvalue $\Lambda$.
\Cref{prop:plt23-splitting-lower-bound} gives
\[
  t(C_m,W)-\bigl(p^m-pq^{m-1}\bigr)
  \ge -R_m+\Lambda_m\gamma^2+M_m\norm{g_s}^2,
\]
while \Cref{thm:plt23-huber-elim} gives
\[
  \Lambda_m\gamma^2+M_m\norm{g_s}^2
  \ge\Omega_m\psi(\xi,\rho).
\]
\Cref{prop:plt23-scalar} proves $R_m\le\Omega_m\psi(\xi,\rho)$.
Therefore
\[
  t(C_m,W)\ge p^m-pq^{m-1}.
\]
This completes the intermediate-density proof for every odd $m\ge9$.


\section{Assembly of the complete proof}\label{sec:assembly}

\begin{proof}[Proof of \Cref{thm:main}]
  Put $q=1-p$.

  If $p\le1/2$, the right side of \eqref{eq:main-goal} is nonpositive, so the claim follows from $t(C_m,W)\ge0$.
  Assume henceforth $p>1/2$.

  The cases $m=3,5,7$ follow at every density from \Cref{prop:C3,prop:C5,prop:C7}.
  Let $m\ge9$ be odd.

  If $p\ge2/3$, or equivalently $q\le1/3$, the result is \Cref{thm:pge23-main}.
  It remains to consider $1/2<p<2/3$, equivalently $1/3<q<1/2$.
  By \Cref{lem:step-reduction}, it is enough to treat a step graphon.
  If the compression $A$ has no eigenvalue above $q$, \Cref{lem:plt23-eigenvalues-above-q} proves the theorem.
  Otherwise there is a unique frontier eigenvalue $\Lambda$, and \Cref{prop:plt23-splitting-lower-bound,thm:plt23-huber-elim} reduce the problem to the scalar target \eqref{eq:plt23-scalar-target}, which is proved in \Cref{prop:plt23-scalar}.
  The step-graphon approximation returns the result to arbitrary graphons.
\end{proof}

\begin{remark}[What is finite and what is uniform]
  The dense-region argument is uniform in every odd $m\ge3$.
  The intermediate-region argument is uniform in every odd $m\ge9$, the only division of the parameter region being $v\le1/2$ against $v\ge1/2$.
  There is no length-by-length computation for $m=11,13$, and no adaptive interval subdivision anywhere in the proof.
\end{remark}

\begin{remark}[Equality and stability]
  The manuscript is organized to prove the inequality rather than classify all equality cases.
  A stability analysis would require tracking near-equality in the cycle-to-path deletion, the gamma moment inequality, the Hilbert--Schmidt bound, the variance lower bound, and the active Huber witness.
  The reductions here identify those loci but do not pursue the classification.
\end{remark}

\section{Transferable mechanisms and the orthogonal-polynomial viewpoint}
\label{sec:transferable}

This section is not needed for logical completeness.
Its purpose is to extract methods that may be useful in other spectral extremal problems.

\subsection{One-step stripping: the Schur complement as a Jacobi recursion}

Let $K$ be a finite-dimensional self-adjoint operator, let $e$ be a unit vector, and decompose
\[
  K=\begin{pmatrix}a&h^*\\h&B\end{pmatrix}
  \quad\text{on}\quad \R e\oplus e^\perp.
\]
The same block inversion used in \eqref{eq:pge23-path-gf} gives
\begin{equation}\label{eq:general-stripping}
  M_{K,e}(z):=\ip e{(I-zK)^{-1}e}
  =\frac{1}{1-az-z^2\ip h{(I-zB)^{-1}h}}.
\end{equation}
If $h\ne0$, let $\nu$ be the probability spectral measure of $B$ at $h/\norm h$.
Then
\begin{equation}\label{eq:general-stripping-normalized}
  M_{K,e}(z)^{-1}
  =1-az-\norm h^2z^2M_\nu(z).
\end{equation}
This is exactly the first step of the Jacobi continued fraction.
Iterating on the cyclic Krylov space generated by $e$ produces a tridiagonal Jacobi matrix and its full J-fraction.
In the centered, variance-one normalization, \eqref{eq:general-stripping-normalized} becomes
\[
  1-M_\mu(z)^{-1}=z^2M_\nu(z),
\]
the classical once-stripped moment relation emphasized by Anshelevich~\cite{Anshelevich2009}; coefficient stripping for scalar and matrix Jacobi operators is surveyed by Damanik--Pushnitski--Simon \cite{DamanikPushnitskiSimon2008}.

In our graphon proof, $e=\one$, $a=q$, $h=g$, and $B=A$.
Thus the complement degree fluctuation $g$ is the first off-diagonal Jacobi coefficient, while $A$ is the once-stripped tail.
The proof does not need to tridiagonalize the tail; it retains its spectral measure $\mu$ directly.
This is useful in infinite-dimensional problems where the first Jacobi step is rigid but the tail is not.

\subsection{First returns, irreducible excursions, and logarithmic cycles}

Equation \eqref{eq:pge23-path-gf} has the first-return form
\[
  M(z)=\frac{1}{1-H(z)},
  \qquad H(z)=qz+z^2R_+(z).
\]
Combinatorially, a closed walk from the distinguished state is a sequence of irreducible returns: either a level step of weight $q$, or a departure into the orthogonal tail followed by a return, encoded by $z^2R_+$.
This is the same structural decomposition that underlies Flajolet's continued fraction for weighted Motzkin paths~\cite{Flajolet1980} and Lehner's use of Jacobi matrices, lattice paths, and irreducible paths in cumulant expansions~\cite{Lehner2003}.

The logarithms $-\log(1-Y_\pm)=\sum_{r\ge1}Y_\pm^r/r$ perform a related but cyclic assembly: $r$ excursions are concatenated with the cyclic weight $1/r$.
The two-sided identity then applies a parity projection: for odd coefficients, the stripped-tail terms from $A$ and $-A$ cancel.
A general design principle suggested by this proof is:
\begin{quote}
  Separate a distinguished state, encode returns by a Schur complement, and combine a model with a signed or complemented companion so that the uncontrolled tail moments cancel in the desired parity class.
\end{quote}
This ``strip, assemble, cancel'' pattern is broader than graphons.

\subsection{Dirichlet diagonalization as probabilistic polarization}

The calculation in \Cref{lem:pge23-dirichlet} is an instance of the following general lemma.

\begin{proposition}[Dirichlet polarization]\label{prop:general-dirichlet}
  Fix $r\ge1$.
  Suppose
  \[
    F(\lambda_1,\ldots,\lambda_r)
    =\sum_{j=0}^d c_jh_j(\lambda_1,\ldots,\lambda_r),
  \]
  and define
  \[
    \widetilde F(\ell)
    =\sum_{j=0}^d c_j\binom{j+r-1}{r-1}\ell^j.
  \]
  If $\Theta\sim\operatorname{Dirichlet}(1,\ldots,1)$, then
  \[
    F(\lambda_1,\ldots,\lambda_r)
    =\E\widetilde F(\Theta\cdot\lambda).
  \]
  Consequently, for every interval $I\subset\R$,
  \[
    \min_{\lambda\in I^r}F(\lambda)
    =\min_{\ell\in I}\widetilde F(\ell).
  \]
\end{proposition}

\begin{proof}
  This is exactly \eqref{eq:pge23-dirichlet-moment}, followed by linearity.
  The minimum statement follows because $\Theta\cdot\lambda\in I$, and the reverse inequality is obtained on the diagonal.
\end{proof}

This lemma may be viewed as a positive, probabilistic version of polarization: the complete homogeneous basis is precisely the basis for which restriction to the diagonal can be inverted by averaging over a simplex.
The law of $\Theta\cdot\lambda$ is a Dirichlet average and, in one dimension, is closely related to B-splines and the Hermite--Genocchi formula for divided differences; see de~Boor~\cite{deBoor2005}.
The key extremal consequence is unusually strong: positivity on a whole spectral cube is equivalent to positivity on the diagonal whenever the polynomial has this basis structure.

\subsection{Beta--gamma algebra and Laplace positivity transfer}

The passage from \eqref{eq:pge23-beta-integral} to \eqref{eq:pge23-laplace-gamma} is an instance of a general positive transform.
For suitable $f$, integers $m>r\ge1$, and $\ell>0$,
\begin{equation}\label{eq:general-laplace-transfer}
  \int_0^\infty\frac{s^{r-1}}{(\ell+s)^m}f(s)\dd s
  =\frac{\Gamma(r)}{\Gamma(m)}
  \int_0^\infty t^{m-r-1}e^{-\ell t}
  \E f(Y/t)\dd t,
  \qquad Y\sim\GammaLaw(r,1).
\end{equation}
The proof is just the Laplace representation of $(\ell+s)^{-m}$, followed by $y=ts$.
Its value is conceptual: a parameter-dependent rational kernel becomes a positive mixture of expectations under a fixed probability family.
Therefore any pointwise-in-$t$ expectation inequality for $f(Y/t)$ transfers automatically to every $\ell$.

In the present proof, the repeated variables in the complete homogeneous polynomials first produce a beta distribution; the projective coordinate $s=\ell(1-x)/x$ turns the beta density into a beta-prime kernel; and the Laplace transform reveals a gamma variable.
This beta--gamma route is a useful search heuristic whenever repeated spectral parameters and rational powers occur together.

\subsection{The gamma Stein identity and a hidden Laguerre equation}

Equation \eqref{eq:pge23-gamma-Stein} is the standard gamma Stein identity; the operator
\[
  \mathcal A_rf(y)=yf'(y)+(r-y)f(y)
\]
characterizes the $\GammaLaw(r,1)$ law and is widely used in Stein's method; see, for example, Gaunt~\cite{Gaunt2016}.
Here it is used not for distributional approximation but as a ladder relation for shifted moments.

There is an exact orthogonal-polynomial shadow.
For an integer $k\ge0$, put
\[
  F_k(b)=\E(Y-b)^k.
\]
Expansion of the moment, or the Rodrigues formula, gives
\begin{equation}\label{eq:Laguerre-identity}
  F_k(b)=(-1)^k k!\,L_k^{(-r-k)}(-b),
\end{equation}
where $L_k^{(\alpha)}$ is the generalized Laguerre polynomial.
For $k=2j$, the Laguerre differential equation is exactly \eqref{eq:pge23-gamma-ODE}.
The parameter $-r-2j$ lies outside the classical orthogonality range $\alpha>-1$, so we are using the algebraic differential and ladder structure of Laguerre polynomials, not a positive Laguerre orthogonality measure.
Standard background is in Szeg\H{o} \cite{Szego1975}.

The function $G$ in \eqref{eq:pge23-gamma-G} is a weighted amplitude whose critical points factor explicitly.
Proving that its left maximum dominates its right maximum and then using \eqref{eq:pge23-gamma-ibp-max} is analogous in spirit to Sonin--Sturm arguments for comparing successive extrema of solutions to second-order equations.
The important transferable move is: construct a first-order multiplier so that the desired linear combination of a polynomial solution and its derivative becomes an integral of a factored derivative; positivity is then reduced to a global maximum comparison.

This explains the connection with the orthogonal-polynomial techniques suggested by the two cited papers.
Anshelevich's once-stripped moment identity is literally our first Schur step, while Lehner's Jacobi-matrix and Motzkin-path picture explains the return decomposition of the moment series.
The later gamma lemma lands in a different classical family---Laguerre---but uses the same general philosophy: a recurrence or differential equation is more effective than coefficientwise estimation of high moments.

\subsection{The one-dangerous-mode principle and Huber duality}

The intermediate region exhibits a complementary mechanism.
The square bound says that two eigenvalues above $q$ would cost more than $pq$ when $q>1/3$, so there is at most one dangerous mode.
Entrywise constraints on the corresponding eigenfunction force a choice: either $g$ couples substantially to that mode, or a quantitatively controlled part of $g$ must lie in the safe subspace.
Retaining both alternatives leads to
\[
  \rho v^2+(\xi-v+v^2)_+,
\]
and minimization over the shape variable $v$ produces the Huber envelope $\psi(\xi,\rho)$.
Its dual variable $\lambda\in[0,1]$ is an active-set multiplier.
The interior optimum $\lambda=2\rho\xi$ and the saturated value $\lambda=1$ are the quadratic and linear regimes, and keeping both at once removes the need to split the parameter region according to which of them is active.

This suggests a general recipe for constrained spectral inequalities:
\begin{enumerate}[leftmargin=2.2em]
  \item use a trace or energy bound to isolate a finite set of dangerous modes;
  \item derive geometric constraints that force compensation either in direct coupling or in the safe complement;
  \item keep the channels simultaneously rather than estimating them separately;
  \item eliminate the shape by convex duality, and choose dual witnesses adapted to small and large activation.
\end{enumerate}
The method should extend to several dangerous modes, where the scalar Huber envelope is replaced by a low-dimensional convex program.

\subsection{Why \texorpdfstring{$q=1/3$}{q=1/3} is the natural interface}

The same threshold appears for independent structural reasons.

On the dense side, centering at $1/2$ gives $\delta=1/2-q$.
The gamma moment lemma pays for the odd term in \eqref{eq:pge23-rho-centered} precisely when
\[
  2\delta\ge\frac{1}{3},
  \qquad\text{i.e.}\qquad q\le\frac{1}{3}.
\]
On the intermediate side, two eigenvalues larger than $q$ would consume more than $2q^2$ of the bound \eqref{eq:plt23-HS-bound}, and uniqueness is guaranteed precisely when
\[
  2q^2>pq=q(1-q),
  \qquad\text{i.e.}\qquad q>\frac{1}{3}.
\]
Thus the two mechanisms meet without overlap or gap.
This matching strongly suggests that $q=1/3$ is an intrinsic phase boundary of the proof problem, not an artifact of constants chosen later in the argument.

\begin{thebibliography}{99}

  \bibitem{Anshelevich2009} M.~Anshelevich, \newblock Appell polynomials and their relatives III: conditionally free theory, \newblock \emph{Illinois Journal of Mathematics} \textbf{53} (2009), 39--66; \newblock arXiv:0803.4279.

  \bibitem{deBoor2005} C.~de~Boor, \newblock Divided differences, \newblock \emph{Surveys in Approximation Theory} \textbf{1} (2005), 46--69.

  \bibitem{DamanikPushnitskiSimon2008} D.~Damanik, A.~Pushnitski, and B.~Simon, \newblock The analytic theory of matrix orthogonal polynomials, \newblock \emph{Surveys in Approximation Theory} \textbf{4} (2008), 1--85.

  \bibitem{Flajolet1980} P.~Flajolet, \newblock Combinatorial aspects of continued fractions, \newblock \emph{Discrete Mathematics} \textbf{32} (1980), 125--161.

  \bibitem{Gaunt2016} R.~E. Gaunt, \newblock Products of normal, beta and gamma random variables: Stein operators and distributional theory, \newblock \emph{Brazilian Journal of Probability and Statistics} \textbf{32} (2018), 437--466; \newblock arXiv:1507.07696.

  \bibitem{Goodman1959} A.~W. Goodman, \newblock On sets of acquaintances and strangers at any party, \newblock \emph{American Mathematical Monthly} \textbf{66} (1959), 778--783.

  \bibitem{Huber1964} P.~J. Huber, \newblock Robust estimation of a location parameter, \newblock \emph{Annals of Mathematical Statistics} \textbf{35} (1964), 73--101.

  \bibitem{Lehner2003} F.~Lehner, \newblock Cumulants, lattice paths, and orthogonal polynomials, \newblock \emph{Discrete Mathematics} \textbf{270} (2003), 177--191; \newblock arXiv:math/0110031.

  \bibitem{Lovasz2012} L.~Lov\'asz, \newblock \emph{Large Networks and Graph Limits}, \newblock American Mathematical Society Colloquium Publications, vol.~60, American Mathematical Society, Providence, RI, 2012.

  \bibitem{PowersReznick2000} V.~Powers and B.~Reznick, \newblock Polynomials that are positive on an interval, \newblock \emph{Transactions of the American Mathematical Society} \textbf{352} (2000), 4677--4692.

  \bibitem{Sion1958} M.~Sion, \newblock On general minimax theorems, \newblock \emph{Pacific Journal of Mathematics} \textbf{8} (1958), 171--176.

  \bibitem{Szego1975} G.~Szeg\H{o}, \newblock \emph{Orthogonal Polynomials}, fourth edition, \newblock American Mathematical Society Colloquium Publications, vol.~23, American Mathematical Society, Providence, RI, 1975.

\end{thebibliography}

\end{document}
